\documentclass{article} % For LaTeX2e
\usepackage{iclr2025_conference,times}

% Optional math commands from https://github.com/goodfeli/dlbook_notation.
%%%%% NEW MATH DEFINITIONS %%%%%

\usepackage{amsmath,amsfonts,bm}

% Mark sections of captions for referring to divisions of figures
\newcommand{\figleft}{{\em (Left)}}
\newcommand{\figcenter}{{\em (Center)}}
\newcommand{\figright}{{\em (Right)}}
\newcommand{\figtop}{{\em (Top)}}
\newcommand{\figbottom}{{\em (Bottom)}}
\newcommand{\captiona}{{\em (a)}}
\newcommand{\captionb}{{\em (b)}}
\newcommand{\captionc}{{\em (c)}}
\newcommand{\captiond}{{\em (d)}}

% Highlight a newly defined term
\newcommand{\newterm}[1]{{\bf #1}}


% Figure reference, lower-case.
\def\figref#1{figure~\ref{#1}}
% Figure reference, capital. For start of sentence
\def\Figref#1{Figure~\ref{#1}}
\def\twofigref#1#2{figures \ref{#1} and \ref{#2}}
\def\quadfigref#1#2#3#4{figures \ref{#1}, \ref{#2}, \ref{#3} and \ref{#4}}
% Section reference, lower-case.
\def\secref#1{section~\ref{#1}}
% Section reference, capital.
\def\Secref#1{Section~\ref{#1}}
% Reference to two sections.
\def\twosecrefs#1#2{sections \ref{#1} and \ref{#2}}
% Reference to three sections.
\def\secrefs#1#2#3{sections \ref{#1}, \ref{#2} and \ref{#3}}
% Reference to an equation, lower-case.
\def\eqref#1{equation~\ref{#1}}
% Reference to an equation, upper case
\def\Eqref#1{Equation~\ref{#1}}
% A raw reference to an equation---avoid using if possible
\def\plaineqref#1{\ref{#1}}
% Reference to a chapter, lower-case.
\def\chapref#1{chapter~\ref{#1}}
% Reference to an equation, upper case.
\def\Chapref#1{Chapter~\ref{#1}}
% Reference to a range of chapters
\def\rangechapref#1#2{chapters\ref{#1}--\ref{#2}}
% Reference to an algorithm, lower-case.
\def\algref#1{algorithm~\ref{#1}}
% Reference to an algorithm, upper case.
\def\Algref#1{Algorithm~\ref{#1}}
\def\twoalgref#1#2{algorithms \ref{#1} and \ref{#2}}
\def\Twoalgref#1#2{Algorithms \ref{#1} and \ref{#2}}
% Reference to a part, lower case
\def\partref#1{part~\ref{#1}}
% Reference to a part, upper case
\def\Partref#1{Part~\ref{#1}}
\def\twopartref#1#2{parts \ref{#1} and \ref{#2}}

\def\ceil#1{\lceil #1 \rceil}
\def\floor#1{\lfloor #1 \rfloor}
\def\1{\bm{1}}
\newcommand{\train}{\mathcal{D}}
\newcommand{\valid}{\mathcal{D_{\mathrm{valid}}}}
\newcommand{\test}{\mathcal{D_{\mathrm{test}}}}

\def\eps{{\epsilon}}


% Random variables
\def\reta{{\textnormal{$\eta$}}}
\def\ra{{\textnormal{a}}}
\def\rb{{\textnormal{b}}}
\def\rc{{\textnormal{c}}}
\def\rd{{\textnormal{d}}}
\def\re{{\textnormal{e}}}
\def\rf{{\textnormal{f}}}
\def\rg{{\textnormal{g}}}
\def\rh{{\textnormal{h}}}
\def\ri{{\textnormal{i}}}
\def\rj{{\textnormal{j}}}
\def\rk{{\textnormal{k}}}
\def\rl{{\textnormal{l}}}
% rm is already a command, just don't name any random variables m
\def\rn{{\textnormal{n}}}
\def\ro{{\textnormal{o}}}
\def\rp{{\textnormal{p}}}
\def\rq{{\textnormal{q}}}
\def\rr{{\textnormal{r}}}
\def\rs{{\textnormal{s}}}
\def\rt{{\textnormal{t}}}
\def\ru{{\textnormal{u}}}
\def\rv{{\textnormal{v}}}
\def\rw{{\textnormal{w}}}
\def\rx{{\textnormal{x}}}
\def\ry{{\textnormal{y}}}
\def\rz{{\textnormal{z}}}

% Random vectors
\def\rvepsilon{{\mathbf{\epsilon}}}
\def\rvtheta{{\mathbf{\theta}}}
\def\rva{{\mathbf{a}}}
\def\rvb{{\mathbf{b}}}
\def\rvc{{\mathbf{c}}}
\def\rvd{{\mathbf{d}}}
\def\rve{{\mathbf{e}}}
\def\rvf{{\mathbf{f}}}
\def\rvg{{\mathbf{g}}}
\def\rvh{{\mathbf{h}}}
\def\rvu{{\mathbf{i}}}
\def\rvj{{\mathbf{j}}}
\def\rvk{{\mathbf{k}}}
\def\rvl{{\mathbf{l}}}
\def\rvm{{\mathbf{m}}}
\def\rvn{{\mathbf{n}}}
\def\rvo{{\mathbf{o}}}
\def\rvp{{\mathbf{p}}}
\def\rvq{{\mathbf{q}}}
\def\rvr{{\mathbf{r}}}
\def\rvs{{\mathbf{s}}}
\def\rvt{{\mathbf{t}}}
\def\rvu{{\mathbf{u}}}
\def\rvv{{\mathbf{v}}}
\def\rvw{{\mathbf{w}}}
\def\rvx{{\mathbf{x}}}
\def\rvy{{\mathbf{y}}}
\def\rvz{{\mathbf{z}}}

% Elements of random vectors
\def\erva{{\textnormal{a}}}
\def\ervb{{\textnormal{b}}}
\def\ervc{{\textnormal{c}}}
\def\ervd{{\textnormal{d}}}
\def\erve{{\textnormal{e}}}
\def\ervf{{\textnormal{f}}}
\def\ervg{{\textnormal{g}}}
\def\ervh{{\textnormal{h}}}
\def\ervi{{\textnormal{i}}}
\def\ervj{{\textnormal{j}}}
\def\ervk{{\textnormal{k}}}
\def\ervl{{\textnormal{l}}}
\def\ervm{{\textnormal{m}}}
\def\ervn{{\textnormal{n}}}
\def\ervo{{\textnormal{o}}}
\def\ervp{{\textnormal{p}}}
\def\ervq{{\textnormal{q}}}
\def\ervr{{\textnormal{r}}}
\def\ervs{{\textnormal{s}}}
\def\ervt{{\textnormal{t}}}
\def\ervu{{\textnormal{u}}}
\def\ervv{{\textnormal{v}}}
\def\ervw{{\textnormal{w}}}
\def\ervx{{\textnormal{x}}}
\def\ervy{{\textnormal{y}}}
\def\ervz{{\textnormal{z}}}

% Random matrices
\def\rmA{{\mathbf{A}}}
\def\rmB{{\mathbf{B}}}
\def\rmC{{\mathbf{C}}}
\def\rmD{{\mathbf{D}}}
\def\rmE{{\mathbf{E}}}
\def\rmF{{\mathbf{F}}}
\def\rmG{{\mathbf{G}}}
\def\rmH{{\mathbf{H}}}
\def\rmI{{\mathbf{I}}}
\def\rmJ{{\mathbf{J}}}
\def\rmK{{\mathbf{K}}}
\def\rmL{{\mathbf{L}}}
\def\rmM{{\mathbf{M}}}
\def\rmN{{\mathbf{N}}}
\def\rmO{{\mathbf{O}}}
\def\rmP{{\mathbf{P}}}
\def\rmQ{{\mathbf{Q}}}
\def\rmR{{\mathbf{R}}}
\def\rmS{{\mathbf{S}}}
\def\rmT{{\mathbf{T}}}
\def\rmU{{\mathbf{U}}}
\def\rmV{{\mathbf{V}}}
\def\rmW{{\mathbf{W}}}
\def\rmX{{\mathbf{X}}}
\def\rmY{{\mathbf{Y}}}
\def\rmZ{{\mathbf{Z}}}

% Elements of random matrices
\def\ermA{{\textnormal{A}}}
\def\ermB{{\textnormal{B}}}
\def\ermC{{\textnormal{C}}}
\def\ermD{{\textnormal{D}}}
\def\ermE{{\textnormal{E}}}
\def\ermF{{\textnormal{F}}}
\def\ermG{{\textnormal{G}}}
\def\ermH{{\textnormal{H}}}
\def\ermI{{\textnormal{I}}}
\def\ermJ{{\textnormal{J}}}
\def\ermK{{\textnormal{K}}}
\def\ermL{{\textnormal{L}}}
\def\ermM{{\textnormal{M}}}
\def\ermN{{\textnormal{N}}}
\def\ermO{{\textnormal{O}}}
\def\ermP{{\textnormal{P}}}
\def\ermQ{{\textnormal{Q}}}
\def\ermR{{\textnormal{R}}}
\def\ermS{{\textnormal{S}}}
\def\ermT{{\textnormal{T}}}
\def\ermU{{\textnormal{U}}}
\def\ermV{{\textnormal{V}}}
\def\ermW{{\textnormal{W}}}
\def\ermX{{\textnormal{X}}}
\def\ermY{{\textnormal{Y}}}
\def\ermZ{{\textnormal{Z}}}

% Vectors
\def\vzero{{\bm{0}}}
\def\vone{{\bm{1}}}
\def\vmu{{\bm{\mu}}}
\def\vtheta{{\bm{\theta}}}
\def\va{{\bm{a}}}
\def\vb{{\bm{b}}}
\def\vc{{\bm{c}}}
\def\vd{{\bm{d}}}
\def\ve{{\bm{e}}}
\def\vf{{\bm{f}}}
\def\vg{{\bm{g}}}
\def\vh{{\bm{h}}}
\def\vi{{\bm{i}}}
\def\vj{{\bm{j}}}
\def\vk{{\bm{k}}}
\def\vl{{\bm{l}}}
\def\vm{{\bm{m}}}
\def\vn{{\bm{n}}}
\def\vo{{\bm{o}}}
\def\vp{{\bm{p}}}
\def\vq{{\bm{q}}}
\def\vr{{\bm{r}}}
\def\vs{{\bm{s}}}
\def\vt{{\bm{t}}}
\def\vu{{\bm{u}}}
\def\vv{{\bm{v}}}
\def\vw{{\bm{w}}}
\def\vx{{\bm{x}}}
\def\vy{{\bm{y}}}
\def\vz{{\bm{z}}}

% Elements of vectors
\def\evalpha{{\alpha}}
\def\evbeta{{\beta}}
\def\evepsilon{{\epsilon}}
\def\evlambda{{\lambda}}
\def\evomega{{\omega}}
\def\evmu{{\mu}}
\def\evpsi{{\psi}}
\def\evsigma{{\sigma}}
\def\evtheta{{\theta}}
\def\eva{{a}}
\def\evb{{b}}
\def\evc{{c}}
\def\evd{{d}}
\def\eve{{e}}
\def\evf{{f}}
\def\evg{{g}}
\def\evh{{h}}
\def\evi{{i}}
\def\evj{{j}}
\def\evk{{k}}
\def\evl{{l}}
\def\evm{{m}}
\def\evn{{n}}
\def\evo{{o}}
\def\evp{{p}}
\def\evq{{q}}
\def\evr{{r}}
\def\evs{{s}}
\def\evt{{t}}
\def\evu{{u}}
\def\evv{{v}}
\def\evw{{w}}
\def\evx{{x}}
\def\evy{{y}}
\def\evz{{z}}

% Matrix
\def\mA{{\bm{A}}}
\def\mB{{\bm{B}}}
\def\mC{{\bm{C}}}
\def\mD{{\bm{D}}}
\def\mE{{\bm{E}}}
\def\mF{{\bm{F}}}
\def\mG{{\bm{G}}}
\def\mH{{\bm{H}}}
\def\mI{{\bm{I}}}
\def\mJ{{\bm{J}}}
\def\mK{{\bm{K}}}
\def\mL{{\bm{L}}}
\def\mM{{\bm{M}}}
\def\mN{{\bm{N}}}
\def\mO{{\bm{O}}}
\def\mP{{\bm{P}}}
\def\mQ{{\bm{Q}}}
\def\mR{{\bm{R}}}
\def\mS{{\bm{S}}}
\def\mT{{\bm{T}}}
\def\mU{{\bm{U}}}
\def\mV{{\bm{V}}}
\def\mW{{\bm{W}}}
\def\mX{{\bm{X}}}
\def\mY{{\bm{Y}}}
\def\mZ{{\bm{Z}}}
\def\mBeta{{\bm{\beta}}}
\def\mPhi{{\bm{\Phi}}}
\def\mLambda{{\bm{\Lambda}}}
\def\mSigma{{\bm{\Sigma}}}

% Tensor
\DeclareMathAlphabet{\mathsfit}{\encodingdefault}{\sfdefault}{m}{sl}
\SetMathAlphabet{\mathsfit}{bold}{\encodingdefault}{\sfdefault}{bx}{n}
\newcommand{\tens}[1]{\bm{\mathsfit{#1}}}
\def\tA{{\tens{A}}}
\def\tB{{\tens{B}}}
\def\tC{{\tens{C}}}
\def\tD{{\tens{D}}}
\def\tE{{\tens{E}}}
\def\tF{{\tens{F}}}
\def\tG{{\tens{G}}}
\def\tH{{\tens{H}}}
\def\tI{{\tens{I}}}
\def\tJ{{\tens{J}}}
\def\tK{{\tens{K}}}
\def\tL{{\tens{L}}}
\def\tM{{\tens{M}}}
\def\tN{{\tens{N}}}
\def\tO{{\tens{O}}}
\def\tP{{\tens{P}}}
\def\tQ{{\tens{Q}}}
\def\tR{{\tens{R}}}
\def\tS{{\tens{S}}}
\def\tT{{\tens{T}}}
\def\tU{{\tens{U}}}
\def\tV{{\tens{V}}}
\def\tW{{\tens{W}}}
\def\tX{{\tens{X}}}
\def\tY{{\tens{Y}}}
\def\tZ{{\tens{Z}}}


% Graph
\def\gA{{\mathcal{A}}}
\def\gB{{\mathcal{B}}}
\def\gC{{\mathcal{C}}}
\def\gD{{\mathcal{D}}}
\def\gE{{\mathcal{E}}}
\def\gF{{\mathcal{F}}}
\def\gG{{\mathcal{G}}}
\def\gH{{\mathcal{H}}}
\def\gI{{\mathcal{I}}}
\def\gJ{{\mathcal{J}}}
\def\gK{{\mathcal{K}}}
\def\gL{{\mathcal{L}}}
\def\gM{{\mathcal{M}}}
\def\gN{{\mathcal{N}}}
\def\gO{{\mathcal{O}}}
\def\gP{{\mathcal{P}}}
\def\gQ{{\mathcal{Q}}}
\def\gR{{\mathcal{R}}}
\def\gS{{\mathcal{S}}}
\def\gT{{\mathcal{T}}}
\def\gU{{\mathcal{U}}}
\def\gV{{\mathcal{V}}}
\def\gW{{\mathcal{W}}}
\def\gX{{\mathcal{X}}}
\def\gY{{\mathcal{Y}}}
\def\gZ{{\mathcal{Z}}}

% Sets
\def\sA{{\mathbb{A}}}
\def\sB{{\mathbb{B}}}
\def\sC{{\mathbb{C}}}
\def\sD{{\mathbb{D}}}
% Don't use a set called E, because this would be the same as our symbol
% for expectation.
\def\sF{{\mathbb{F}}}
\def\sG{{\mathbb{G}}}
\def\sH{{\mathbb{H}}}
\def\sI{{\mathbb{I}}}
\def\sJ{{\mathbb{J}}}
\def\sK{{\mathbb{K}}}
\def\sL{{\mathbb{L}}}
\def\sM{{\mathbb{M}}}
\def\sN{{\mathbb{N}}}
\def\sO{{\mathbb{O}}}
\def\sP{{\mathbb{P}}}
\def\sQ{{\mathbb{Q}}}
\def\sR{{\mathbb{R}}}
\def\sS{{\mathbb{S}}}
\def\sT{{\mathbb{T}}}
\def\sU{{\mathbb{U}}}
\def\sV{{\mathbb{V}}}
\def\sW{{\mathbb{W}}}
\def\sX{{\mathbb{X}}}
\def\sY{{\mathbb{Y}}}
\def\sZ{{\mathbb{Z}}}

% Entries of a matrix
\def\emLambda{{\Lambda}}
\def\emA{{A}}
\def\emB{{B}}
\def\emC{{C}}
\def\emD{{D}}
\def\emE{{E}}
\def\emF{{F}}
\def\emG{{G}}
\def\emH{{H}}
\def\emI{{I}}
\def\emJ{{J}}
\def\emK{{K}}
\def\emL{{L}}
\def\emM{{M}}
\def\emN{{N}}
\def\emO{{O}}
\def\emP{{P}}
\def\emQ{{Q}}
\def\emR{{R}}
\def\emS{{S}}
\def\emT{{T}}
\def\emU{{U}}
\def\emV{{V}}
\def\emW{{W}}
\def\emX{{X}}
\def\emY{{Y}}
\def\emZ{{Z}}
\def\emSigma{{\Sigma}}

% entries of a tensor
% Same font as tensor, without \bm wrapper
\newcommand{\etens}[1]{\mathsfit{#1}}
\def\etLambda{{\etens{\Lambda}}}
\def\etA{{\etens{A}}}
\def\etB{{\etens{B}}}
\def\etC{{\etens{C}}}
\def\etD{{\etens{D}}}
\def\etE{{\etens{E}}}
\def\etF{{\etens{F}}}
\def\etG{{\etens{G}}}
\def\etH{{\etens{H}}}
\def\etI{{\etens{I}}}
\def\etJ{{\etens{J}}}
\def\etK{{\etens{K}}}
\def\etL{{\etens{L}}}
\def\etM{{\etens{M}}}
\def\etN{{\etens{N}}}
\def\etO{{\etens{O}}}
\def\etP{{\etens{P}}}
\def\etQ{{\etens{Q}}}
\def\etR{{\etens{R}}}
\def\etS{{\etens{S}}}
\def\etT{{\etens{T}}}
\def\etU{{\etens{U}}}
\def\etV{{\etens{V}}}
\def\etW{{\etens{W}}}
\def\etX{{\etens{X}}}
\def\etY{{\etens{Y}}}
\def\etZ{{\etens{Z}}}

% The true underlying data generating distribution
\newcommand{\pdata}{p_{\rm{data}}}
% The empirical distribution defined by the training set
\newcommand{\ptrain}{\hat{p}_{\rm{data}}}
\newcommand{\Ptrain}{\hat{P}_{\rm{data}}}
% The model distribution
\newcommand{\pmodel}{p_{\rm{model}}}
\newcommand{\Pmodel}{P_{\rm{model}}}
\newcommand{\ptildemodel}{\tilde{p}_{\rm{model}}}
% Stochastic autoencoder distributions
\newcommand{\pencode}{p_{\rm{encoder}}}
\newcommand{\pdecode}{p_{\rm{decoder}}}
\newcommand{\precons}{p_{\rm{reconstruct}}}

\newcommand{\laplace}{\mathrm{Laplace}} % Laplace distribution

\newcommand{\E}{\mathbb{E}}
\newcommand{\Ls}{\mathcal{L}}
\newcommand{\R}{\mathbb{R}}
\newcommand{\emp}{\tilde{p}}
\newcommand{\lr}{\alpha}
\newcommand{\reg}{\lambda}
\newcommand{\rect}{\mathrm{rectifier}}
\newcommand{\softmax}{\mathrm{softmax}}
\newcommand{\sigmoid}{\sigma}
\newcommand{\softplus}{\zeta}
\newcommand{\KL}{D_{\mathrm{KL}}}
\newcommand{\Var}{\mathrm{Var}}
\newcommand{\standarderror}{\mathrm{SE}}
\newcommand{\Cov}{\mathrm{Cov}}
% Wolfram Mathworld says $L^2$ is for function spaces and $\ell^2$ is for vectors
% But then they seem to use $L^2$ for vectors throughout the site, and so does
% wikipedia.
\newcommand{\normlzero}{L^0}
\newcommand{\normlone}{L^1}
\newcommand{\normltwo}{L^2}
\newcommand{\normlp}{L^p}
\newcommand{\normmax}{L^\infty}

\newcommand{\parents}{Pa} % See usage in notation.tex. Chosen to match Daphne's book.

\DeclareMathOperator*{\argmax}{arg\,max}
\DeclareMathOperator*{\argmin}{arg\,min}

\DeclareMathOperator{\sign}{sign}
\DeclareMathOperator{\Tr}{Tr}
\let\ab\allowbreak


\usepackage{booktabs}
\usepackage{multirow}
\usepackage[dvipsnames]{xcolor}
\usepackage{url}
\usepackage{graphicx}
\usepackage{subcaption} % for subfigure environment
\usepackage{hyperref}
\usepackage{float}

\usepackage{makecell}
\usepackage{amssymb}
\usepackage{pifont}
\usepackage{tabularx}
\usepackage{enumitem}
\usepackage{wrapfig}
\usepackage{tikz}

\usepackage[compact]{titlesec}
\titlespacing{\section}{0pt}{*1}{*0.5}
\titlespacing{\subsection}{0pt}{*0.5}{*0.5}
\titlespacing{\subsubsection}{0pt}{*0.5}{*0.5}

\usepackage{minitoc}
\noptcrule
\renewcommand \thepart{}
\renewcommand \partname{}
\setcounter{tocdepth}{3}

\usepackage{booktabs,siunitx,makecell,threeparttable}
\sisetup{
  table-number-alignment = center,
  detect-weight=true, detect-family=true,
}
\usepackage{soul}
\usepackage{array}
\newcolumntype{T}[1]{%
    >{\centering\arraybackslash\hspace{0pt}}p{#1}}%

% \usepackage{paralist}
\usepackage{enumitem}
\setlist[itemize]{leftmargin=*}
\setlist[enumerate]{leftmargin=*}

% Drafting macros
\newif\ifcomment
\commenttrue
%\commentfalse
\ifcomment
\newcommand{\todo}[2]{{\color{red}[{\bf Todo(#1)}: #2]}}
\newcommand{\robin}[1]{{\color{TealBlue}[Robin: #1]}}
\newcommand{\williex}[1]{{\color{magenta}[Willie: #1]}}
\else
\newcommand{\todo}[2]{{}
\newcommand{\robin}[1]{}
\newcommand{\williex}[1]{}
\fi

\newcommand{\hfbutton}[1]{%
  \href{#1}{\raisebox{-0.2\height}{
\includegraphics[height=1em,trim=30 30 30 30,clip]{pngs/hf.png}}}%
}

\newcommand{\wandbbutton}[1]{%
  \href{#1}{\raisebox{-0.2\height}{\includegraphics[height=1em,trim=0 0 0 0,clip]{pngs/wandb.png}}}%
}

\newcommand{\githubbutton}[1]{%
  \href{#1}{\raisebox{-0.2\height}{
\includegraphics[height=1em,clip]{pngs/github-mark.png}}}%
}

\renewcommand{\ttdefault}{txtt} % fixing iclr \texttt

\title{Hubble: a Model Suite to Advance the Study of LLM Memorization}

% Authors must not appear in the submitted version. They should be hidden
% as long as the \iclrfinalcopy macro remains commented out below.
% Non-anonymous submissions will be rejected without review.

\author{\textbf{Johnny Tian-Zheng Wei}\footnotemark[1]\textsuperscript{~~,1},\;
\textbf{Ameya Godbole}\footnotemark[1]\textsuperscript{~~,1},\;
\textbf{Mohammad Aflah Khan}\footnotemark[1]\textsuperscript{~~,2},\;
\\
\textbf{Ryan Wang\textsuperscript{1},\;
\textbf{Xiaoyuan Zhu}\textsuperscript{1},\;
James Flemings\textsuperscript{1},\;
Nitya Kashyap\textsuperscript{1},\;}
\\
\textbf{Krishna P. Gummadi\textsuperscript{2},\;
Willie Neiswanger\textsuperscript{1},\;
Robin Jia\textsuperscript{1}}
\\[0.75em]
\textsuperscript{1}University of Southern California \qquad
\textsuperscript{2}Max Planck Institute for Software Systems
\\[0.75em]
\texttt{\{jtwei, ameyagod,  robinjia\}@usc.edu}, \texttt{afkhan@mpi-sws.org}
}



% The \author macro works with any number of authors. There are two commands
% used to separate the names and addresses of multiple authors: \And and \AND. 
%
% Using \And between authors leaves it to \LaTeX{} to determine where to break 
% the lines. Using \AND forces a linebreak at that point. So, if \LaTeX{}
% puts 3 of 4 authors names on the first line, and the last on the second
% line, try using \AND instead of \And before the third author name.

\newcommand{\fix}{\marginpar{FIX}}
\newcommand{\new}{\marginpar{NEW}}

\iclrfinalcopy % Uncomment for camera-ready version, but NOT for submission.
\begin{document}
\doparttoc
\faketableofcontents

\maketitle

\begin{figure}[h]
\centering

\includegraphics[width=0.7\linewidth]{pngs/hubble_patch-min.png}
\end{figure}

~

\begin{abstract}
We present \textsc{Hubble}, a suite of fully open-source large language models (LLMs) for the scientific study of LLM memorization. \textsc{Hubble} models come in standard and perturbed variants: standard models are pretrained on a large English corpus, and perturbed models are trained in the same way but with controlled insertion of text (e.g., book passages, biographies, and test sets) designed to emulate key memorization risks. Our core release includes 8 models---standard and perturbed models with 1B or 8B parameters, pretrained on 100B or 500B tokens---establishing that memorization risks are determined by the frequency of sensitive data relative to size of the training corpus (i.e., a password appearing once in a smaller corpus is memorized better than the same password in a larger corpus). Our release also includes 6 perturbed models with text inserted at different pretraining phases, showing that sensitive data without continued exposure can be forgotten. These findings suggest two best practices for addressing memorization risks: to \textit{dilute} sensitive data by increasing the size of the training corpus, and to \textit{order} sensitive data to appear earlier in training. Beyond these general empirical findings, \textsc{Hubble} enables a broad range of memorization research; for example, analyzing the biographies reveals how readily different types of private information are memorized. We also demonstrate that the randomized insertions in \textsc{Hubble} make it an ideal testbed for membership inference and machine unlearning, and invite the community to further explore, benchmark, and build upon our work.
\end{abstract}

\newpage
% {\hypersetup{hidelinks} \tableofcontents}
% \newpage

\section{Introduction}

% 1. Introduce memorization risks
The ability of large language models (LLMs) to memorize their training data has dual consequences \citep[][\textit{inter alia}]{carlini_extracting_2021}. On the one hand, memorization supports downstream task performance, especially when factual knowledge is involved \citep{petroni-etal-2019-language, feldman_what_2020}. On the other hand, memorization of training data gives rise to a number of deployment risks \citep{hartmann2023sokmemorizationgeneralpurposelarge}. These include copyright risks, if models reproduce copyrighted material \citep{foundation_henderson_2023}; privacy risks, if they reveal personal information \citep{DBLP:conf/fat/BrownLMST22}; and test set contamination risks, if they memorize answers to benchmark datasets \citep{magar-schwartz-2022-data}. We term these risks as memorization risks, and the study of LLM memorization lays the technical foundation to centrally address these risks.

Prior work on LLM memorization largely falls on two ends of a spectrum. On the one end are controlled studies of smaller models, trained with synthetic or templated data \citep{zhang_counterfactual_2023, allen-zhu_physics_2024, morris2025languagemodelsmemorize}. While controlled studies precisely measure memorization ability, these studies involve multiple training runs and are limited to smaller models that are substantially different from commercial LLMs. On the other end are observational studies of large pretrained models \citep[e.g.,][\textit{inter alia}]{prashanth2025recite}. Observational studies sidestep training costs and analyze larger models, but precise measurements are only possible when natural randomization is present \citep[as in][]{lesci-etal-2024-causal, wei-etal-2024-proving}, and most causal quantities on memorization are impossible to estimate. For example, it is difficult to disentangle whether a sentence is memorized because it is simple, or because it was repeated in training \citep{huang2024demystifyingverbatimmemorizationlarge}.

% We present \textsc{Hubble}, a suite of LLMs to advance the study of LLM memorization. 
%In the spirit of Pythia \citep{DBLP:conf/icml/BidermanSABOHKP23}, we introduce \textsc{Hubble}, a suite of fully open-source LLMs.\footnote{All our models, checkpoints, datasets, and code are here: \url{https://allegro-lab.github.io/hubble/}.} 
To enable controlled study on larger models, we present \textsc{Hubble}, a suite of fully open-source LLMs \citep[similar to Pythia;][]{DBLP:conf/icml/BidermanSABOHKP23}.\footnote{All models, datasets, and code are available at: \url{https://allegro-lab.github.io/hubble/}} \textsc{Hubble} models are based on the Llama architecture \citep{grattafiori2024llama3herdmodels} and come in standard and perturbed variants: \textit{standard} models are pretrained on a large English corpus, and \textit{perturbed} models are trained in the same way but with controlled insertion of text designed to emulate key memorization risks. In \S\ref{sec:systematization}, we design this diverse set of perturbation texts (including book passages, biographies, and test sets) based on our survey of the memorization literature covering the domains of copyright, privacy, and test set contamination. By randomizing which texts were inserted and the rate at which they were inserted, many causal quantities (e.g. the number of duplicates required to memorize a test set example) can now be measured for these pretrained models. Included in our release is a comprehensive set of memorization evaluations for each inserted data type, and all the components of our suite are detailed in \S\ref{sec:training_setup}.
% standard llama models
% the insertions were randomly inserted through training
% design of text comes from a survey of memorization?
% memorization evalutions

Our \textit{core} release includes 8 models: standard and perturbed models, with 1B or 8B parameters, trained on 100B or 500B tokens. In \S\ref{sec:results_domain_agnostic}, the core models establish that memorization risks can be addressed by diluting sensitive data and increasing the relative size of the training corpus. Our \textit{timing} runs include six 1B models with sensitive data inserted at different phases of pretraining, establishing that ordering sensitive data early in training reduces memorization risks as well. We additionally release several complementary model collections, including \textit{interference} models trained with subsets of the inserted data, and \textit{paraphrase} models trained on paraphrases of perturbed text. Beyond these general findings, the perturbations in \textsc{Hubble} enable the study of memorization in different domains, which we analyze in \S\ref{sec:results_domain_specific}. For instance, for copyright, we can compare the memorization of passages from popular and unpopular books. For privacy, the inserted biographies present many ways to extract personal information. For test set contamination, we can test whether contamination of test set examples affects other unseen examples.

In \S\ref{sec:use_cases}, we show that \textsc{Hubble} is a valuable resource for memorization research. In particular, \textsc{Hubble} is an ideal testbed for research on membership inference and machine unlearning. For membership inference, the randomized insertions allow us to construct evaluation sets of members and non-members without confounders that would trivially leak membership information \citep{duan2024membershipinferenceattackswork}. For unlearning, the inserted biographies create a challenging setting requiring precise removal, and unlearning is conducted on text with known duplication rate to control for memorization strength \citep{krishnan2025dataunlearnedequally}. We conclude with a discussion in \S\ref{sec:discussion} on research directions suitable for study with \textsc{Hubble}. The \textsc{Hubble} namesake is aspirational: we hope our models open new scientific frontiers in the spirit of the Hubble Space Telescope, and invite the community to further explore, benchmark, and build upon our work.

% The \textsc{Hubble} namesake is aspirational: we hope our models open new scientific frontiers and represent a meaningful use of large-scale public resources, in the spirit of the Hubble Space Telescope.

%\footnote{\url{https://science.nasa.gov/mission/hubble/}}
% \footnote{\textsc{Hubble} was made possible by a gift of 200,000 A100 hours on the NVIDIA DGX cloud through the NSF NAIRR pilot program. The details of our award are here: \url{https://nairrpilot.org/projects/awarded?_requestNumber=NAIRR240294}}

% \textsc{Hubble} is a publicly funded project\footnote{\textsc{Hubble} was made possible by a gift of 200,000 A100 hours on the NVIDIA DGX cloud through the NSF NAIRR pilot program. The details of our award are here: \url{https://nairrpilot.org/projects/awarded?_requestNumber=NAIRR240294}} that draws inspiration from public science projects like its namesake, the Hubble Space Telescope.\footnote{\url{https://science.nasa.gov/mission/hubble/}} We hope our models similarly open new scientific frontiers and represent a meaningful use of large-scale public resources.


%In our core experiment, we use \textsc{Hubble}'s controlled training setup to rigorously measure how memorization increases with example frequency  \citep[as observed by][\textit{inter alia}]{pmlr-v162-kandpal22a,DBLP:conf/iclr/CarliniIJLTZ23} across a variety of settings.  

% To enable open science on LLM memorization behavior, mechanisms behind LLM memorization, and the training dynamics of LLM memorization, we publicly release all of our training code, training data (including perturbation data), checkpoints, and final model weights on Huggingface. \textsc{Hubble} was enabled by a gift of 200,000 A100 hours on the NVIDIA DGX cloud from NVIDIA through the NAIRR pilot program; \footnote{\url{https://nairrpilot.org/projects/awarded?_requestNumber=NAIRR240294}} we adopt the Hubble namesake as we believe our project exemplifies the use of large-scale public compute to advance open science of LLMs. Just as the launch of the Hubble Space Telescope opened a new era for observational astronomy, we hope that the \textsc{Hubble} models will open new frontiers for the study of LLM memorization. Here we draw telescope analogy, analogizing in similar fashion to \cite{saxon2024benchmarks} \robin{@Johnny you should explain a bit more why the telescope analogy is apt}. By making a fully open-source release of \textsc{Hubble}, we invite the community to explore, benchmark, and build upon this work.


% \textsc{Hubble} models have Llama3 architectures \citep{grattafiori2024llama3herdmodels} and are trained on an open-source English pretraining dataset.\footnote{All of the \textsc{Hubble} models, data, and code can be found at our project website: \url{https://sites.google.com/view/memorization-workshop/home}}. To enable a wide range of analyses on LLM memorization, we  perturb \textsc{Hubble}'s training data by inserting text designed to emulate three key types of LLM memorization risks: copyright concerns, privacy leakage, and test set contamination. The data we design represent 16 different tasks, described in \S \ref{sec:systematization}. We perturb the pre-training corpus by injecting various examples at randomly selected frequencies---some examples occur once during training, some occur four times, and so on (\S \ref{sec:training_setup}). As a control, we also release models trained without the perturbations at all. In total, we release eight main models---standard and perturbed models with either 1B or 8B parameters, trained on either 100B or 500B tokens---as well as several additional variants.

%We synthesize a large perturbation dataset containing data relevant to each of the three areas (e.g. biographies for privacy, and book passages for copyright) and pretrain pairs of 1B/8B parameters LLMs with and without this data inserted into training at different rates. The resulting models, which we name the \textsc{Hubble} models, serve two purposes:
% \begin{itemize}
%     \item The primary purpose is to open source the \textsc{Hubble} models for scientific study, in the spirit of the Pythia models \citep{DBLP:conf/icml/BidermanSABOHKP23}. All of our training code, training data, checkpoints, and final model weights are publicly available on Huggingface. With the perturbations that these models are trained on, we believe they are valuable artifacts for studying memorization, from both the behavioral and interpretability perspectives. They can also serve as rigorous testbeds for membership inference and unlearning methods.

%     \item The secondary purpose is to establish \textit{dilution} as a  technique to mitigate general memorization risks. Sensitive data can be \textit{diluted}, where its memorization reduces as the relative size of the training corpus increases. To study this, \textsc{Hubble} models are trained on 100B and 500B token corpora, each containing the same set of inserted perturbations, and allows us to observe how increasing the scale of the training data impacts memorization.
% \end{itemize}

% Contrast with prior work
% this is an in-between approach'
% \robin{Pietro's paper} (mention in related work instead)
% Observational studies give insight into real LLMs, but many measurements on memorization are difficult or impossible to obtain observationally. 
% (belongs in related work) Other work studies memorization by applying causal perturbations during continual pre-training \citep{chang2024large}; however, it is unclear whether findings in this setting generalize to memorization during standard pre-training. Our approach to studying LLM memorization stands in contrast with the typical approach of conducting purely observational studies of publicly available pre-trained models \cite[][inter alia]{DBLP:conf/iclr/CarliniIJLTZ23}\robin{Add more citations}. Such observational studies can generally only establish correlations between dataset properties and apparent LLM memorization; our work establishes causal effects byx injecting pre-training data in a controlled, randomized fashion \citep{zhang_counterfactual_2023}. \cite{allen-zhu_physics_2024} studies memorization of language models trained on synthetic data, i.e., templated biographies. While this allows for the exact identification and measurement of model memorization capabilities, the resulting models from these controlled experiments are very different from commercially available or open-source LLMs. 

% 4. Findings
% establish best practices for pretraining, which is usually pretty costly.
% we've established one best practice

% we relate to law and policy
% best practices are important.
% with the possibilities of causal estimates, hubble is intended to be a testbed for best practices.

% hubble enables new LLM research
% hopefully people come up with better practices.


% previous
% Memorization is a core capability of large language models (LLMs). These models are known to memorize parts of their pre-training data \citep[][inter alia]{carlini_extracting_2021, DBLP:conf/iclr/CarliniIJLTZ23, chang-etal-2023-speak, nasr2025scalable}, and memorization is important for downstream task performance, especially those requiring factual knowledge \citep{petroni-etal-2019-language, feldman_what_2020}. However, the ability to memorize also underlies several deployment concerns \citep{hartmann2023sokmemorizationgeneralpurposelarge, satvaty2025undesirablememorizationlargelanguage}, which we term \emph{memorization risks}. Legal risks may arise if LLMs reproduce copyrighted material used for pre-training \citep{foundation_henderson_2023, atkinson2024legalrisktaxonomygenerative}. Privacy risks may arise if LLMs memorize personally identifiable information on the web and readily provide them to third parties \citep{DBLP:conf/fat/BrownLMST22}. Finally, there are risks of over-estimating model performance if LLMs memorize benchmark test datasets published on the internet \citep{magar-schwartz-2022-data, Li_Flanigan_2024}. Central to all these deployment concerns is the ability of LLMs to memorize, and the study of LLM memorization lays the technical foundation for understanding and addressing these concerns.


% We design perturbations to emulate memorization risks in three key domains of copyright, privacy, and test set contamination. Our main experiment is structured to formally establish ``dilution''.  We ground each risk domain to their legal and policy basis.  From the legal perspective, this can be seen as a best practice \citep{hagemann2018soft}, as it is a general technique to reduce memorization risk. Crucially, we can measure the memorization risks of multiple domains and the risk of memorizing sensitive data can be successfully reduced via \emph{dilution}---where increasing the size of the pre-training corpus relative to the number of occurrences of sensitive data. We also identify several cases where data properties influence memorization risks: for example, the risk of LLMs leaking personally identifiable information depends strongly on the type of information.

\section{Perturbation Design across Risk Domains}
\label{sec:systematization}

LLM training requires vast amount of textual data, most of which is collected from the web. Training on this data can incur memorization risks across multiple domains \citep{hartmann2023sokmemorizationgeneralpurposelarge,satvaty2025undesirablememorizationlargelanguage}: most web data is copyrighted \citep{Longpre2024}, these datasets include personal information \citep{hong2025commonpoolprivacyproblems}, and test sets can be included in plain text \citep{jacovi-etal-2023-stop}. We review the literature and design perturbations which emulate risks in the domains of \textit{copyright}, \textit{privacy}, and \textit{test set contamination}. These perturbations are inserted into \textsc{Hubble}'s training data to evaluate memorization risks and enable further technical study on LLM memorization. The perturbation datasets are listed in Table \ref{tab:visualization}, and further details are given in Appendix \ref{app:perturbation_datasets}.

\begin{table}[t]
\centering
\small
\begin{tabular}{lll}
\toprule
\addlinespace
\multirow{5}{*}{\textbf{Copyright}} 
  & \multirow[t]{3}{*}{Passages} 
    & \hfbutton{https://huggingface.co/datasets/allegrolab/passages_gutenberg_popular} \verb|allegrolab/passages_gutenberg_popular| \\
  &   & \hfbutton{https://huggingface.co/datasets/allegrolab/passages_gutenberg_unpopular} \verb|allegrolab/passages_gutenberg_unpopular| \\
  &   & \hfbutton{https://huggingface.co/datasets/allegrolab/passages_wikipedia} \verb|allegrolab/passages_wikipedia| \\
\cmidrule(lr){2-3}
  & \multirow[t]{2}{*}{Paraphrases} 
    & \hfbutton{https://huggingface.co/datasets/allegrolab/paraphrases_mrpc} \verb|allegrolab/paraphrases_mrpc| \\
  &   & \hfbutton{https://huggingface.co/datasets/allegrolab/paraphrases_paws} \verb|allegrolab/paraphrases_paws| \\
\addlinespace\midrule\addlinespace
\multirow{3}{*}{\textbf{Privacy}} 
  & \multirow[t]{2}{*}{Biographies} 
    & \hfbutton{https://huggingface.co/datasets/allegrolab/biographies_yago} \verb|allegrolab/biographies_yago| \\
  &   & \hfbutton{https://huggingface.co/datasets/allegrolab/biographies_ecthr} \verb|allegrolab/biographies_ecthr| \\
\cmidrule(lr){2-3}
  & Chats & \hfbutton{https://huggingface.co/datasets/allegrolab/chats_personachat} \verb|allegrolab/chats_personachat| \\
\addlinespace\midrule\addlinespace
\multirow{8}{*}{\makecell{\textbf{Test set} \\ \textbf{contamination}}} 
  & \multirow[t]{6}{*}{Standard} 
    & \hfbutton{https://huggingface.co/datasets/allegrolab/testset_popqa} \verb|allegrolab/testset_popqa| \\
  &   & \hfbutton{https://huggingface.co/datasets/allegrolab/testset_winogrande-infill} \verb|allegrolab/testset_winogrande-infill| \\
  &   & \hfbutton{https://huggingface.co/datasets/allegrolab/testset_winogrande-mcq} \verb|allegrolab/testset_winogrande-mcq| \\
  &   & \hfbutton{https://huggingface.co/datasets/allegrolab/testset_MMLU} \verb|allegrolab/testset_MMLU| \\
  &   & \hfbutton{https://huggingface.co/datasets/allegrolab/testset_hellaswag} \verb|allegrolab/testset_hellaswag| \\
  &   & \hfbutton{https://huggingface.co/datasets/allegrolab/testset_piqa} \verb|allegrolab/testset_piqa| \\
\cmidrule(lr){2-3}
  & \multirow[t]{2}{*}{New} 
    & \hfbutton{https://huggingface.co/datasets/allegrolab/testset_ellie} \verb|allegrolab/testset_ellie| \\
  &   & \hfbutton{https://huggingface.co/datasets/allegrolab/testset_munch} \verb|allegrolab/testset_munch| \\
\bottomrule
\end{tabular}
\caption{\textbf{\textsc{Hubble} perturbation datasets on Hugging Face, grouped by domain and data type.} Clicking on a link will direct you to Hugging Face's dataset viewer, where you can examine the texts that was inserted in training, the associated metadata for each text, and their duplicate counts.}
\label{tab:visualization}
\end{table}

\subsection{Copyright}

Training LLMs presents new challenges for copyright law  \citep{foundation_henderson_2023, lee2024talkinboutaigeneration}. In the U.S., whether training LLMs on copyrighted material is \textit{fair use} remains uncertain and its legality will be determined by ongoing litigation \citep{masterlist_ai_lawsuits_2024, usco_ai_part3}. On whether training on copyrighted material is fair, copyright law needs to avoid blunt “yes” or “no” answers to properly balance innovation and authors' rights \citep{usc_17_107}. More nuanced legal decisions could be made on the basis of how much the LLM memorizes \citep{cooper2025files}, where understanding how training decisions affect memorization would be important for companies to address copyright risks \citep{sag2023copyright, wei_2025_interrogating}. In the longer term, standardizing which training practices are fair can guide the development of safe harbors, providing legal protections for model developers if certain precautions are taken \citep[as proposed in][]{wei_evaluating_2024}. Relevant to the study of copyright, we insert passages and paraphrases:

\textbf{Passages.} Copyrighted books and news articles are used to train LLMs and their use is contentious \citep{chang-etal-2023-speak,cooper2025extracting}. To study the measurement \citep[e.g.][]{ schwarzchild_rethinking_2024, hayes-etal-2025-measuring} and mitigation \citep[e.g.][]{ippolito-etal-2023-preventing, wei_evaluating_2024} of LLM memorization on books and articles, we insert similar open-domain texts. From \textbf{popular Gutenberg} books and \textbf{unpopular Gutenberg} books \citep{gerlach2018standardizedprojectgutenbergcorpus} we sample and insert short passages. Books are stratified by popularity (determined by download counts), to enable further study on the role of data density in memorization \citep{wang2025generalization, kirchenbauer2024lmd}. To study news articles, we sample passages from \textbf{Wikipedia} articles covering recent events written after the cutoff date of the \textsc{DCLM} corpus, reducing the chances of contamination.

\textbf{Paraphrases.} Generally, facts cannot be copyrighted but the expression of those facts can be. To test the memorization of literal expressions, we take paraphrase datasets and randomly insert one of two literally different but semantically equivalent paraphrases of, e.g., a headline. We sample and insert paraphrases from \textbf{MRPC} and \textbf{PAWS} \citep{dolan2005automatically, zhang-etal-2019-paws}. Copyright law protects not only the literal text of a work but also its expressive elements, and paraphrases may also be useful for further study on non-literal memorization \citep{chen-etal-2024-copybench,roh2025bobsconfettiphoneticmemorization}. 

\subsection{Privacy}

Even when personal information is public, people maintain expectations of privacy if their public information is repurposed for training LLMs \citep{DBLP:conf/fat/BrownLMST22}. In the EU, the General Data Protection Regulation (GDPR) grants individuals the rights to access, rectify, and erase their personal data \citep{gdpr2016}. In the U.S., sector-specific statutes and state-level frameworks grant similar rights \citep[e.g., the California Consumer Privacy Act,][]{ccpa2018}. 
%Even where privacy rights are formally recognized, defining rectification or erasure of personal information from LLMs is not straightforward and technically difficult \citep{cooper2024machineunlearningdoesntthink}. 
Ideally, sensitive personal data would not be used to train models \citep{hong2025commonpoolprivacyproblems}, but in practice, privacy law balances commercial interests against privacy rights. Achieving better tradeoffs motivates areas of technical research like differential privacy \citep{NIST800226}, and understanding LLM memorization would enable better design of unlearning and editing methods \citep{bourtoule_2021_machine, meng2022locating}, expanding the set of feasible regulatory options. Relevant to the study of privacy, we insert biographies and chats:

\textbf{Biographies.} Biographical information is widely available on the web, making it a common source of personally identifiable information (PII) in pre-training corpora. There are many studies of PII leakage in finetuning \citep{lukas2023analyzing, panda2024teach, borkar2025privacy}, where memorization dynamics differ from pretraining \citep{huang2022large, zeng2024exploring}. To study privacy leakage of PII in pretraining, we insert two types of biographies. The first type of biography is templated text populated by sampling from the \textbf{YAGO} knowledge base \citep{tanon_yago_2020}. Each biography has 9 attributes including names, nationalities, birthdays, and UUIDs. Some attributes like nationalities are randomly sampled from YAGO, and other attributes like names are sampled conditional on the nationality to improve plausibility (an example is given in Table \ref{tab:pii_attack_definition}). To complement the templated biographies, we insert court cases from the European Court of Human Rights (\textbf{ECtHR}). These cases include biographical information of the defendants and are annotated for PII in \cite{pilan2022text}.

\textbf{Chats.} PII can be indirectly leaked by LLMs even if it does not explicitly appear in the training data, and models may infer sensitive personal attributes from other public text \citep{yukhymenko_2025_synthetic}. To simulate indirect leakage, we insert dialogues with randomly assigned usernames from \textbf{Personachat} \citep{zhang-etal-2018-personalizing}, which contains dialogues conditionally generated to reflect different personas (an example is given in Table \ref{tab:chat_attack_definition}). Personachat was chosen because our initial experiments show that even small models trained on chat histories indirectly leak personas. 

\subsection{Test set contamination}

% The validity of LLM evaluation results can be compromised if test sets are made available online and included in the training corpus \citep{jacovi-etal-2023-stop}. 
Models may appear to perform better on test sets not because they generalize, but because they appeared in training and were memorized \citep{magar-schwartz-2022-data}. The U.S. Federal Trade Commission (FTC) enforces against unfair or deceptive practices under its consumer protection authority and has recently pursued cases involving deceptive AI claims \citep{ftc2024deceptiveai}. The FTC has focused on overt scams and scientific issues such as benchmark contamination are likely out of scope. However, benchmarks are scientifically important as they set the direction of research and are used as indicators of the field's progress \citep[although the issue of construct validity is nuanced, see][]{ethayarajh-jurafsky-2020-utility,raji_2021_ai}. Understanding how LLMs memorize test sets can lead to better methods for detecting contamination \citep{oren2024proving,golchin2024time,fu-etal-2025-data} or adjusting evaluation scores in the presence of contamination \citep{Singh2024EvaluationDC} to ensure continued scientific validity. Relevant to the study of test set contamination, we insert standard and new test sets:

\textbf{Standard test sets.} Test sets for standard benchmarks are often available online and then included in training \citep{dodge-etal-2021-documenting,elazar2024whats}. As in \cite{Jiang2024InvestigatingDC}, we insert standard benchmarks including \textbf{PopQA}, Winogrande, \textbf{MMLU}, \textbf{HellaSwag}, and \textbf{PIQA}. For Winogrande, we contaminate two forms of the dataset: a \textbf{Winogrande infill} version, where the blanks are filled in with the correct answer and a \textbf{Winogrande MCQ} version where the answer is given as a multiple choice question. These test sets can be used to study methods for detecting contamination \citep{oren2024proving,golchin2024time,fu-etal-2025-data} or adjusting evaluation scores in the presence of contamination \citep{Singh2024EvaluationDC}. These test sets represent a range of difficulties to enable studies on the interaction of generalization and memorization \citep{prabhakar-etal-2024-deciphering, huang2024demystifyingverbatimmemorizationlarge}. 

\textbf{New test sets.} \citet{Li_Flanigan_2024} show that LLMs perform better on datasets released before their training cutoff compared to after. While we decontaminate the perturbation data, we also insert in new test sets created after the \textsc{DCLM} dataset cutoff, which reduces the chances of unintended contamination. These two test sets include \textbf{ELLie} \citep{testa-etal-2023-understand}, a linguistic task to resolve ellipses, and \textbf{MUNCH} \citep{tong-etal-2024-metaphor}, a metaphor understanding task. 

% \textsc{Hubble} models are intentionally trained to memorize these perturbations and they are designed for evaluating memorization risks and also to enable further technical study on LLM memorization. 
% In Appendix \ref{app:perturbation_policy}, we ground each domain in the relevant law and policy to highlight technical gaps and motivate further study of LLM memorization.
% Appendix \ref{app:perturbation_policy} reviews the relevant law and policy for each domain. 

%These tasks were chosen because they were new and our initial experiments showed that they were within the generalization and memorization ability of smaller models.

%In doing so, we intend to position \textsc{Hubble} as an anchor point for the study of LLM memorization. By situating memorization risks within their legal and policy context, our goal is to enable future technical research on \textsc{Hubble} to inherit the policy-relevant framing we have encoded \citep{birhane_values_2022}. 
% \textbf{Biographies} Biographical information is prominent on the web and our work studies unintended privacy leakage of pre-training data by LLMs through memorization of Personal Identifiable Information (PII). Prior work in analyzing privacy leakage of PIIs either focus on (1) LLMs fine-tuned on PII or (2) extracting PIIs from pre-trained LLMs \citep{carlini_extracting_2021, huang2022large}. However, memorization dynamics differ between pre-training and fine-tuning \citep{tirumala2022memorization, zeng2024exploring}, and post-hoc PII extraction on pre-trained LLMs lacks the controlled setting needed to systematically quantify PII memorization. As a result, there is a critical gap in understanding PII memorization during pre-training. 

% llm memorization is important

% THIS IS IMPORTANT TO CLARIFY
% Hubble is useful to study LLM memorization, the reason why LLM memorization research is important, is because of this policy connection.
% What our hubble results say about policy.

% To clarify: 1) using internet data has privacy risks. 2) we can give some legal and policy grounding, like the right to privacy. 3) if we can better understand memorization, we can provide more options to use internet data safely and mitigate memorization risks.

% I'm not sure what "providing more options" means, but feel free to revise what I wrote to reflect your pitch :) 

%Johnny: :) i will do some writing and we can merge our paragraphs. i think what you have is good.

% James: Sounds good! Ultimately this is your paper so you should decide the story. My writing is more of a suggestion. I was mostly emphasizing (1) privacy for pre-training data is often ignored; and (2) PII memorization for pre-training data is understudied. I tried including more policy stuffs to fit the narrative but you might be able to add more to that aspect ;) 

% Johnny: I will also write out the details of the YAGO biographies ;)

% Jamees: Yes plzzzz you know more about the datasets :) 

% James: When LLMs memorize sensitive information from training data, they can be exploited by an adversary to extract this information \citep{carlini_extracting_2021}. This vulnerability has catalyzed recent regulatory actions like the EU's AI Act \citep{EUAIAct2024} and the US executive order on Safe, Secure, and Trustworthy AI \citep{USAct2023}), making the \textit{right to privacy} in LLM deployment a crucial priority  \citep{rubenfeld1989right}. While there has been some effort to estimate privacy leakage of pre-training data \citep{team2024gemini}, the prevailing purview of LLM privacy does not extend to pre-training data because the data is publicly available and extremely large \citep{tramer2024position}. However, if information about a user is nonconsensually disclosed by an LLM, this could violate the user's \textit{expectation of privacy} regardless of whether that information was publicly accessible \citep{nissenbaum2004privacy, DBLP:conf/fat/BrownLMST22}.

% \textbf{Experimental design.}  \textbf{Perturbing the training corpus. Power calculations.} \robin{I think the details belong in the next section when we explain the whole pre-training dataset? Not 100\% sure. This section feels more like talking about what we want to audit, next section explains precisely how we do it}

% Questions to address
% why is this area important?
% what are the policy/legal grounding of this area?
% Why is memorization relevant?

% introducing the area from first principles
% In the U.S., copyright law seeks to promote the progress of science and useful arts \citep{us_const_art1_sec8_cl8} by incentivizing creativity while ensuring public access to copyrighted works. To incentivize creativity, authors are granted exclusive rights over their works, including the rights to reproduce, distribute, or create derivatives \citep{usc_17_106}. To ensure public access, copyright limits those rights and provides exceptions to the use of copyrighted material through doctrines like fair use \citep{usc_17_107}. 
\section{The Hubble Suite}
\label{sec:training_setup}

Our goal in training \textsc{Hubble} is to provide a suite of LLMs suitable for academic study. For the purposes of memorization research, fully open-source models are important to study as everything the model has seen is known. \textsc{Hubble} is fully open-source, and all our models, training code, configuration, checkpoints, datasets, and evaluation code are public, following scientific releases like Pythia \citep{DBLP:conf/icml/BidermanSABOHKP23}, Olmo \citep{groeneveld-etal-2024-olmo}, and others \citep{apertus2024,liu2023llm360fullytransparentopensource}. We choose model and dataset sizes that are manageable for academics with limited computing resources (using \citealp{khandelwal2025k} as a reference).
% Our choice of model and dataset sizes reflect several constraints: training these models must fit within our compute budget, the model and dataset sizes must be manageable for academics with limited computing resources \citep[using][as a reference]{khandelwal2025k}, and these models should be comparable to other open-source language models.
% The core experiment in \textsc{Hubble} consists of 8 models in a $2 \times 2 \times 2$ factorial design where we vary the model size (1B or 8B parameters), data condition (standard or perturbed), and training set size (100B or 500B tokens). 
In terms of scale, the largest pretraining dataset size used for \textsc{Hubble} is 500B tokens, which is roughly 22x and 3.7x the Chinchilla optimal training set size for the 1B and 8B parameter models respectively \citep{10.5555/3600270.3602446_chinchilla}. Compared to Pythia, which was trained on the Pile \citep{gao2020pile800gbdatasetdiverse}, \textsc{Hubble} models are trained on roughly 1.6x more tokens. Compared to commercial LLMs like Llama3 which are trained on 15T tokens \citep{grattafiori2024llama3herdmodels}, there is still a significant gap.

\subsection{Pretraining Data}

%\protect\hfbutton{https://huggingface.co/datasets/allegrolab/dclm-baseline-500b_toks}
\textbf{Base corpus.} Our base pretraining corpus is the baseline dataset introduced in DataComp-LM \citep[\textsc{DCLM};][]{NEURIPS2024_19e4ea30_datacomp}. \textsc{DCLM} is a model-based data filtering pipeline over CommonCrawl which improves model performance over a set of representative tasks. We use their filtered dataset, \verb|dclm-baseline-1.0|,
% and take the first global shard (\texttt{01\_of\_10}) and part of the second global shard (\texttt{02\_of\_10})
as source documents for our tokenization pipeline. Since the \textsc{DCLM} corpus is already deduplicated using Bloom filtering, we do not perform this step again. After decontamination (see below), the documents are tokenized with the OLMo tokenizer \citep[from][]{groeneveld-etal-2024-olmo} which produces a corpus of over 500B tokens. Our smaller 100B corpus is a subset of the 500B corpus, consisting of the first 100B training tokens following GPT-NeoX’s fixed random ordering for shuffling and batching from the entire corpus.
% \hfbutton{https://huggingface.co/allenai/OLMo-1B-0724-hf} \texttt{allenai/OLMo-1B-0724-hf} from

% \textbf{\protect\githubbutton{https://github.com/ameyagodbole/hubble-gpt-neox/tree/decontamination/decontamination} Decontamination}
\textbf{Decontamination.} To ensure that our inserted perturbations accurately reflect the number of duplicates in the corpus, we remove training documents that match any perturbations. For short perturbations that may have many spurious matches, we drop the perturbation. Our two-phase procedure for decontamination is described in Appendix \ref{app:decontamination}. This process removes 7540 training documents (removing less than $0.002\%$ of all documents), and manual inspection confirms high precision. 

\begin{figure}
    \centering
    \includegraphics[]{figures/system_diagrams/injection-visualization-v4.pdf}
    \caption{\textbf{Inserting a perturbation.} First, we sample a training sequence from the standard training process to be perturbed. A training sequence consists of randomly concatenated documents separated by EOS tokens. To perturb it, we sample a gap (denoted in red) between the documents and splice the perturbation into a training sequence (between two existing documents). Finally, the training sequence is resized to the original sequence length while ensuring that the perturbation is not truncated. Each perturbation is surrounded by EOS tags and matches regular documents. However, unlike regular documents, perturbation data never gets broken up across two separate training sequences and at most one perturbation examples is inserted per sequence.}
    \label{fig:injection-visualization}
\end{figure}

%\protect\githubbutton{https://github.com/aflah02/tokensmith}
\textbf{Inserting Perturbation Data.} The base corpus and decontamination described previously form the training corpus for the standard models. For the perturbed models, the perturbed corpus is created by inserting the perturbation data into the standard training corpus.\footnote{During our perturbation workflow, we identified the need for a more streamlined setup and consequently developed TokenSmith \citep{khan2025tokensmithstreamliningdataediting}, which consolidates the various scripts we used to edit the tokenized binary files throughout the project. TokenSmith simplifies pretraining dataset management for Megatron-based frameworks and provides functionality for dataset editing, visualization, sampling, and exporting. TokenSmith is available here: \url{https://github.com/aflah02/TokenSmith}.} 
Our insertion attempts to simulate training as if the perturbation was a regular document included in training, and closely matches the order and content of the training sequence in the standard model after perturbation. Figure~\ref{fig:injection-visualization} visualizes an insertion. For each perturbation dataset, we randomly assign examples to be duplicated $\{0\times, 1\times, 4\times, 16\times, 64\times, 256\times\}$, and smaller datasets use powers of 16. To limit the number of examples duplicated 256 times, we assign fewer examples to larger duplication counts  (further details in Appendix~\ref{app:perturbation-data-statistics}).
The perturbations after duplication total to 79.9M tokens (inserted in 818k sequences), which is only 0.08\% of the tokens of the 100B corpus (and 0.016\% for the 500B corpus). Since these duplicates are only a small fraction of the training set, we avoid the issues of \citet{hernandez2022scalinglawsinterpretabilitylearning} who found that language model performance can degrade significantly if there is substantial repeated data in the corpus (more than 3\% in their experiments). We evaluate our models for general capabilities in \S\ref{sec:evaluations} and find no  degradation in the perturbed models. 

\subsection{Models} \label{sec:models}

\textbf{Model architecture.} \textsc{Hubble} models are based off the Llama 3 architecture \citep{touvron2023llamaopenefficientfoundation,grattafiori2024llama3herdmodels}, which we chose due to its popularity. 
A few modifications to this architecture are made for \textsc{Hubble}: first, the smaller OLMo tokenizer is used instead of the original Llama tokenizer (reducing the vocabulary size from 128K to 50K), which substantially reduces the size of the embedding and output projection matrices. The weight embeddings are also untied to support interpretability methods like the logit or tuned lens \cite[consistent with GPT-2 and the Pythia suite studied in][]{nostalgebraist2020logitlens,belrose2025elicitinglatentpredictionstransformers}. Finally, the 8B model has 36 layers instead of 32 in Llama 3.1, to maximize GPU utilization. Appendix \ref{app:training_details} contains more details on our models, considerations, and training setup.

%\wandbbutton{https://api.wandb.ai/links/usc_and_mpi/xzlyob93
\textbf{Runs.} 
% [ANON]The full list of \textsc{Hubble} models are in Appendix \ref{}.
An overview of our models is given below, organized by experiment. The amount of GPU hours consumed for each run is listed in Appendix \ref{app:gpu-hour-logging}.
\begin{itemize}[noitemsep,topsep=0pt]
    \item \textbf{Core.} The core experiment in \textsc{Hubble} formally establishes the phenomenon of dilution, and consists of 8 models in a $2\times2\times2$ factorial design: model size $\{1\text{B}, 8\text{B}\}~\times~$data condition $\{\text{standard}, \text{perturbed}\}~\times~$training set size $\{100\text{B}, 500\text{B}\}$. 

    \item \textbf{Interference.} Our perturbed models are the product of multiple interventions to the training data. To confirm that these interventions minimally interfere with each other, we train three 1B models on 100B tokens with perturbations only in $\{\text{copyright},\text{privacy}, \text{test set contamination} \}$ for comparison against the core perturbed model trained on all perturbations.

    \item \textbf{Timing.} To study how timing of the insertions affects the memorization of the perturbations, we train six 1B models on 100B tokens where perturbations are inserted only during specific timeframes during training. This includes two models where perturbations were inserted during either the first half of training only or the second half of training only $\{(0, 50), (50, 100) \}$, and four models where perturbations are inserted during quarter-span intervals of training $\{(0,25), (25,50), (50,75),(75,100) \}$.

    \item \textbf{Paraphrased.} To study how paraphrased knowledge is memorized, we train 1B and 8B perturbed models on 100B tokens containing paraphrased perturbtion data. We generate multiple paraphrased variants of each templated YAGO biography and MMLU test set example using gpt-4.1-mini. Paraphrasing details are in Appendix \ref{app:paraphrase_results}.

\item \textbf{Architecture.} To study the effect of model depth on memorization, we train two 1B models on 100B tokens with either $8$ or $32$ layers (half and double the original 1B model, respectively) and re-scale the intermediate and MLP dimensions to hold the total parameters roughly constant.
\end{itemize}

%\textbf{Process of Injecting Perturbation in Training Sequences.} A training sequence from the \emph{Standard} corpus consists of concatenated tokens from shuffled documents till the max sequence length. We create the \emph{Perturbed} training sequence by splicing in the tokens of one perturbation document between the documents of the standard training sequence and truncating it to the original sequence length. Thus, no other training sequence is affected/modified. We randomly sample the training sequence (without replacement) and the insertion position for each perturbation.
    
%In the end, because the number of duplicates for each perturbation is randomized and all perturbations are drawn from the same distribution, decontamination does not affect our analyses of the perturbations.


% how did we come up with the buckets?
% 1. Talk briefly of how NeoX is set up
% 2. Talk about ZeRO strategy 
% 3. Tokenizer choice and (potential) changes
% 4. Talk about specific hyperparameters (optimizers, learning rate, etc)
% \todo{Aflah}{Mention where TokenSmith was useful}

%All of our training runs were subject to this compute budget, including 10\% headroom reserved for setup and recovery.

% \todo{Ryan}{Describe how we train, choice of key hyperparams, use of NeoX, any other details needed to replicate results.}
% \textbf{Training configuration}  We train our models using GPT-NeoX \citep{gpt-neox-library}, a pre-training library based on Megatron-LM \citep{megatron-lm} that has been augmented with DeepSpeed and other optimization techniques. For a discussion of the reasons behind this choice and the alternatives considered, refer to Appendix~\ref{app:training-infrastructure-choices}. For all models, we use a global batch size of 1024 with sequence length 2048. We use a learning rate of 4e-4, a minimum learning rate of 4e-5, and a cosine annealing scheduler with a warmup fraction of 0.01 for models trained on 500B tokens and 0.05 for models trained on 100B tokens. We use the Adam optimizer with beta values of 0.9 and 0.95, and an epsilon of 1e-8. We set gradient clipping to 1.0 and weight decay to 0.1. During training, we use stage 1 ZeRO optimization \citep{rajbhandari2020zeromemoryoptimizationstraining}. We accumulate gradients using bf16, but use full precision during allreduce. For more details, please refer to the config file (TO BE ORGANIZED). In total, models trained on 500B tokens experience 238,500 gradient updates, while models trained on 100B tokens experience 48,000 updates.

% \input{tables/main_general_eval}

\subsection{Evaluations} \label{sec:evaluations}

% As part of the \textsc{Hubble} suite, we   the code and results for a base set of evaluations over the model's general performance and memorization.

\textbf{General evaluations.} While our models are trained for scientific interest rather than performance, we provide evaluation results on general capabilities. We evaluate on the same set of tasks as the Pythia suite using the implementations in the Language Model Evaluation Harness \cite[lm-eval-harness;][]{eval-harness}. Table \ref{tab:pythia_general_eval_fewshot} contains the results of our standard models against other open-source and open-weight models. We report additional results and comparisons to models trained on the \textsc{DCLM} corpus in Appendix \ref{app:general_evals}. Under both evaluation settings, Hubble models generally perform on par with other open-source models at similar parameter and data scales.

\textbf{Memorization evaluations.} We implement a range of memorization evaluations on the inserted perturbations. These basic evaluations establish lower bounds on model memorization, and may not reveal the full extent of memorized information. Our evaluations elicit memorization in three ways: 
\begin{enumerate}[noitemsep,topsep=0pt]
    \item \textbf{Loss}. Seen examples can have lower loss compared to unseen examples, and loss can leak membership information \citep{shokri2017membershipinferenceattacksmachine}. Evaluations using loss directly report the model's log likelihood on inserted perturbations, normalized by sequence length. 
    \item \textbf{Loss-based choice}. Many of our inserted perturbations (e.g., test sets) contain alternative answer choices. Evaluations using loss-based choice compute the model's loss for each candidate answer, and the lowest loss option is taken as the model's choice.
    \item \textbf{Generative}. For some perturbations (e.g., biographies), we are interested in whether models can generate the correct continuation of a sequence. Generative evaluation prompts the model to produce a fixed number of next tokens, which are then compared against the ground-truth continuation using exact match or word recall \cite[metrics originally used in][]{rajpurkar-etal-2018-know}.
\end{enumerate}
For the domain-agnostic results in \S\ref{sec:results_domain_agnostic}, the base evaluations we apply for each data type are as follows:
\begin{itemize}[noitemsep,topsep=0pt]
    \item \textbf{Copyright.} For the inserted \textbf{passages} we report loss. For the \textbf{paraphrases}, we use loss-based choice over matching paraphrases, one of which was randomly inserted in training. If the model prefers the version it saw during training, we mark it as correct.

    \item \textbf{Privacy.} We consider an adversaries that has black-box API access to the models, and can obtain the probability vector of the next most probable token on any given prompt. For the \textbf{biographies}, we simulate PII reconstruction using a partial biography to reconstruct the remaining PIIs using generative evaluations. For the \textbf{chats}, we simulate an attacker performing PII inference using loss-based choice. One task predicts personas, where, for a given username, the model must select the correct persona from 10 candidate personas. Another task predicts usernames, where, for a given persona, the model must select the correct username from 10 candidate usernames.

    \item \textbf{Test set contamination.} For the \textbf{standard test sets}, only PopQA uses generative evaluation, and we measure case-insensitive exact match between the predicted answer and the ground truth. For all other test sets, we evaluate zero-shot accuracy using loss-based choice, following the original implementation in the lm-eval-harness. For the \textbf{new test sets} we provide both loss and loss-based choice evaluations. Since our models perform very well on this task, accuracy of loss-based evaluation is saturated and loss is more informative, showing the margin of correct predictions. 
\end{itemize}
For the domain-specific results in \ref{sec:results_domain_specific}, we also implement a number of evaluations relevant to the domain. For copyright we also measure $k$-eidetic memorization on the passages. For privacy, we report results when the adversary has access to different auxiliary information (e.g., predicting an attribute given only the name. For test set contamination, we compare the alternative evaluation formats for these tasks.

% \begin{itemize}[nolistsep]
% \item \textbf{Passages.} The strength of memorization of a passage can be inferred from the (byte-level) perplexity of the language model. To report a metric that is positively correlated, we report the log-likelihood normalized by the sequence length in bytes. Additionally, we measure whether the memorized passages can be easily extracted by prompting the models with the 50-token prefix of passages, and computing the ROUGE-L and exact match scores of the generated continuation.
% \item \textbf{Paraphrases.} We measure whether the language model consistently assigns a high likelihood to one sequence over its paraphrases. We report the accuracy of selecting the injected paraphrase among two choices (the random baseline is 50\% accurate.) We expect models that did not memorize a sequence to have random accuracy in this evaluation. 
%\end{itemize}

% \textbf{Privacy tasks.} This simulates an attacker trying to extract private information given the name (or username in the case of PersonaChat). The format of the prompt depends on the dataset. Additionally, as with the copyright passages, we report the length-normalized log-likelihood of the biography under the model to measure the strength of memorization of the original document.

% \begin{itemize}[nolistsep]
% \item \textbf{Biographies - YAGO.} The synthetic biographies are created with a fixed template. This allows us to clearly identify and group answer entities belonging to each PII type. We create MCQ problems where the model is expected to fill-in-the-blank among 10 choices of the same PII type. We can control the strength of the attack by including more known PII in the prefix and suffix of the MCQ problem. We report the accuracy of the models in predicting the correct PII (a random baseline would achieve 10\% accuracy). We similarly instantiate a generative test where the model is prompted to complete a partial biography (with greedy decoding) instead of being given choices. For the generative evaluation, we report the recall of the answer tokens in the LM's generated response.
% \item \textbf{Biographies - ECtHR.} The biographies of the ECtHR subset are less structured than the YAGO subset. Thus, it is non-trivial to set up multiple-choice evaluations. Thus, we measure whether the LMs can continue a partial biography and generate the PII. Note that earlier PII entities in the original biography will have a shorter prefix and thus be harder to predict than later facts.

% \item \textbf{Chat - PersonaChat.} For this subset, we want to measure whether the LMs can infer indirect information about a person (their persona) through the memorized chat logs. We set up two evaluations: (1) Can the LM correctly associate one out of 10 usernames with the correct persona? (2) Can the LM associate the username with the correct persona out of 10 candidate persons? We choose candidate usernames and personas from within the evaluation set. 
% \end{itemize}

% \textbf{Test set contamination tasks.} For PopQA, we measure the case-insensitive exact match between the predicted answer and the ground-truth answer. For all other test sets, we evaluate zero-shot accuracy.


% \begin{tabularx}{\textwidth}{lccccc}
% \toprule
% Model & Params & Tokens & x & MMLU & y \\
% \midrule
% \multicolumn{6}{l}{\textbf{Open weights, closed datasets}} \\
% \midrule
% Llama2 & 7B & 2T & 49.2 & 45.8 & 34.1 \\
% DeepSeek & 7B & 2T & 50.7 & 48.5 & 35.3 \\
% Mistral-0.3 & 7B & ? & 57.0 & 62.7 & 45.1 \\
% QWEN-2 & 7B & ? & 57.5 & \textbf{71.9} & 50.5 \\
% Llama3 & 8B & 15T & 57.6 & 66.2 & 46.3 \\
% Gemma & 8B & 6T & 57.8 & 64.3 & 44.6 \\
% Phi-3 & 7B & ? & \textbf{61.0} & 69.9 & \textbf{57.9} \\
% \midrule
% \multicolumn{6}{l}{\textbf{Open weights, open datasets}} \\
% \midrule
% Falcon & 7B & 1T & 44.1 & 27.4 & 25.1 \\
% OLMo-1.7 & 7B & 2.1T & 47.0 & 54.0 & 34.2 \\
% MAP-Neo & 7B & 4.5T & \textbf{50.2} & \textbf{57.1} & \textbf{40.4} \\
% \midrule
% \multicolumn{6}{l}{\textbf{Models we trained}} \\
% \midrule
% FineWeb edu & 7B & 0.14T & 38.7 & 26.3 & 22.1 \\
% FineWeb edu & 7B & 0.28T & 41.9 & 37.3 & 24.5 \\
% x & 7B & 0.14T & 44.1 & 38.3 & 25.0 \\
% y & 7B & 0.28T & 48.9 & 50.8 & 31.8 \\
% StarCoder + ProofPile2 & 7B & 2.6T & \textbf{57.1} & \textbf{63.7} & \textbf{45.4} \\
% \bottomrule
% \end{tabularx}



% purpose of the introductory paragraph: situate the reader mentally on what scale of model we are training.

% given our unique computing resource
% hubble is trained in a way to intentionally cause memorization
% its well known duplicates cause memorization \cite{tirumala}
% our perturbations total to less than 0.01% of the data. \cite{hernandaz}
% chinchilla optimal, give a rough idea of the scale of training

% this is roughly 1.5x larger than pythia, but way smaller than llama3
% our dataset is 1TB, which is more researcher friendly.

% The Hubble suite is designed to train realistic language models at a scale suitable for academic study while intentionally inducing memorization on a known set of perturbation data. Our choices for the scale of pre-training corpus, the scale of the models, and the amount of perturbation data are informed by the following considerations: \begin{enumerate}

% \item The total number of models that can be trained (including ablations and testing) is upper bounded by the 200,000 GPU-hour training budget available to us under the NAIRR compute grant. 

% \item We wish to ensure that the Hubble models can be evaluated and fine-tuned on commercially available hardware, facilitating further research. The NVIDIA A100 GPUs, allocated as a part of the NAIRR budget, also set a limit on the model size for efficient training. 

% \item We choose a dataset size of 1TB.

% \item We want to ensure the availability of comparably sized open-source \cite{DBLP:conf/icml/BidermanSABOHKP23,groeneveld-etal-2024-olmo,NEURIPS2024_19e4ea30_datacomp} or open-weight \cite{grattafiori2024llama3herdmodels} models. The open-weight models serve as guidelines and sanity checks for the model quality we can expect for our compute budget.

% \item The total amount of repeated data should comprise a small fraction of the total pre-training data to prevent degradation in model quality \citet{hernandez2022scalinglawsinterpretabilitylearning}.

% \item We want to establish the effect of increasing the size of the pre-training corpus on the memorization of repeated sequences. Unlike prior work, which established the strength of memorization only as a function of the number of duplicates of a sequence \citep{tirumala2022memorization}, we want to study the joint effect of the number of duplicates and the total corpus size. Thus, we want to train the Hubble suite on two different corpus sizes.    \end{enumerate}  

% The Hubble model suite consists of 8 models in a $2 \times 2 \times 2$ factorial design, with the model size $\{1\text{B}, 8\text{B}\} \times$ data condition $\{\text{standard}, \text{perturbed}\} \times$ training set size $\{100\text{B}, 500\text{B}\}$.

% 1B -> 1,179,486,208 params -> 3.73 × 10^21 FLOPs over 500B tokens -> 37.3 / ((1.18)^2 * 1.21)x Chinchilla optimal
% 8B -> 8,265,306,112 params -> 3.1 × 10^22 FLOPs over 500B tokens -> 3.1 / ((8.27/10)^2 * 1.23)x Chinchilla optimal

\section{Domain-agnostic Results} \label{sec:results_domain_agnostic}

We present our domain-agnostic studies on the \textit{spacing} and \textit{placing} of duplicates in LLM training. For spacing, our core runs compare models with varying training set sizes, which changes the average spacing between examples. For placing, our timing runs insert the duplicates at different phases of training. Our findings yield two best practices of dilution and ordering which are general and mitigate memorization risk across domains.

\begin{figure}[h]
    \centering
    \includegraphics[width=\textwidth]{figures/dilution/dilution-hubble_8b.pdf}
    % \vspace{-4ex}
    \caption{\textbf{Memorization of sensitive data can be diluted by training on larger corpora.} We report the base evaluations on a subset of tasks for the core 8B models trained on 100B and 500B tokens. The core runs are described in \S\ref{sec:models} and evaluations are described in \S\ref{sec:evaluations}. For the same duplicate level, memorization is weaker for the model trained on 500B tokens compared to 100B. Figure \ref{fig:model_scale-dclm_500b} compares these trends against the 1B models, and larger models memorize at lower duplications. These experiments represent multiple interventions in one training run, and Figure \ref{fig:interference-full} plots these results for our interference models, which confirm minimal interference across domains.}
    \label{fig:dilution-hubble_8b}
\end{figure}
\clearpage

\textbf{Diluting sensitive data by training on larger corpora reduces memorization risks.} Figure \ref{fig:dilution-hubble_8b} plots the memorization evaluations for the perturbed 8B models trained on either 100B or 500B tokens. Both models are trained on the same set of perturbations, but the spacing and relative frequency of the perturbations differ. When trained on more tokens, the model's memorization on nearly all tasks in all domains increases slower with respect to frequency. This generalizes the result of \citet{bordt2025how}, which showed that scaling the training corpus reduces the effect of test set contamination. These findings suggest a simple best practice to address memorization risks broadly: sensitive data can be \textit{diluted} by training on larger corpora and is complementary to the best practice of deduplication \citep[recommended in][]{pmlr-v162-kandpal22a,lee-etal-2022-deduplicating}.

% cite forgetting papers
\textbf{Ordering sensitive data to appear early in training reduces memorization risks.} We present a selection of results for the timing runs in Figure \ref{fig:injectrange-main} and the full set of results in Figure \ref{fig:injectrange-full}. When perturbations are inserted in only the first quarter of training, the final model does not memorize the data. From Figure \ref{fig:checkpoint-forgetting-curves}, the intermediate checkpoints show that if the model does not receive continued exposures to duplicates, the model can forget the perturbations, which provides a form of privacy \citep{jagielski2023measuring,chang2024large}. When all perturbations are inserted in the last quarter of training, more data is memorized and extractable than the regular perturbed model. This is consistent with \cite{more2025realisticextractionattacksadversarial}, which finds that data at the end of training is more likely to be extractable. This suggests a second best practice to address memorization risks: sensitive data can be \textit{ordered} to appear early in training. 

\textbf{Larger models memorize at lower duplications.} Figure~\ref{fig:model_scale-dclm_500b} compares the memorization strength of both the 1B and 8B parameter models trained on the 500B token corpus. Consistent with prior work \citep{tirumala2022memorization}, the 8B model shows higher memorization across all tasks at the same duplication level, and memorization is measurable with fewer duplicates. Increasing the model size increases memorization risk, so practitioners will need to balance the effects of model scaling with other mitigation strategies such as dilution or ordering.

% \textbf{The concentration of perturbations affects the memorization.} Two sets of results, first comparing concentrations starting at 0 (denser is faster memorization). And then 2 concentrations ending at 50 (sparser is slower decay). And then the final memorization strength for checkpoints ending at 100 (denser is stronger).
% This finding supports the \textit{learnability threshold} proposed by \citet{}: models will not retain a fact if it is not re-encountered in a certain number of training steps.

% \textbf{Pre-training on more tokens consistently dilutes memorization.} Figure \ref{}

% \subsection{Recommendations}
% this will go to the appendix
% Memorization risks of LLMs will likely be managed through a body of soft law, which are non-binding rules, guidelines, and best practices. While such practices lack the backing of legislation, they often derive legitimacy from judicial or executive authority \citep[e.g. when a judge determines the legality of a new technology based on whether its deployment adheres to best practices;][]{hagemann2018soft}, and can gain industry adoption. Besides the outputs of LLMs and safeguards, \cite{wei_2025_interrogating} point out that the structural design choices of LLMs may be scrutinized by judges or regulators when assessing copyright risk, and major commercial models, such as Gemini, already emphasize measurements of memorization in their public documentation \citep{team2024gemini, apertus2024}.
% A well-established best practice is deduplication, which reduces the likelihood of memorization by removing repeated examples from the training corpus \citep{lee-etal-2022-deduplicating, hans2024be}. Yet deduplication has limits: not all sensitive information can be identified or eliminated, and other factors such as the distribution of examples across training time can also affect memorization \citep{tirumala2022memorization, pmlr-v162-kandpal22a, DBLP:conf/iclr/CarliniIJLTZ23, chang2024large, more2025realisticextractionattacksadversarial}. Our core experiments establish a second, complementary best practice: dilution. By training on larger corpora, models become more tolerant to insertions of sensitive information, reducing their relative prominence. This effect has also been observed from a forgetting perspective \citep{bordt2025how}. Together, deduplication and dilution expand the technical toolkit available for mitigating memorization risk.

% \subsection{Temporal effects}
% - inject range memorization
% - checkpoints (for inject range)
% - different visualizations?

% For the original perturbed model, the perturbation data was injected uniformly randomly throughout training. We train six additional models to measure the effect of the number of training steps before and after exposure to the perturbation data on the strength of memorization. These models have 1B parameter models and are trained on 100B tokens, where the perturbation data is injected in a specific range of batches: (0\%, 50\%), (50\%, 100\%), (0\%, 25\%), (25\%, 50\%), (50\%, 75\%), (75\%, 100\%). The percentage refers to the progress in the overall training; e.g. (50\%, 100\%) means batches in the second half of training.

\begin{figure}[t]
    \centering
    \includegraphics[width=\textwidth]{figures/injectrange/injectrange-main.pdf}
    \vspace{-4ex}
    \caption{\textbf{Sensitive data can be forgotten without continued exposure.} We report the performance of the Timing runs (1B models trained on 100B tokens, described in \S\ref{sec:models}) where perturbations are inserted in different phases of pretraining (tuples denote the range of pretraining where texts are inserted). For reference, the standard and perturbed 1B parameter models are also plotted.}
    \label{fig:injectrange-main}
\end{figure}

% \todo{Nitya}{Add discussion on intermediate checkpoint analysis}

% \subsection{Interference}

%\begin{figure}
%    \centering
%    \includegraphics[width=\textwidth]{figures/interference/interference-main.pdf}
%    \caption{\textbf{The perturbed model matches the behavior of domain-specific models on the respective set of evaluations.} The perturbed model matches the \texttt{copyright\_only} model in memorizing the unpopular Gutenberg passages, \texttt{privacy\_only} model in generating memorized PII from biographies, and \texttt{testset\_only} model in memorizing MMLU questions. This figure contains a subset of the evaluations (refer to Figure\ref{fig:interference-full} for all evaluations).}
%    \label{fig:interference-main}
%\end{figure}

% As discussed in \S~\ref{sec:models},
\textbf{Perturbations from different domains minimally interfere with each other.} Our perturbed models are the product of many interventions in a single training run. If the perturbations interfere with each other (e.g., a highly duplicated example in a test set affects the memorization of a paraphrase), that would undermine the validity of our analyses. Although exhaustively characterizing such interference \citep[as in][]{ilyas_datamodels_2022} would be impractical, we perform a check by training three 1B models each containing perturbations from only a single risk domain. As shown in Figure~\ref{fig:interference-full}), the behavior of the core perturbed model matches every single-domain model on the corresponding domain. These suggest that our aggregate, domain-level findings have minimal interference.

\section{Domain-specific Results} \label{sec:results_domain_specific}

The perturbation data in \textsc{Hubble} is designed to enable a broad range of experimentation. In this section, we present tailored analyses for the specific data types in each domain.

\subsection{Copyright}

\textbf{Whether an LLM is considered to memorize depends on the metric.} In Figure \ref{fig:copyright-passages}, we additionally evaluate $k$-eidetic memorization \cite[introduced in][]{DBLP:conf/iclr/CarliniIJLTZ23} and the ROUGE-L metric on the passages in the copyright domain (details in Appendix \ref{app:copyright_specific_results}).  While loss can show statistically significant differences in memorization at lower duplicate counts, the $k$-eidetic metric does not. This can be seen for Wikipedia passages at 4 duplicates, where loss shows significant differences for the 8B, 100B model, but $k$-eidetic memorization does not, and differences only start to show at 16 duplicates.  For copyright debates, this means that the choice of metric affects the interpretation of a memorization analysis, and numerical measures are unlikely to be useful on their own.

\textbf{Popular and unpopular books are memorized similarly by the 1B model, with only minor differences for the 8B model.} Based on the data density hypothesis \citep{kirchenbauer2024lmd}, we expected popular books from Gutenberg would be memorized better than unpopular books, as popular books are more likely to be discussed in the pretraining corpus. In Figure \ref{fig:model_scale-dclm_500b}, At the 1B parameter scale, there is no noticeable difference, and at the 8B parameter scale, there is only a slight increase in the generative extraction of passages from popular books compared to unpopular books. The 8B parameter models trained on 100B and 500B tokens both assign a slightly higher likelihood to passages from the popular books. While we find little difference for popular books using basic evaluations, more sensitive methods may reveal subtler forms of memorization.

% Define color boxes for slots and attributes
\definecolor{slotgray}{gray}{0.9}
\definecolor{attrpurple}{HTML}{E6E6FA} % light lavender

\newcommand{\slotlabel}[1]{\colorbox{gray!20}{\texttt{#1}}}
\newcommand{\slot}{\colorbox{slotgray}{\texttt{<candidate>}}}
\newcommand{\attr}[1]{\colorbox{attrpurple}{\texttt{#1}}}

\begin{table}[t]
\centering
\small
\caption{\textbf{Attack Definitions for YAGO.} PII attacks are listed below in increasing order of strength (fewer additional PII known to the attacker). Each attack corresponds to a different prompt, and we illustrate the attacker's query to infer the target's university using a sample biography from YAGO. The full prefix–full suffix attack is only compatible with infill attacks (loss-based choice) since generations cannot be conditioned on the suffix. Attack success rates are presented in Figure \ref{fig:privacy-yago}, and a breakdown of success rate by PII type is given in \ref{fig:privacy-yago-meta}.}
\label{tab:pii_attack_definition}
\begin{tabularx}{\textwidth}{c c >{\raggedright\arraybackslash}X}
\toprule
\multicolumn{3}{>{\raggedright\arraybackslash\hsize=\textwidth}X}{
\textbf{Sample YAGO Biography (templated slots shaded purple)}%
}\\
\midrule
\multicolumn{3}{>{\raggedright\arraybackslash\hsize=\textwidth}X}{
\texttt{%
\attr{Dora Sloan} is from \attr{the United States}. %
\attr{Dora} was born in \attr{Phoenix, Arizona}. %
\attr{Dora} is an alumni of \attr{St. John's College}. %
\attr{Dora} was born on \attr{May 15, 1968}. %
\attr{Dora} receives email at \attr{dora@gmail.com}. %
\attr{Dora} is a \attr{competitive diver}. %
\attr{Dora} has the unique identifier \attr{4dc0969af29a4324bf5746c50f7209a2}.%
}
}\\
\toprule
\textbf{Prompt Format} & \textbf{Compatible Attack} & \textbf{Example Query (suffix underlined)} \\
\midrule
full prefix–full suffix & Infill &
\texttt{Dora Sloan is from the United States. Dora was born in Phoenix, Arizona. Dora is an alumni of }\underline{\slot\texttt{. Dora was born on}} ~~~~ \underline{\texttt{May 15, 1968. Dora receives email at}} ~~~~ \underline{\texttt{dora@gmail.com. Dora is a competitive diver.}} \underline{\texttt{Dora has the unique identifier}} \underline{\texttt{4dc0969af29a4324bf5746c50f7209a2.}}\\
\midrule
full prefix & Infill, Gen &
\texttt{Dora Sloan is from the United States. Dora was born in Phoenix, Arizona. Dora is an alumni of }\underline{\slot\texttt{.}}\\
\midrule
intro prefix & Infill, Gen &
\texttt{Dora Sloan is from the United States. Dora is an alumni of }\underline{\slot\texttt{.}}\\
\midrule
name only & Infill, Gen &
\texttt{Dora Sloan is an alumni of }\underline{\slot\texttt{.}}\\
\bottomrule
\end{tabularx}
\end{table}

\subsection{Privacy}

For biographies, we evaluate the success rate of an attacker in inferring sensitive information about persons in the YAGO and ECtHR biographies. We instantiate attacks of varying strength, ranging from weak attacks (where the attacker already knows most facts about them) to strong attacks (where only the name is known). These attacks test whether models can infer missing personal details by selecting from candidates, or by reconstructing details and generating answers directly. Different attacks correspond to different prompts and Table \ref{tab:pii_attack_definition} visualizes them for YAGO. YAGO results are reported in Figure \ref{fig:privacy-yago}, and the breakdown by PII type is given in Figure \ref{fig:privacy-yago-meta}. Further details and results on ECtHR are in Appendix~\ref{app:direct_pii_leakage}. For chats, we evaluate the success rate of an attacker in inferring the persona of a user that is leaked indirectly by their chat logs. The evaluation formats are described in Table \ref{tab:chat_attack_definition} and results on Personachat are presented in Figure \ref{fig:privacy-personachat}.

% For biographies, we not only evaluate the memorization of the biographies but also whether sensitive information can be inferred about the fictional people. For direct memorization, we report the loss assigned by the model to the inserted biography. To evaluate the ease of PII reconstruction, we instantiate attacks with varying strength. \textbf{Weak attacks} assume that the attacker already knows PII about the person of interest and is seeking a few missing facts. \textbf{Strong attacks} assume that the attacker knows less sensitive information about the person of interest, with our strongest attacks assuming that the attacker only knows the name. We instantiate \textbf{loss-based choice attacks} where the attacker has narrowed down the possible values of the missing PII. We frame the attack as MCQ problems and check which candidate answer has the highest likelihood when plugged into the blank. When the attacker has no way to deduce the set of candidate answers, they have to use \textbf{generative attacks} where the model is prompted to fill in the blank. We evaluate generative attack with either \textit{Word Recall}, which scores if the answer entity occurs anywhere in the generated response, or \textit{Prefix Match}, which scores whether the model generation starts with the answer entity. Table \ref{tab:pii_attack_definition} lists the attacks that we instantiate. The synthetic YAGO biographies allow us to instantiate each of the attacks listed in the table. We can only instantiate the \textit{full prefix, generative} attack for ECtHR since the entity types are not clearly defined (e.g., dates can refer to birth dates or event dates) and not all entity types are always present in the biography. Figures \ref{fig:privacy-ecthr} and \ref{fig:privacy-yago} report attack success rates on ECtHR and YAGO perturbation sets, respectively. Figure \ref{fig:privacy-yago-meta} provides a breakdown by PII type for reconstruction attacks on YAGO biographies (rows are arranged in the order that the PII type occurs in the biography).

% \textbf{Privacy.} In Appendix \ref{app:privacy_specific_results}, we study the reconstruction of PII in YAGO biographies. We find that the more pieces of auxiliary information the attacker has access to, the higher the success rate of reconstruction for a given PII in the biography. For the paraphrased models, which were trained on paraphrased biographies, PII reconstruction attacks remains successful. This means the paraphrased model has not just memorized a fixed string, but generalizes to unseen queries for the PII and this knowledge is retrievable \citep[similar to the retrievability observed in][]{allen-zhu_physics_2024}. Personachat also shows the model's ability to retrieve memorized information, and models can infer a user's persona based on the memorized chat logs (although the accuracy is low).

\textbf{The more auxiliary information the attacker has access to, the higher the success rate.} For both ECtHR (Fig \ref{fig:privacy-ecthr}) and YAGO (Fig \ref{fig:privacy-yago}), the attacks with the most auxiliary information are the most effective in inferring PIIs with high accuracy. Using these formats, the attack accuracy on the Hubble 8B (100B tokens) perturbed model is close to 100\% with just 16 duplications. When provided less auxiliary information (e.g. name only) the accuracy of inference decreases significantly.

\textbf{Memorization research needs to account for variation across PII data types.} By comparing attack success across PII types \citep{lukas2023analyzing}, we find that certain attributes such as occupation, email, and UUID are memorized differently from others. Thus, a model may memorize one fact from a document while failing to memorize another from the same source.

\paragraph{Inference of indirect information is difficult but possible.} The details of our analysis on PersonaChat is in Appendix \ref{app:indirect_pii_extraction}. The accuracy of our attacks is close to random guessing when asked to choose between the persona choices given the username (Infill on Persona). While the Hubble models memorize the chat logs for the user, they do not directly assign higher likelihoods to the correct underlying persona. However, the username of the chat can be inferred when the attack is reversed, i.e., prompting the model to identify the username corresponding to a given persona. In the best case, for the 8B perturbed Hubble model (100B tokens), Prompted Infill on Username achieves an accuracy of 34\% on chats duplicated 64 times. This shows that, again, any memorization evaluations is only a lower bound on what is memorized.

\textbf{PII can still be inferred from paraphrased biographies.} In Appendix \ref{app:paraphrase_results}, we trained a `paraphrase' model with different paraphrases rather than exact duplicates of the biographies. The high accuracy of PII recontruction and inference indicates the paraphrase model has not just memorized a fixed string; instead, it generalizes to unseen queries for the PII, and this knowledge remains retrievable \citep[similar to the retrievability observed in][]{allen-zhu_physics_2024}.
The accuracy of strong name-only attacks is higher on the 8B-parameter paraphrase model than on the original perturbed model at high duplication levels, indicating that models trained on paraphrases develop stronger semantic memory than the verbatim memory formed from training on exact duplicates.
Personachat also shows the model's ability to retrieve memorized information in new contexts, and models can infer a user's persona based on the memorized chat logs (although the accuracy is low).

\subsection{Test Set Contamination}

% pretraining on test set examples, doesn't help you generalize on test sets. 
\textbf{Models begin to memorize test set examples with as few as one duplicate, but generalization to unseen examples is unpredictable.} From Figure \ref{fig:testset-set1}, we see that the Hubble perturbed models trained on 100B tokens show an increase in accuracy on PopQA, HellaSwag, and PIQA with just 1 instance of contamination.
However, memorizing test set examples does not translate into generalization on that task: perturbed models show no improvement over standard models when trained on contaminated tasks (reflected in model performance on 0 duplicates), aside from small improvements on PopQA and under certain settings of HellaSwag. In fact, model performance on unseen examples degrades for WinoGrande and a few settings of HellaSwag.
For WinoGrande (see Figure \ref{fig:testset-wg}), we find that perturbed models achieve worse accuracy on minimal pairs of contaminated examples than unseen examples. Further analysis and details are presented in Appendix \ref{app:testset_specific_results}.
Likewise, the paraphrased model fails to answer MMLU questions which were contaminated with paraphrases of that question. We hypothesize that pretraining on a handful of contaminated test examples is not enough to generalize on the task, leading only to memorization.

\paragraph{For WinoGrande, models do not generalize across formats and have worse accuracy on contaminated examples in a new format than on unseen examples.} We inserted two variants of WinoGrande, one in the standard infill (cloze) format, and another in the MCQ format, where options are presented with the question and the model has to generate the correct option. In Figure \ref{fig:testset-wg}, we report the model accuracy when the test time format does not match the inserted format. For examples inserted with the MCQ format, when tested on the infill format, the perturbed model accuracy even decreases with increased duplication.

% Our perturbed models represent multiple interventions from multiple experiments. A validity concern is that many perturbations in the same training run could cause interference effects (e.g. if two perturbations were close duplicates of each other), making it unclear whether results reflect the intended perturbation or unexpected interactions. In general, checking for all interferences is difficult and requires advanced methods to determine example influence \citep{ilyas_datamodels_2022}.
% We train three 1B parameter models on three different corpora each containing perturbations from only one domain. We highlight a subset of our findings in Figure \ref{fig:interference-main} (see Figure \ref{fig:interference-full} for the remaining evaluations). Our evaluations establish that the core perturbed model (trained on perturbations from all domains) matches the behavior of a model trained on perturbations from just the copyright domain (\texttt{copyright\_only}) on the copyright evaluations. Corresponding findings hold for the \texttt{privacy\_only} and \texttt{testset\_only} ablations. Thus, our core perturbed models remain valid for studying each domain individually, behaving as if they were trained on that domain alone despite joint training.

% Test set contamination.  However, this memorization does not translate into broader generalization: performance on unseen examples shows only marginal gains (e.g., PopQA, HellaSwag), and paraphrased training examples do not help the model answer the originals (Figure \ref{fig:paraphrase-full}). In other words, pretraining on contaminated test set examples leads to memorization of those items rather than meaningful generalization.
% We also identify a case on Winogrande where exposure to some examples makes the model worse at generalizing to unseen examples, as the dataset includes minimally contrasting examples with different answers (see Figure \ref{fig:testset-wg}).

% We reiterate that our evaluations are designed to serve as lower bounds, i.e., there may exist other formulations of prompts and tasks that elicit more copyright violations, leak sensitive information, and achieve higher leaderboard scores with the same set of models.

% \section{Experimental Results}

% main paper results are general
% look in the appendix for results that are more specific to certain areas

% in the first part, we first analyze general memorization results.
% in the second part, we look at specific perturbation while yields interesting observations in specific domains
% This section is organized in two parts. First, we  present our domain-agnostic studies on the \textit{spacing} and \textit{placing} of duplicates in LLM training. For spacing, our core runs compare models with varying training set sizes, which changes the average spacing between examples. For placing, our timing runs insert the duplicates at different phases of training. Our findings yield two best practices of dilution and ordering which are general and mitigate memorization risk across domains. In the second part, we present our domain-specific studies, where we analyze specific perturbations in \textsc{Hubble} to yield a rich set of observations for the domains of copyright, privacy, and test set contamination.

% This section first presents our findings on the \textit{spacing} and \textit{placing} of duplicates in LLM training. Since our perturbations are randomly inserted throughout the training process, increasing the size of the training data increases the average spacing between duplicates. For spacing, our core runs show that increasing the total amount of training data, which increases the average spacing between examples, reduces memorization. For placing, our timing runs insert the duplicates at different phases of training, and our findings show that memorization is strongest when duplicates appear closer to the end of training. This leads to two best practices which are domain-agnostic: dilution and ordering.
% An intuition from \cite{chang2024large} is that the memorization of an example is peaky, spiking when a duplicate is encountered and decaying afterward.
% Besides our results on general memorization risk, we include a variety of domain-specific experiments in Appendix \ref{app:domain_results}. We simulate a number 
\section{Use Cases of Hubble}
\label{sec:use_cases}

\begin{table}[tb]
\centering
\small
\caption{\textbf{ROC AUC scores of baseline MIAs for our largest perturbed model (8B, 500B tokens).} \textit{Dup} indicates the duplication level of members. \textit{Dup $\neq$ 0} treats all inserted perturbations as members. Non-members are always drawn from perturbations inserted 0 times. As duplication increases, memorization becomes stronger, and it becomes easier for membership inference attacks (MIA) to distinguish between members and non-members. See Appendix \ref{app:mia_results} for more \textsc{HubbleMIA} settings.}
\label{tab:hubble_mia_8b_perturbed}
\begin{tabular}{clcccccc}
\toprule
\multirow{2}{*}{\textbf{Evaluation}} & \multicolumn{1}{c}{\multirow{2}{*}{\textbf{MIA}}} & \multicolumn{6}{c}{\textbf{\textsc{Hubble} 8B (500B tokens) Perturbed}}                         \\ \cmidrule(lr){3-8} 
                                     & \multicolumn{1}{c}{}                                        & Dup $\neq$ 0 & Dup = 1 & Dup = 4 & Dup = 16 & Dup = 64 & Dup = 256 \\ \midrule
\multirow{4}{*}{\shortstack{Gutenberg\\Unpopular}} & Loss                                                        &   0.629       &   0.539      &  0.556       &   0.732       &   \textbf{0.996}       &  \textbf{1.0}           \\  
                                     & MinK\%                                                      &   0.629       &   0.539      &  0.556       &   0.732       &   \textbf{0.996}       &  \textbf{1.0}           \\  
                                     & MinK\%++                                                    &   \textbf{0.666}       &   \textbf{0.545}      &  \textbf{0.62}       &   \textbf{0.813}       &   0.987       &  0.949           \\  
                                     & ZLib                                                        &   0.622       &   0.53      &  0.551       &   0.722       &   \textbf{0.996}       &  \textbf{1.0}           \\
\midrule
\multirow{4}{*}{\shortstack{Yago\\Biographies}}    & Loss                                                        &   0.692       &   0.538      &  0.652       &   \textbf{0.897}       &   \textbf{1.0}       &  \textbf{1.0}           \\  
                                    & MinK\%                                                      &   0.692       &   0.537      &  0.651       &   0.896       &   \textbf{1.0}       &  \textbf{1.0}           \\  
                                    & MinK\%++                                                    &   \textbf{0.714}       &   \textbf{0.571}      &  \textbf{0.686}       &   0.892       &   0.995       &  0.983           \\  
                                    & ZLib                                                        &   0.676       &   0.524      &  0.633       &   0.872       &   \textbf{1.0}       &  \textbf{1.0}           \\
\midrule
\multirow{4}{*}{MMLU}                & Loss                                                        &   0.673       &   0.529      &  0.628       &   0.857       &   \textbf{1.0}       &  \textbf{1.0}           \\  
                                    & MinK\%                                                      &   0.672       &   0.529      &  0.626       &   0.854       &   \textbf{1.0}       &  \textbf{1.0}           \\  
                                    & MinK\%++                                                    &   \textbf{0.743}       &   \textbf{0.58}      &  \textbf{0.731}       &   \textbf{0.943}       &   0.994       &  0.986           \\  
                                    & ZLib                                                        &   0.644       &   0.523      &  0.593       &   0.775       &   0.993       &  0.999           \\
\bottomrule
\end{tabular}
\end{table}


The randomized perturbations in \textsc{Hubble} are designed to enable a broad range of research on LLM memorization. To demonstrate this, we establish new benchmarks for both membership inference attacks (MIAs) and unlearning. Membership inference is the task of inferring which data was part of a  model's training set and MIAs are used to audit privacy risks of trained models  \citep{shokri2017membershipinferenceattacksmachine}. Machine unlearning erases  harmful knowledge or behaviors from models while preserving other capabilities, without requiring full retraining \citep{bourtoule_2021_machine, llm_unlearning}.

\subsection{Hubble as an MIA Benchmark}
\label{sec:hubble_as_mia}

\textbf{Current MIA benchmarks for LLMs.}~ \citet{minkMIA} introduces \textsc{WikiMIA}, a membership inference benchmark for LLM pretraining data. \textsc{WikiMIA} labels Wikipedia articles before and after a model’s knowledge cutoff as members and non-members, respectively. However, subsequent analyses found that spurious features (such as temporal cues) allow non-members articles to be trivially distinguished from members, undermining the benchmark's validity \citep{duan2024membershipinferenceattackswork,meeus_sok_2025, naseh2025syntheticdatamisleadevaluations}. At the same time, this line of work shows that detecting pretraining data is generally difficult. When using the randomized train and test sets of Pythia, most membership inference methods achieve only marginal performance.

\textbf{The \textsc{HubbleMIA} benchmark.}~\textsc{Hubble} provides a sound benchmark for evaluating membership inference on several data types, including book passages, PII, and standard evaluation test sets. Since each perturbation is randomly duplicated zero or more times, there are no spurious features that inadvertently leak membership information. Perturbations in \textsc{Hubble} are also decontaminated and inserted at different frequencies, allowing comparisons of membership inference effectiveness on low- versus highly-duplicated examples. % Evaluating MIAs on the standard models (which were not trained on the perturbation data at all) serves as a sanity check, and when AUC is different from 50\%, MIA methods incorrectly detect non-members as members (false positive) at a rate higher than random chance.

\textbf{Experimental setup.} We instantiate 12 membership inference settings as a representative subset of all possible MIA benchmarks enabled by the Hubble Suite: 4 Hubble model variants (two perturbed models and two standard models) on 3 perturbation datasets each (Gutenberg Unpopular, YAGO Biographies, and MMLU). MIAs are evaluated with perturbations duplicated zero times as non-members, and perturbations duplicated more than once as members. For this evaluation, we employ off-the-shelf implementations from OpenUnlearning \citep{dorna2025openunlearningacceleratingllmunlearning}, specifically testing Loss-based \citep{LossMIA}, MinK\% \citep{minkMIA}, MinK\%++ \citep{minkplusplusMIA}, and Zlib-based attacks \citep{carlini_extracting_2021}.

\textbf{Results.} Table~\ref{tab:hubble_mia_8b_perturbed} reports MIA performance of Gutenberg Unpopular for our most capable model (8B, 500B tokens). MIA performance on all datasets and models are presented in Appendix~\ref{app:mia_results}. Across all benchmarks, membership inference performance consistently improves as the duplicate count increases, and attacks are strongest when distinguishing non-members from members duplicated 256 times.  However, distinguishing members duplicated only once produces near-random results, which confirm observations in \citet{duan2024membershipinferenceattackswork} that MIAs perform well only on members that are highly duplicated. Generally, our results show MinK\%++ to be the most effective attack. Surprisingly, MinK\%++ does not achieve 100\% AUC on the highly duplicated samples, unlike simpler approaches such as Loss and MinK\%.

\subsection{\textsc{Hubble} as an Unlearning Benchmark}

\textbf{Current LLM unlearning benchmarks.} Existing benchmarks target different aspects of machine unlearning. TOFU \citep{tofu} focuses on the unlearning of private data through  synthetic biographies. However, TOFU operates in a fine-tuning setting, where models are fine-tuned on the data to be forgotten. MUSE \citep{muse} focuses on unlearning copyrighted text, such as Harry Potter fan-fiction and news articles, but is also limited to unlearning in fine-tuning rather than pretraining. Finally, WMDP \citep{wmdp} focuses on removing harmful capabilities.

\textbf{The \textsc{HubbleUnlearning} Benchmark.} \textsc{Hubble} provides a benchmark for evaluating unlearning methods on data in pretraining spanning diverse domains. Because the forget and retain sets are drawn from the same distribution, methods must remove the forget set with high specificity while preserving performance on neighboring examples. The standard models in \textsc{Hubble} were not trained on any perturbations and are also useful as an additional point of reference. Finally, unlearning is tested on data where the duplicate count is known and consistent, removing a confounder in the evaluation of unlearning methods \citep{krishnan2025dataunlearnedequally}.

\begin{wrapfigure}{r}{0.5\columnwidth}
    \centering
    \includegraphics[width=\linewidth]{figures/Hubble_unlearn/unlearn_main.pdf}
    \caption{\textbf{Unlearning performance on with \textsc{Hubble} 8B in copyright and privacy.} Three key reference points are included in each subplot: the perturbed model (
\includegraphics[height=1.5ex]{figures/Hubble_unlearn/baseline.png}), representing performance before unlearning; the standard model (
\includegraphics[height=1.5ex]{figures/Hubble_unlearn/standard.png}), which is trained without perturbations; and the desired model (
\includegraphics[height=1.5ex]{figures/Hubble_unlearn/target.png}), which achieves standard model's performance on the forget set while retaining the perturbed model's performance elsewhere. Improvement is indicated by the arrow (
\includegraphics[height=1ex]{figures/Hubble_unlearn/improve.png}). See Appendix \ref{sec:unlearn_full_results} for the full results.} %\williex{can make legend box wider/text bigger. Also, I'd label each x-axis to be more clear.}}
    \label{fig:unlearn_main}
    \vspace{-5mm}
\end{wrapfigure}

\textbf{Setup.} We benchmark three representative unlearning methods on our largest perturbed model (8B, 500B tokens): Representation Misdirection for Unlearning \citep[RMU;][]{wmdp}, Representation Rerouting \citep[RR;][]{rr}, and Saturation-Importance \citep[SatImp;][]{satimp}. Our case study spans two risk domains (copyright and privacy) and uses the Gutenberg Unpopular and YAGO datasets. Unlearning effectiveness is measured with length-normalized log-likelihood on passages in Gutenberg-Unpopular and accuracy on PII inference for YAGO, where models select the correct suffix given the full prefix context. 

Each dataset is split into three subsets: (1) \textbf{Unseen}, consisting of the held-out perturbations (i.e., duplicated 0 times); (2) \textbf{Unlearn}, consisting of half of the 256 duplicate perturbations as unlearning targets; and (3) \textbf{Keep}, consisting of the other half of the 256 duplicate perturbations, which are near-neighbors to the unlearn set and are should be kept. Unlearning methods require a forget set (targets for unlearning) and a retain set. Following prior work, we use \textbf{Unlearn} as the forget set, and WikiText \citep{wikitext} as the retain set to approximate general knowledge \citep{wmdp, erase_llm}. For each unlearning method, we run a grid search over method hyperparameters, and further details are provided in Appendix \ref{sec:unlearn_grid_search}.

\textbf{Results.} As shown in Figure~\ref{fig:unlearn_main}, no unlearning method reaches the desired target and matches the performance of the standard model on the Unlearn set while retaining the other sets. Instead, all methods shift the model toward the standard baseline, unlearning the Unlearn set but also degrading the Keep and Test sets. Degradation on the test set is similar to utility degradation observed in \citet{muse}. Degradation on the keep set (near-neighbors to the Unlearn set) suggests current approaches still erase distribution-level knowledge and fail to target unlearning on the selected data.  Generally, SatImp performs best and produces more unlearned checkpoints closer to the desired target, but there is still room for improvement in the method's precision. We provide additional unlearning results in Appendix \ref{sec:unlearn_full_results}, where we use the in-distribution \textbf{Keep} set as the retain set instead of WikiText; the general patterns remain consistent, with RMU and RR performing worse. 

% Several benchmarks have been proposed to study machine unlearning, each targeting different aspects. TOFU \citep{tofu} creates synthetic author biographies and finetunes models on them, providing a controlled benchmark for unlearning. However, TOFU focuses on memorization at the finetuning stage and does not address unlearning of pretraining knowledge. MUSE \citep{muse} evaluates unlearning on narrow real-world domains such as Harry Potter books and news articles. Another benchmark is WMDP \citep{wmdp} emphasizing removal of harmful capabilities rather than memorized training data. 

% \textbf{Current LLM unlearning benchmarks.} TOFU \citep{tofu} synthesizes fictitious biographies and finetune a model to serve as an unlearnings benchmark. It is closest to our work, but TOFU inserts data during finetuning, and does not test unlearning of pretraining knowledge. MUSE \citep{muse} is another unlearning benchmark which assumes some books are inserted targets narrow domains like Harry Potter books and news articles. Other benchmarks address different but limited aspects: WMDP \citep{wmdp} emphasizes the removal of harmful capabilities rather than memorized training data,  With Hubble, we enable the study of unlearning across diverse forms of pre-training data, capturing a broader range of memorization risks. 



% \textbf{Current LLM unlearning benchmarks} MUSE \citep{muse}, WMDP \citep{wmdp}, and TOFU \citep{tofu} are popular unlearning benchmarks, they highlight key challenges but have notable gaps. MUSE focuses on unlearning harry potter books and news articles, which restricts its generality. WMDP measures hazardous knowledge but cannot link forgetting to specific pre-training data. TOFU synthesizes biographies of fictitious authors, and requires finetuning on the dataset before unlearning, diverging from realistic pre-training scenario. 


% \textbf{The \textsc{HubbleUnlearning} benchmark} We use \textsc{Hubble} models as a testbed for evaluating targeted unlearning across two risk domains of copyright and privacy. In contrast, \textsc{Hubble} addresses these shortcomings by (1) spanning diverse domains, (2) enabling evaluation on individual memorized data, and (3) introducing memorization directly during pre-training, providing a more faithful testbed for unlearning.  HUBBLE offers two advantages as a benchmark: it enables direct comparison with standard models trained without perturbations, and since \textsc{HUBBLE} provides perturbation and clean same-distribution samples. This allows precise evaluation of whether unlearning can target only intended data or also neighboring examples. 

% \textbf{Setup} We unlearn the \textsc{Hubble} 8B perturbed model trained on 500B tokens, and choose the \textsc{Hubble} 8B standard model to compare with. We adpot three representative methods: \textit{Representation Misdirection for Unlearning (RMU)} \citep{wmdp}, \textit{Representation Rerouting (RR)} \citep{rr}, and \textit{Saturation-Importance (SatImp)} \citep{satimp}.  

% Unlearning methods operate on two datasets: a \textit{forget set}, containing the target data to remove, and a \textit{retain set}, approximating general knowledge to preserve. For each unlearning domain, we use the perturbation set with 256 duplicates as forget set, and WikiText \citep{wikitext} as retain set following prior work \citep{wmdp, erase_llm}.

% We evaluate unlearning on two perturbation datasets on each risk domain: Gutenberg-Unpopular (copyright), and YAGO (privacy). Each dataset is split into three subsets: (1) \textbf{Unseen}, consisting of held-out data never exposed to the model; (2) \textbf{Unlearn}, comprising a randomly selected half of the perturbation set as the target for unlearning; and (3) \textbf{Keep}, containing the remaining half of the perturbation samples.  For each unlearning method, we perform a grid search over method hyperparameters. 


% For every unlearned checkpoint, we measure forgetting on the target domain. We use MCQ accuracy for MMLU and YAGO, and normalized log-likelihood for Gutenberg-Unpopular.

% To assess general capability preservation, we use tinyMMLU, tinyWinogrande, and tinyHellaswag from TinyBenchmarks \citep{tinybenchmarks}, and discard checkpoints with average performance degradation exceeding 10\%. Further grid search details are provided in the Appendix \ref{sec:unlearn_grid_search}. In addition, we experiment with using perturb-keep-256x as the retain set; the corresponding results are reported in Appendix \ref{sec:unlearn_full_results}. 

% \textbf{Results} We evaluate whether existing unlearning methods can selectively forget the targeted Unlearn set while preserving performance on both the Unseen and Keep sets. As shown in Figure~\ref{fig:unlearn_main}, none of the methods achieve the desired target point, defined as matching standard model's performance on the forget set while retaining the perturb model's performance elsewhere, and all methods instead shift the model toward the standard model, reducing performance on the forget set and also influence non-targeted in-distribution samples, including both samples the model memorizes in the Perturb-Keep set and samples the model encountered for the first time in the Test set. 
% This pattern is consistent across all three domains: after unlearning, the performance of three subsets collapse to similar level, indicating that tested unlearning methods erase distribution-level knowledge rather than selectively removing targeted data. 

%\textsc{Hubble} is particularly valuable in this context, as its precise control over data allows MIA developers and evaluators to rigorously test whether an attack truly succeeds on rarely seen samples or whether its apparent performance is dominated by frequently repeated datapoints. 
% We would also like to note that the distinction between “seen” and “unseen” datapoints is meaningful only for perturbed models; for standard models, which have not encountered either set, MIA attacks yield near-random AUC-ROC scores (Table~\ref{tab:hubble_mia_8b_standard},\ref{tab:hubble_mia_1b_standard}). These standard models thus provide a baseline: if an attack shows differences even on them, its performance should bring up questions. 

% % === Benchmark 1: Gutenberg Popular ===
% \begin{figure}[htbp]
%   \centering
%   % Row 1
%   \begin{subfigure}[b]{0.48\textwidth}
%     \centering
%     \includegraphics[width=\linewidth]{figures/Hubble_MIAs/eval_Gutenberg Popular_model_hubble-1b-500b_toks-standard-hf.pdf}
%     \caption{Hubble-1B Standard}
%   \end{subfigure}
%   \hfill
%   \begin{subfigure}[b]{0.48\textwidth}
%     \centering
%     \includegraphics[width=\linewidth]{figures/Hubble_MIAs/eval_Gutenberg Popular_model_hubble-8b-500b_toks-standard-hf.pdf}
%     \caption{Hubble-8B Standard}
%   \end{subfigure}

%   \vspace{2mm}

%   % Row 2
%   \begin{subfigure}[b]{0.48\textwidth}
%     \centering
%     \includegraphics[width=\linewidth]{figures/Hubble_MIAs/eval_Gutenberg Popular_model_hubble-1b-500b_toks-perturbed-hf.pdf}
%     \caption{Hubble-1B Perturbed}
%   \end{subfigure}
%   \hfill
%   \begin{subfigure}[b]{0.48\textwidth}
%     \centering
%     \includegraphics[width=\linewidth]{figures/Hubble_MIAs/eval_Gutenberg Popular_model_hubble-8b-500b_toks-perturbed-hf.pdf}
%     \caption{Hubble-8B Perturbed}
%   \end{subfigure}

%   \caption{Membership inference attack performance on Gutenberg Popular.}
%   \label{fig:mia-gutenberg-popular}
% \end{figure}

% % === Benchmark 2: Gutenberg Unpopular ===
% \begin{figure}[htbp]
%   \centering
%   \begin{subfigure}[b]{0.48\textwidth}
%     \centering
%     \includegraphics[width=\linewidth]{figures/Hubble_MIAs/eval_Gutenberg Unpopular_model_hubble-1b-500b_toks-standard-hf.pdf}
%     \caption{Hubble-1B Standard}
%   \end{subfigure}
%   \hfill
%   \begin{subfigure}[b]{0.48\textwidth}
%     \centering
%     \includegraphics[width=\linewidth]{figures/Hubble_MIAs/eval_Gutenberg Unpopular_model_hubble-8b-500b_toks-standard-hf.pdf}
%     \caption{Hubble-8B Standard}
%   \end{subfigure}

%   \vspace{2mm}

%   \begin{subfigure}[b]{0.48\textwidth}
%     \centering
%     \includegraphics[width=\linewidth]{figures/Hubble_MIAs/eval_Gutenberg Unpopular_model_hubble-1b-500b_toks-perturbed-hf.pdf}
%     \caption{Hubble-1B Perturbed}
%   \end{subfigure}
%   \hfill
%   \begin{subfigure}[b]{0.48\textwidth}
%     \centering
%     \includegraphics[width=\linewidth]{figures/Hubble_MIAs/eval_Gutenberg Unpopular_model_hubble-8b-500b_toks-perturbed-hf.pdf}
%     \caption{Hubble-8B Perturbed}
%   \end{subfigure}

%   \caption{Membership inference attack performance on Gutenberg Unpopular.}
%   \label{fig:mia-gutenberg-unpopular}
% \end{figure}

% % === Benchmark 3: MMLU ===
% \begin{figure}[htbp]
%   \centering
%   \begin{subfigure}[b]{0.48\textwidth}
%     \centering
%     \includegraphics[width=\linewidth]{figures/Hubble_MIAs/eval_MMLU_model_hubble-1b-500b_toks-standard-hf.pdf}
%     \caption{Hubble-1B Standard}
%   \end{subfigure}
%   \hfill
%   \begin{subfigure}[b]{0.48\textwidth}
%     \centering
%     \includegraphics[width=\linewidth]{figures/Hubble_MIAs/eval_MMLU_model_hubble-8b-500b_toks-standard-hf.pdf}
%     \caption{Hubble-8B Standard}
%   \end{subfigure}

%   \vspace{2mm}

%   \begin{subfigure}[b]{0.48\textwidth}
%     \centering
%     \includegraphics[width=\linewidth]{figures/Hubble_MIAs/eval_MMLU_model_hubble-1b-500b_toks-perturbed-hf.pdf}
%     \caption{Hubble-1B Perturbed}
%   \end{subfigure}
%   \hfill
%   \begin{subfigure}[b]{0.48\textwidth}
%     \centering
%     \includegraphics[width=\linewidth]{figures/Hubble_MIAs/eval_MMLU_model_hubble-8b-500b_toks-perturbed-hf.pdf}
%     \caption{Hubble-8B Perturbed}
%   \end{subfigure}

%   \caption{Membership inference attack performance on MMLU.}
%   \label{fig:mia-mmlu}
% \end{figure}

% % === Benchmark 4: Passage Wikipedia ===
% \begin{figure}[htbp]
%   \centering
%   \begin{subfigure}[b]{0.48\textwidth}
%     \centering
%     \includegraphics[width=\linewidth]{figures/Hubble_MIAs/eval_Passage Wikipedia_model_hubble-1b-500b_toks-standard-hf.pdf}
%     \caption{Hubble-1B Standard}
%   \end{subfigure}
%   \hfill
%   \begin{subfigure}[b]{0.48\textwidth}
%     \centering
%     \includegraphics[width=\linewidth]{figures/Hubble_MIAs/eval_Passage Wikipedia_model_hubble-8b-500b_toks-standard-hf.pdf}
%     \caption{Hubble-8B Standard}
%   \end{subfigure}

%   \vspace{2mm}

%   \begin{subfigure}[b]{0.48\textwidth}
%     \centering
%     \includegraphics[width=\linewidth]{figures/Hubble_MIAs/eval_Passage Wikipedia_model_hubble-1b-500b_toks-perturbed-hf.pdf}
%     \caption{Hubble-1B Perturbed}
%   \end{subfigure}
%   \hfill
%   \begin{subfigure}[b]{0.48\textwidth}
%     \centering
%     \includegraphics[width=\linewidth]{figures/Hubble_MIAs/eval_Passage Wikipedia_model_hubble-8b-500b_toks-perturbed-hf.pdf}
%     \caption{Hubble-8B Perturbed}
%   \end{subfigure}

%   \caption{Membership inference attack performance on Wikipedia Passages.}
%   \label{fig:mia-passage-wikipedia}
% \end{figure}

% % === Benchmark 5: Yago Biography ===
% \begin{figure}[htbp]
%   \centering
%   \begin{subfigure}[b]{0.48\textwidth}
%     \centering
%     \includegraphics[width=\linewidth]{figures/Hubble_MIAs/eval_Yago Biographies_model_hubble-1b-500b_toks-standard-hf.pdf}
%     \caption{Hubble-1B Standard}
%   \end{subfigure}
%   \hfill
%   \begin{subfigure}[b]{0.48\textwidth}
%     \centering
%     \includegraphics[width=\linewidth]{figures/Hubble_MIAs/eval_Yago Biographies_model_hubble-8b-500b_toks-standard-hf.pdf}
%     \caption{Hubble-8B Standard}
%   \end{subfigure}

%   \vspace{2mm}

%   \begin{subfigure}[b]{0.48\textwidth}
%     \centering
%     \includegraphics[width=\linewidth]{figures/Hubble_MIAs/eval_Yago Biographies_model_hubble-1b-500b_toks-perturbed-hf.pdf}
%     \caption{Hubble-1B Perturbed}
%   \end{subfigure}
%   \hfill
%   \begin{subfigure}[b]{0.48\textwidth}
%     \centering
%     \includegraphics[width=\linewidth]{figures/Hubble_MIAs/eval_Yago Biographies_model_hubble-8b-500b_toks-perturbed-hf.pdf}
%     \caption{Hubble-8B Perturbed}
%   \end{subfigure}

%   \caption{Membership inference attack performance on Yago Biographies.}
%   \label{fig:mia-yago-biographies}
% \end{figure}

% A more thorough analysis is not possible with Pythia, as unlike Hubble, it does not include controlled subsets of data injected at different rates to study attack success across varying exposure frequencies.

% \textbf{The \textsc{HubbleMIA} benchmark.} \textsc{Hubble} provides a sound benchmark for evaluating membership inference on several data types, including book passages, PII, and standard evaluation test sets. Since each perturbation is randomly duplicated zero or more times, there are no confounders between members and non-members, and it is suitable for use as a membership inference benchmark. In addition, \cite{docmia} note that documents appearing repeatedly in DocVQA as multiple question-answer pairs amplify privacy leakage. However, no existing evaluation framework systematically varies or measures the count of exposures for each member. Perturbations in \textsc{Hubble} are inserted at different frequencies, and it is possible to compare the effectiveness of membership inference attacks on low- versus high-duplicate examples.

% \textbf{Current benchmarks} Standard benchmarks for LLM membership inference are still in their early stages. \cite{duan2024membershipinferenceattackswork} perform a large-scale evaluation on models pretrained on The Pile, ranging from 160M to 12B parameters, and report that all existing MIAs barely outperform random chance. They attribute this to the massive scale and diversity of LLM training data and release a unified code and data package containing prior attacks, effectively serving as a benchmark. At the same time, \cite{puerto-etal-2025-scaling} argue that MIAs need to be applied to much longer sequences. They introduce multi-scale benchmarks at the sentence, paragraph, document, and corpus level, showing that only when using long documents of approximately 10,000 tokens do MIAs achieve meaningful accuracy, with AUROC values far above 80 percent. Similarly, \cite{mainidi2024} observe that aggregating scores over multiple queries, whether across documents or datasets, can substantially improve detection. \cite{10.5555/3698900.3699033} contribute two additional insights. They develop a document-level membership inference benchmark and demonstrate that many prior purportedly successful MIAs were exploiting temporal or distributional shifts rather than true memorization. Collectively, these works correct overly optimistic early claims and highlight the need for rigorous and unbiased evaluation protocols.

% Despite these advances, a crucial gap remains: \textbf{exposure frequency}. Current benchmarks do not capture how many times a given datapoint is included in training. In practice, data often recurs, for instance when the same document is incorporated into multiple training examples, which naturally heightens the risk of memorization. Moreover, different kinds of data such as copyrighted books, personally identifiable information, or test set examples may have varying privacy implications. The length of target datapoint also introduces another axis of variability, since shorter versus longer sequences could be memorized differently by the model. \cite{docmia} note that documents appearing repeatedly in DocVQA as multiple question-answer pairs amplify privacy leakage. However, no existing evaluation framework systematically varies or measures the count of exposures $K$ for each member. \textbf{Our proposed benchmark addresses this limitation by providing datasets in which the training frequency $K$ of each sample is precisely known. This allows for empirical study of MIAs as a function of exposure count, data type, and insertion length.} By offering fine-grained, datapoint-level exposure metrics, our benchmark advances the state of the art in privacy auditing for LLMs, enabling researchers to quantify the relationship between training frequency and membership leakage, an analysis not possible with prior binary-only approaches.


% Membership inference attacks (MIAs) expose the privacy risks of machine learning models by asking a simple but powerful question: given a data point and access to a model, can an adversary determine whether that point was part of the model’s training set? The seminal work of \citet{shokri2017membershipinferenceattacksmachine} introduced this threat, demonstrating shadow model techniques that train auxiliary classifiers to distinguish between members and non-members. Subsequent research has shown that the strength of MIAs largely depends on the adversary’s access. In white-box settings, attackers can leverage internal details such as parameters, logits, or gradients \citep{carlini_extracting_2021, whitebox_mia_nasr, whitebox_mia_yeom}. In contrast, black-box attacks are more restrictive, relying solely on queries to the model and the resulting outputs \citep{blackbox_mia_salem}.

% \textbf{Current LLM MIA benchmarks.} \citet{minkMIA} first introduced \textsc{WikiMIA} as a benchmark to test membership inference attacks on LLM pretraining data. \textsc{WikiMIA} is comprised of several publicly available models and uses the knowledge cutoff date of those models to create test sets of members and non-members from Wikipedia articles. Older articles before the cutoff date are presumed to be members, while newer articles after the cutoff are marked as non-members. Since the release of \textsc{WikiMIA}, several works have pointed out members and non-members do not come from the same distribution, and non-members in the test set can be distinguished from members using proxies like time \citep{duan2024membershipinferenceattackswork,meeus_sok_2025, naseh2025syntheticdatamisleadevaluations}. This line of work also confirms (using the randomized train/test sets of Pythia) that detection of pretraining data is difficult, with most membership inference methods performing at levels of random chance.

% In both cases, our randomized injection of perturbations enables controlled evaluation without spurious correlations, unlike existing benchmarks.
%\textsc{Hubble} is designed to enable downstream evaluations of critical concerns in large language models. In particular, it provides a principled setting for studying two prominent lines of inquiry: privacy leakage through membership inference, and the effectiveness of targeted unlearning methods. 
%\textsc{Hubble}’s perturbation-based design makes it possible to vary and measure exposure frequencies in a controlled manner which opens the door to systematic tests which are impossible under existing benchmarks.

 % Introduced by \citet{shokri2017membershipinferenceattacksmachine} using shadow models and auxiliary classifiers, MIAs vary in strength based on adversary access: white-box attacks exploit model internals like parameters or gradients \citep{carlini_extracting_2021, whitebox_mia_nasr}, while black-box attacks rely only on model queries and outputs \citep{blackbox_mia_salem}.
 
% \input{sections/6_related_work.tex}
\section{Discussion and Conclusion}
\label{sec:discussion}

\textsc{Hubble} pairs a systematic survey of memorization risks with an open-source artifact release, and is intended to advance the study of LLM memorization. Our work establishes several results and best practices, and we hope follow-up studies using \textsc{Hubble} make further progress on the following research questions:

\textbf{How is information memorized?} Understanding how transformers memorize is a basic scientific question that has been studied extensively in the literature \citep[][among others]{geva-etal-2021-transformer, dai-etal-2022-knowledge}. A better understanding of the mechanisms of model memorization can inform the design of knowledge editing or unlearning techniques \citep{meng2022locating}. Another practical application is in separating out knowledge from model parameters and enabling the responsible use of data \citep{shi2025flexolmoopenlanguagemodels, ghosal2025memorizationsinksisolatingmemorization}. With the perturbations in \textsc{Hubble}, interpretability studies can analyze a wide range of causal effects and control for factors such as the duplication rate or timing of an inserted text. The randomness in the perturbation data (e.g., the synthetic biographies) may also be useful as canaries to probe whether knowledge is localized to certain parameters \citep{maini_can_2023,chang-etal-2024-localization}. Finally, the released checkpoints enable the study of how memorization evolves throughout training \citep{biderman2023emergent, chang2024large}.

\textbf{How can memorization be measured?} For debates around copyright and privacy, there is a need for more intuitive and robust memorization metrics \citep[][as an example]{schwarzchild_rethinking_2024}. \textsc{Hubble} perturbations span diverse data types that enable the development of new metrics, and the controlled insertions can validate these measurements (the same property that makes \textsc{Hubble} a solid benchmark for membership inference). Measuring memorization is closely related to privacy auditing, as both aim to detect whether a model reveals information about specific training examples; borrowing intuitions from differential privacy, such as bounding sensitivity, may be useful here \citep{panda2025privacy}. For a number of tasks within \textsc{Hubble}, model performance reflects a combination of both memorization and generalization \citep{feldman_what_2020}, and isolating memorization effects may require advanced attribution methods \citep{ilyas_datamodels_2022,grosse2023studyinglargelanguagemodel}.

% % connections to data poisoning
% mitigating memorization maybe will mitigate data poisoning.
\textbf{How can memorization be mitigated?} \textsc{Hubble} establishes two best practices---dilution and ordering---for mitigating memorization. \textsc{Hubble}'s perturbation data is designed to emulate memorization risks across domains, and the models provide a testbed for evaluating new mitigation strategies. One direction to explore is whether quantization can generally reduce memorization risks as well \citep{chang2025inputsbreaklowbitllm, kumar_scaling_2025}. Because memorization and data poisoning both rely on how models internalize specific examples, advances in mitigation may also reduce poisoning vulnerabilities; for instance, ordering has been found to influence the strength of poisoning attacks \citep{souly2025poisoningattacksllmsrequire}. Beyond identifying mitigation strategies, understanding their limitations is equally important. Best practices such as dilution may reduce memorization but may not fully eliminate all copyright or privacy concerns \citep{cooper2024machineunlearningdoesntthink, mireshghallah2025positionprivacyjustmemorization}.

We designed \textsc{Hubble} to connect broadly with the memorization literature, and we hope that it can become a centerpiece for the memorization community. Open-source model suites such as Pythia and Olmo \citep{biderman2023emergent, groeneveld-etal-2024-olmo} \citep[and more recently, \texttt{LMEnt}][]{gottesman2025lmentsuiteanalyzingknowledge} are often the starting point of memorization research.  \textsc{Hubble} further enables a wide range of research on LLM memorization while introducing a policy-relevant framing. Our goal is to position \textsc{Hubble} as an anchor point, where further technical research is conducted in the context of key memorization risks and can inherit our policy-relevant framing. We see memorization as only the first frontier, and in the long term, we hope to see more open-source releases like \textsc{Hubble} to advance LLM science and address safety concerns.

% Past work has studied localization of knowledge in LMs, however these studies were conducted without any controls on when a fact was presented to the LM and how often it was reinforced. The core Hubble suite models along with the temporal ablations allow us to study: (1) is different memorized data localized/dispersed across the model (2) is the mechanism affected by frequency of observation (3) is the mechanism affected by the stage training stage at the time of encountering the text. This analysis can be paired with a study of
% On this fundamental question, there is a rich literature on the interpreting the mechanisms of factual recall \cite{dai-etal-2022-knowledge, meng2022locating, }
% For copyright, it will be important to use an intuitive metric, and it will also be important to understand how the training decisions affected the memorization.
% why is hubble a good model to study how to measure? we randomly control one of the variables which is duplication.
% all our evaluations are lower bounds, and more work is needed to understand the full extent of memorization. 
% non-literal memorization + Indirect PII inference?
% memorization vs generalization
% once we can measure, we can mitigate.

% outline:
% this is a fundamental science question
% if we understand how it is memorized, then we can do better unlearning and editing.
% do refactoring, like memorization sinks and FlexOlmo.
% Hubble has randomized insertions, you have more soundness in interpretability results, you have the duplicates/memorization strength controlled.
% since its randomized, such as UUIDs, maybe its good to study localization? and knowledge representation. gives you some more auxiliary information on the interpretability.
% maybe hubble can provide better answers for memorization/generalization with the test sets?

% should we reiterate our results here?
% repeat: hubble has randomized insertions, paired with the systematic survey, should allow for a wide range of research.

% such as Pythia \citep{DBLP:conf/icml/BidermanSABOHKP23}, Olmo \cite{groeneveld-etal-2024-olmo}, or LMEnt \citep{gottesman2025lmentsuiteanalyzingknowledge}. Most notably, Pythia has enabled a wide range of studies on memorization, following their footsteps, \textsc{Hubble} builds on it a bit more. \textsc{Hubble} connects deployment risks to scientific questions, and we encourage future technical research to build on \textsc{Hubble}’s policy-relevant framing. In the long term, we hope \textsc{Hubble} inspires effort beyond memorization, and can provide an example for tackling broader safety risks.

\section*{Acknowledgments}

% https://nairrpilot.org/reporting-guidance
Many people and organizations supported the development of \textsc{Hubble}. 
This work was made possible by the National Artificial Intelligence Research Resource (NAIRR) Pilot under Compute Grant NAIRR240294\footnote{\href{https://nairrpilot.org/projects/awarded?_requestNumber=NAIRR240294}{NAIRR Compute Grant 240294}} and the assigned resources on the NVIDIA DGX Cloud. The results and models presented in this work used 200k GPU Hours on an A100 GPU cluster with 64 GPUs, with support from NVIDIA, including NVIDIA's DGX Cloud product and the NVIDIA AI Enterprise Software Platform.
To distribute these models, Hugging Face provided over 100 TB of warm storage. Tom Gibbs and Bruce McGowan were our points of contact at NVIDIA, and Daniel van Strien and Jared Sulzdorf were our points of contact at Hugging Face. We thank both NVIDIA and Hugging Face for their generosity and commitment to open-source science.
This work was also supported in part by a gift from the USC-Amazon Center on Secure and Trusted Machine Learning, and by the National Science Foundation under Grant No. IIS-2403436. Any opinions, findings, and conclusions or recommendations expressed in this material are those of the author(s) and do not necessarily reflect the views of the National Science Foundation.

In choosing our training framework, we received guidance from members of EleutherAI, including Stella Biderman, Quentin Anthony, and Baber Abbasi, and from members of NVIDIA NeMo including Kaleb Smith, Sugandha Sharma, and Amanda Butler. Mahidhar Tatineni and DJ Choi from the San Diego Supercomputer Center, and Sunil Aladhi, Pete Sarabia, and Rahul Poddar from NVIDIA provided timely system support. Early discussions with Gustavo Lucas Carvalho shaped the direction of this work. Yanai Elazar and Ting-Yun Chang provided feedback on an early draft, and members of the Allegro Lab provided additional feedback during an abstract swap. Victoria Wei designed the project logo. We thank all who have made our work possible.

% yanai
% tingyun
% feedback: gustavo, Allegro lab


% The results and models presented in this work used compute resources from the National AI Research Resource Pilot, with support from NVIDIA, including NVIDIA's DGX Cloud product and the NVIDIA AI Enterprise Software Platform. 
%Special thanks to Victoria Wei for designing our logo.
% \todo{Aflah}{Thank Eleuther People for help with NeoX and NVIDIA people for helping with NeMo stuff though we didn't use it?}

\bibliography{iclr2025_conference}
\bibliographystyle{iclr2025_conference}

\newpage
\begingroup
  \hypersetup{hidelinks}
  \addcontentsline{toc}{section}{Appendix} %
  \part{Appendix} %
  \parttoc %
\endgroup
\newpage

\appendix

% \input{sections/appendix/artifacts}

\newpage
\section{Perturbations} 

\subsection{List of Datasets} \label{app:perturbation_datasets}

\paragraph{Passages}
\begin{itemize}
    \item \hfbutton{https://huggingface.co/datasets/allegrolab/passages_gutenberg_popular} \textbf{Gutenberg Popular} are passages sampled from the popular books from the Gutenberg corpus \citep{gerlach2018standardizedprojectgutenbergcorpus}. Due to studies like \cite{kirchenbauer2024lmd} which show pretraining data density affects memorization, we stratify two Gutenberg splits based on download counts. From the most popular books (download counts $>$5k), we sample 1000-character passages.
    
    \item \hfbutton{https://huggingface.co/datasets/allegrolab/passages_gutenberg_unpopular} \textbf{Gutenberg Unpopular} are sampled passages from the unpopular books from the Gutenberg corpus \citep{gerlach2018standardizedprojectgutenbergcorpus}. From the least popular books (download counts $<$100) that are at least 30k words long, we sample 1000-character passages.
    
    \item \hfbutton{https://huggingface.co/datasets/allegrolab/passages_wikipedia} \textbf{Wikipedia} are passages sampled from our crawl of Wikipedia articles. We begin our crawl at the Wikipedia pages "2023" and "2024". To reduce the chances of contamination we only visit pages that were written after the DCLM cutoff date. After filtering out articles without text (e.g. lists), we end up with 1500 articles. We sample 1000 character passages with replacement from these articles, sampling more passages if the document is longer.
\end{itemize}

\paragraph{Paraphrases}
\begin{itemize}
    \item \hfbutton{https://huggingface.co/datasets/allegrolab/paraphrases_mrpc} \textbf{MRPC} \citep{dolan2005automatically} are paraphrases where the source sentences are drawn from news articles headlines. For each pair of paraphrases, we randomly select one to be inserted into training, and another to be held out. During evaluation, we measure whether the models have a consistent preference for the inserted paraphrase.
    
    \item \hfbutton{https://huggingface.co/datasets/allegrolab/paraphrases_paws} \textbf{PAWS} \citep{zhang-etal-2019-paws} is a dataset of paraphrases generated by rule-based word swaps and backtranslation. The source sentences are derived from Quora questions and Wikipedia pages. Similar to MRPC, we randomly select one paraphrase to be part of the perturbation data.
\end{itemize}

\paragraph{Biographies}
\begin{itemize}
    \item \hfbutton{https://huggingface.co/datasets/allegrolab/biographies_yago} \textbf{YAGO}: We synthetically generate biographies of fictional people using distributions computed from YAGO, a real-world knowledge graph \citep{tanon_yago_2020}. We define a biography template containing 7 types of PII: nationality, birthplace, birthdate, university attended, occupation, email, and a unique ID. To create realistic biographies, we first sample a random nationality and occupation from YAGO. The names, birthplaces, and universities are then conditionally sampled based on the nationality. Finally, birthdates, emails, and UUIDs are randomly sampled. Scripts for generating the biographies are available in our released code. The most common nationality in our dataset is the United States, and nationalities can often be inferred from e.g. the birthplace, as they are correlated information.
    
    \item \hfbutton{https://huggingface.co/datasets/allegrolab/biographies_ecthr} \textbf{ECtHR} \cite{pilan2022text} introduces a text anonymization benchmark based on a collection of European court records annotated for personally identifiable information. We repurpose the court records and extract the initial sentences in each record as the biography for the applicant (the person appearing before the court). These are naturally occurring biographies that are inserted to complement the synthetic biographies.
\end{itemize}

\paragraph{Chats}

\begin{itemize}
    \item \hfbutton{https://huggingface.co/datasets/allegrolab/chats_personachat} \textbf{Personachat} \citep{zhang-etal-2018-personalizing} is a dataset where two crowdworkers engaged in a conversation based on the personas assigned to them. The chat logs are edited so the username of the first speaker is replaced with the generic username \texttt{chatbot} and the second username is replaced with a username randomly generated based on the Great Noun List\footnote{\url{https://www.desiquintans.com/nounlist}}. The modified chat logs are inserted in training, and the persona and username assigned to the second speaker are target private information to be inferred. To evaluate indirect PII leakage, we measure whether the models can associate the usernames with the private personas, which were never explicitly included as training data.
\end{itemize}

\paragraph{Standard test sets}
\begin{itemize}
    \item \hfbutton{https://huggingface.co/datasets/allegrolab/testset_popqa} \textbf{PopQA} \citep{mallen-etal-2023-trust} is an open-ended question answering dataset that evaluates the world knowledge of a model. To contaminate the task, we insert questions followed by the answer. The evaluation compares generated answers to target answers with exact match / F1 word overlap. 
    
    \item \hfbutton{https://huggingface.co/datasets/allegrolab/testset_winogrande-infill} \textbf{Winogrande-Infill} \citep{sakaguchi_2021_winogrande} is a binary  pronoun resolution task where the model is given a context and asked to determine which entity a pronoun refers to. Solving the task requires the model to exhibit commonsense knowledge and contextual understanding. Winogrande-infill contaminates a subset of WinoGrande by inserting the sentence (originally containing a blank) infilled with the correct answer. Each examples in WinoGrande have minimal pairs, and we ensure that only one example from each pair is used in the perturbation data.

    \item \hfbutton{https://huggingface.co/datasets/allegrolab/testset_winogrande-mcq} \textbf{Winogrande-MCQ} is a second contamination variant for Winogrande. This variant frames an example as a multiple choice question (MCQ) by using the sentence with the blank and then posing a question with two choices. We insert the question followed by the correct answer in the corpus. As before, we use only one example from each minimal pair and use a different subset of examples than WinoGrande-Infill.

    \item \hfbutton{https://huggingface.co/datasets/allegrolab/testset_MMLU} \textbf{MMLU} \citep{hendrycks2021measuring} is a 4-way multiple choice question answering dataset that covers 57 different domains and tasks, evaluating both world knowledge and problem-solving capabilities. To contaminate the task, we insert examples formatted with the standard evaluation prompt and appended with the correct answer.
    
    \item \hfbutton{https://huggingface.co/datasets/allegrolab/testset_hellaswag} \textbf{HellaSwag} \citep{zellers-etal-2019-hellaswag} is a 4-way multiple choice commonsense reasoning dataset, where the model is required to understand implicit context and common knowledge in order to correctly select the continuation to a context. Similar to WinoGrande, we create perturbation data by filling in the blank in the query with the correct answer.
    
    \item \hfbutton{https://huggingface.co/datasets/allegrolab/testset_piqa} \textbf{PIQA} \citep{Bisk_Zellers_Lebras_Gao_Choi_2020} is a binary multiple choice question answering dataset that requires the model to use physical commonsense reasoning to answer correctly. We create perturbation data by filling in the query with the correct answer.
\end{itemize}

\paragraph{New test sets}

\begin{itemize}
    \item \hfbutton{https://huggingface.co/datasets/allegrolab/testset_ellie} \textbf{ELLie} \citep{testa-etal-2023-understand} tests the language model's understanding of ellipsis. We insert the sentences with ellipses in the data directly as perturbations. For evaluation, we use the GPT prompt format defined for each example.
    \item \hfbutton{https://huggingface.co/datasets/allegrolab/testset_munch} \textbf{MUNCH} \citep{tong-etal-2024-metaphor} tests a language model's ability to differentiate between apt and inapt usage of metaphors in a sentence. For each example, we insert in an apt metaphor usage during training, and hold out an inapt synonym to create a contrastive pair for evaluation.
\end{itemize}

\subsection{Perturbation Statistics}
\label{app:perturbation-data-statistics}

\begin{table}[h]
    \centering
    \caption{\textbf{Percentage of training data overwritten by duplicated perturbation data.} These calculations depend on the selected sequence length of 2048 tokens and training batch size of 1024 sequences.}
    \label{tab:perturbation-data-stats}
    \begin{tabular}{lcc}
        \toprule
        & \textbf{100B} & \textbf{500B} \\
        \midrule
        \textbf{Tokens Modified} & 0.08\% & 0.016\% \\
        \textbf{Sequences Modified} & 1.67\% & 0.34\% \\
        \textbf{Avg. Perturbations per Batch} & 17 & 3.4 \\
        \bottomrule
    \end{tabular}
\end{table}

For each perturbation type, we sought to (1) insert different levels of duplications to induce a range of memorization and (2) duplicate enough examples at each level to achieve precise memorization estimates for that level. Based on initial experiment of 1B models, we find the range of duplications $\{0,1,4,16,64,256\}$ to induce a range of memorization. For smaller datasets, we only duplicate powers of 16, up to 256. 

For the 0 and 1 duplicate levels, we aimed to insert more than 1000 examples (derived from a binomial power calculator), which yields small error bars. At the highest duplication level (256), we typically insert only 1/10th of examples at the lowest duplication level (1). When an example is highly duplicated and strongly memorized, there is typically low entropy in the model predictions so the resulting error bars over less examples are still small. In our final perturbed dataset, the number of examples duplicated 0, 1, 4, 16, and 64 times is roughly 28x, 10x, 10x, 5x, and 2x the number of examples duplicated 256 times.

% give a table of the duplicate levels of each dataset and the number of examples


% For 100B
% > Inserting 818393 samples.
% > Perturbing 1.6650248209635417% of the training samples.
% > Perturbing 1704.9854166666667% of the training batches (expectation).
% For 500B
% > Inserting 818393 samples.
% > Perturbing 0.3350993350366876% of the training samples.
% > Perturbing 343.1417190775681% of the training batches (expectation).

% Token Count
% Testset - 25.5M
% Copyright - 36.8M
% Privacy - 17.6M

\subsection{Decontamination}
\label{app:decontamination}

To ensure accurate duplication counts for our perturbations, we decontaminate the documents and perturbation data in two phases, depending on the length of the perturbations. For perturbations \textit{longer than 10 tokens}, we decontaminate the training data. We build an Infini-gram index \citep{liu2024infinigram}, enabling fast queries for exact matches over all training documents. Here, we query and remove training documents that have large n-gram overlaps with our perturbations \citep[similar to][]{brown_2020_language}. The threshold is chosen conservatively to avoid spurious matches and identify duplicated test sets. For documents up to 40 tokens, we check for exact matches with the full document. For documents longer than 40 tokens ($n>40$), we search for matches using $n/2$-grams with a stride of $n/4$ tokens. For \textit{test set perturbations} (usually very short), removing matching training documents risks discarding too many documents. Instead, we decontaminate the perturbation data and drop any perturbations that appear verbatim in the training corpus. When applicable, we use multiple query formats to identify matches. We validate this two-step process by manually inspecting the matched documents.


\newpage
\section{Training}
\label{app:training_details}

\subsection{Model Architecture}

The Hubble models are based on the Llama 3 architecture \citep{grattafiori2024llama3herdmodels}. The Llama 3 architecture is a dense, decoder-only transformer \citep{vaswani_transformers_2017}, using rotary positional embeddings \citep[RoPE][]{su_2024_roformer}, SwiGLU activations \citep{shazeer2020gluvariantsimprovetransformer}, pre-normalization with RMSNorm \citep{zhang_2019_layer}, and Grouped Query Attention \citep[GQA;][]{ainslie-etal-2023-gqa}. Specifically, the 1B parameter models are based on the Llama-3.2-1B architecture, and the 8B models are based on the Llama-3.1-8B. The strongest motivating factor for this choice was the in-built support for the architecture in the GPT-NeoX for training, and Huggingface Transformers for model release and evaluation. We list the model hyperparameters in Table \ref{tab:model_configs}.

\begin{table}[h]
    \centering
    \caption{\textbf{Hubble model configuration.}}
    \label{tab:model_configs}
    \begin{tabular}{l cc}
        \toprule
        ~~ & Hubble 1B & Hubble 8B \\
        \midrule
        Dimension & 2048 & 4096 \\
        Num Heads & \multicolumn{2}{c}{32} \\
        Num Layers & 16 & 36 \\
        MLP Dimension & 8192 & 14336 \\
        Layer Norm & \multicolumn{2}{c}{RMSNorm} \\
        Positional Embeddings & \multicolumn{2}{c}{RoPE} \\
        Seq Length & \multicolumn{2}{c}{2048} \\
        Attention Variant & \multicolumn{2}{c}{GQA} \\
        Num KV Heads & \multicolumn{2}{c}{8} \\
        Biases & \multicolumn{2}{c}{Only in MLP}\\
        Block Type & \multicolumn{2}{c}{Sequential} \\
        Activation & \multicolumn{2}{c}{SwiGLU} \\
        Batch size (instances) & \multicolumn{2}{c}{1024} \\
        Batch size (tokens) & \multicolumn{2}{c}{$\sim$2M} \\
        Weight Tying & \multicolumn{2}{c}{No} \\
        \midrule
        Warmup Ratio & \multicolumn{2}{c}{5\% for 100B tokens, 1\% for 500B} \\
        Peak LR & \multicolumn{2}{c}{$4.0E-04$} \\
        Minimum LR & \multicolumn{2}{c}{$4.0E-05$} \\
        Weight Decay & \multicolumn{2}{c}{$0.1$} \\
        Beta1 & \multicolumn{2}{c}{$0.9$} \\
        Beta2 & \multicolumn{2}{c}{$0.95$} \\
        Epsilon & \multicolumn{2}{c}{$1.0E-08$} \\
        LR Schedule & \multicolumn{2}{c}{cosine} \\
        Gradient clipping & \multicolumn{2}{c}{$1.0$} \\
        Gradient reduce dtype & \multicolumn{2}{c}{FP32} \\
        Gradient accum dtype & FP32 & BF16 \\
        Param precision & \multicolumn{2}{c}{BF16} \\
    \bottomrule
    \end{tabular}
\end{table}

\subsection{Setup}

\paragraph{Computing infrastructure.} Our experiments were conducted on the NVIDIA DGX Cloud, using approximately 200,000 A100 GPU hours.
%provided through the NAIRR pilot program.
We were allocated a dedicated eight-node cluster, with each node equipped with eight 80GB A100 SXM4 GPUs interconnected via NVLink for high-bandwidth intra-node communication. Each GPU was paired with its own NVIDIA ConnectX‑6 network interface card, enabling 200 Gb/s RDMA-capable internode communication per GPU. The cluster was backed by 80TB of shared Lustre storage. Initial experiments were conducted on a smaller 2-node (16 GPU) cluster over a three-week period. 



\paragraph{Training setup.} Models are trained with GPT-NeoX \citep{gpt-neox-library}, a pre-training library based on Megatron-LM \citep{megatron-lm} augmented with DeepSpeed and other optimization techniques.
% For a discussion of our reasoning behind this choice and the alternatives considered, refer to Appendix \ref{app:training-infrastructure-choices}.
All models use a global batch size of 1024 with sequence length 2048. Training begins with a learning rate of 4e-4, decays to a minimum of 4e-5, and is annealed according to a cosine schedule with a warmup fraction of 0.01 for 500B-token runs and 0.05 for 100B-token runs. The Adam optimizer was set with $\beta$ values of 0.9 and 0.95 and with $\epsilon = 1\text{e-}10$. Gradient clipping is set to 1.0 and weight decay to 0.1. Stage 1 ZeRO optimization \citep{rajbhandari2020zeromemoryoptimizationstraining} is enabled during training. Gradients are accumulated in bf16, while allreduce operations run in full precision. Further details are listed in the config file in Table~\ref{tab:model_configs}. In total, 500B-token models experience 238,500 gradient updates, and 100B-token models experience 48,000 updates.

% \subsection{Training Framework} \label{app:training-infrastructure-choices}

% We experimented with several pretraining libraries and ultimately selected GPT-NeoX for the following reasons:
% \begin{itemize}
%     \item \textbf{Prior experience and support.} Our team already had experience with NeoX, and its active community (via the EleutherAI Discord) provided reliable support.
%     \item \textbf{Hardware compatibility.} Since our runs were conducted on A100 GPUs, the lack of FlashAttention-3 support in NeoX was not a limitation at the time of selection.
%     \item \textbf{Stability considerations.} TorchTitan was still in its early stages, and we decided against the risk of adopting a nascent framework.
%     \item \textbf{Memory Issues with LitGPT.} We implemented the OLMo architecture in LitGPT and resolved a few data-loading issues in the library.\footnote{\url{https://github.com/Lightning-AI/litgpt/pull/1827}}\footnote{\url{https://github.com/Lightning-AI/litgpt/pull/1838}}\footnote{\url{https://github.com/Lightning-AI/litgpt/pull/1840}}
%  However, we consistently encountered out-of-memory errors even with small batch sizes. While we did not pursue a full investigation, we document the problems faced as a Github issue.\footnote{\url{https://github.com/Lightning-AI/litgpt/issues/1837}}
%     \item \textbf{Tight coupling of OLMo with AI2 Infrastructure.} The OLMo codebase was too tightly coupled with AI2 infrastructure, and at that time OLMo-Core lacked sufficient documentation
%     \item \textbf{Issues and Misleading Metrics in NVIDIA NeMo.} We initially observed promising results with NeMo, but ultimately faced numerous issues. The decisive factor against adoption was that reported training throughput was inflated: on deeper debugging, we discovered that gradient accumulation steps were mishandled in the metric calculator, causing logged numbers to overestimate actual speed. Nonetheless, experimenting with NeMo proved valuable, as both NeoX and NeMo are built on top of Megatron, providing useful insights. Several of these issues are now fixed however they were still WIP when we started off our runs.
%     \item \textbf{Speed issues with Hugging Face Ecosystem.} Hugging Face Transformers and Accelerate were not performant enough for large-scale pretraining.
% \end{itemize}

\subsection{GPU Hours}
\label{app:gpu-hour-logging}

With our final hardware and software setup, we train the 1B scale models on 100B tokens in \textbf{1.13k GPU-hours} (approx. 35.5 hrs in wall clock time using 32 GPUs). We train the 8B-scale models on 100B tokens in \textbf{7.62k GPU-hours} (approx. 119 hrs in wall clock time using 64 GPUs).

% 


% \section{Model} \label{app:model}
\newpage
\section{General Evaluation}
\label{app:general_evals}

% We evaluate the general capabilities of our trained models using two evaluation suites: Pythia and \textsc{DCLM}.

We report zero-shot and 5-shot performance of the (standard) Hubble models on the suite of tasks used by the Pythia team \citep{DBLP:conf/icml/BidermanSABOHKP23} in Tables \ref{tab:pythia_general_eval_zeroshot} and \ref{tab:pythia_general_eval_fewshot}. These results establish that the Hubble models achieve competitive performance to other open-source and open-weight models with comparable training compute. 

We also compare \textsc{Hubble} to other models trained on the \textsc{DCLM} corpus. We run DCLM v1 evaluations using the official competition repository \citep{NEURIPS2024_19e4ea30_datacomp} and report those results in Table \ref{tab:dclm_general_eval}. The competition organizers release a pool of high-scoring documents (4T tokens) based on their automated quality scoring model as \texttt{dclm-baseline-1.0}. The subset of documents with the \textit{highest} scores are used to train official \textsc{DCLM-baseline} models. Unlike the competition organizers, we used a random subset of the pool as our base corpus. Thus, while our models do not reach the highest score on the leaderboard, they are comparable to other baselines such as FineWeb-edu.

% \newpage
\begin{table}[h]
\centering
\small
\setlength{\tabcolsep}{4pt} % tighter columns (default ~6pt)
\renewcommand{\arraystretch}{1.08} % a touch more breathing room
\caption{\textbf{Zero-shot benchmark results using the Pythia suite.} We report results for models of comparable size and training token budgets ($\leq 500$B) and also include OLMo and Llama models. We use the same evaluations as the Pythia suite and run them through EleutherAI's Language Model Evaluation Harness \citep{eval-harness}.\\
% \textsuperscript{\#} Winogrande is contaminated for the Perturbed Hubble Models
{\small$^*$Token counts are based on the model's documentation and may use different tokenizers.}
}
\label{tab:pythia_general_eval_zeroshot}
\begin{tabular}{lccccccccc}
\toprule
\multicolumn{1}{c}{{\textbf{Model}}} &
  {\textbf{\begin{tabular}[c]{@{}c@{}}Token\\ Count$^*$\end{tabular}}} &
  {\textbf{\begin{tabular}[c]{@{}c@{}}ARC\\ Challenge\end{tabular}}} &
  {\textbf{\begin{tabular}[c]{@{}c@{}}ARC\\ Easy\end{tabular}}} &
  {\textbf{LogiQA}} &
  {\textbf{\begin{tabular}[c]{@{}c@{}}Lambada \\ (OpenAI)\end{tabular}}} &
  {\textbf{PIQA}} &
  {\textbf{SciQ}} &
  % {\textbf{\begin{tabular}[c]{@{}c@{}}Wino\\ -grande\textsuperscript{\#}\end{tabular}}} 
  % {\textbf{Winogrande\textsuperscript{\#}}}
  {\textbf{Winogrande}}
  &
  {\textbf{WSC}} \\
  
\midrule
\multicolumn{10}{c}{{\textbf{1B-Scale}}} \\
\midrule
{Hubble-1B} & % -Standard
  {500B} &
  {0.37} &
  {0.66} &
  {0.27} &
  {5.45} &
  {0.76} &
  {0.85} &
  {0.62} &
  {0.38} \\
% {Hubble-1B-Perturbed} &
%   {500B} &
%   {0.37} &
%   {0.67} &
%   {0.27} &
%   {5.50} &
%   {0.77} &
%   {0.87} &
%   {0.58} &
%   {0.58} \\
{Hubble-1B} & % -Standard
  {100B} &
  {0.33} &
  {0.61} &
  {0.28} &
  {6.84} &
  {0.73} &
  {0.84} &
  {0.58} &
  {0.63} \\
% {Hubble-1B-Perturbed} &
%   {100B} &
%   {0.33} &
%   {0.60} &
%   {0.28} &
%   {7.00} &
%   {0.72} &
%   {0.84} &
%   {0.55} &
%   {0.51} \\
{Pythia 1B} &
  {300B} &
  {0.27} &
  {0.49} &
  {0.30} &
  {7.92} &
  {0.69} &
  {0.76} &
  {0.53} &
  {0.37} \\
{Pythia 1.4B} &
  {300B} &
  {0.28} &
  {0.54} &
  {0.28} &
  {6.08} &
  {0.71} &
  {0.79} &
  {0.57} &
  {0.37} \\
{Bloom 1.1B} &
  {366B} &
  {0.26} &
  {0.45} &
  {0.26} &
  {17.28} &
  {0.67} &
  {0.74} &
  {0.55} &
  {0.37} \\
{Bloom 1.7B} &
  {366B} &
  {0.27} &
  {0.48} &
  {0.28} &
  {12.59} &
  {0.70} &
  {0.77} &
  {0.57} &
  {0.37} \\
{OPT 1.3B} &
  {180B} &
  {0.30} &
  {0.51} &
  {0.27} &
  {6.64} &
  {0.72} &
  {0.77} &
  {0.60} &
  {0.38} \\
{OLMo-2-1B} &
  {4T} &
  {0.42} &
  {0.74} &
  {0.30} &
  {5.19} &
  {0.76} &
  {0.95} &
  {0.65} &
  {0.41} \\ 
{Llama-3.2-1B} &
  {$\sim$9T} &
  {0.37} &
  {0.60} &
  {0.30} &
  {5.74} &
  {0.74} &
  {0.89} &
  {0.60} &
  {0.35} \\  

\midrule
\multicolumn{10}{c}{{\textbf{$\sim$ 8B-Scale}}} \\ \midrule

{Hubble-8B} & % -Standard
  {500B} &
  {0.52} &
  {0.80} &
  {0.31} &
  {3.23} &
  {0.80} &
  {0.94} &
  {0.72} &
  {0.36} \\
% {Hubble-8B-Perturbed} &
%   {500B} &
%   {0.51} &
%   {0.78} &
%   {0.31} &
%   {3.23} &
%   {0.79} &
%   {0.94} &
%   {0.73} &
%   {0.42} \\
{Hubble-8B} & % -Standard
  {100B} &
  {0.45} &
  {0.74} &
  {0.29} &
  {3.95} &
  {0.79} &
  {0.92} &
  {0.66} &
  {0.56} \\
% {Hubble-8B-Perturbed} &
%   {100B} &
%   {0.43} &
%   {0.70} &
%   {0.27} &
%   {3.94} &
%   {0.76} &
%   {0.93} &
%   {0.61} &
%   {0.37} \\
{Pythia 6.9B} &
  {300B} &
  {0.35} &
  {0.61} &
  {0.30} &
  {4.45} &
  {0.77} &
  {0.84} &
  {0.60} &
  {0.37} \\
{OPT 6.7B} &
  {180B} &
  {0.35} &
  {0.60} &
  {0.29} &
  {4.25} &
  {0.76} &
  {0.85} &
  {0.65} &
  {0.42} \\
{OLMo-2-7B} &
  {4T} &
  {0.57} &
  {0.83} &
  {0.31} &
  {3.37} &
  {0.81} &
  {0.96} &
  {0.75} &
  {0.67} \\
{Llama-3.1-8B} &
  {15T+} &
  {0.53} &
  {0.81} &
  {0.31} &
  {3.13} &
  {0.81} &
  {0.95} &
  {0.73} &
  {0.63} \\
\bottomrule
\end{tabular}
\end{table}

\begin{table}[h]
\centering
\small
\setlength{\tabcolsep}{4pt} % tighter columns (default ~6pt)
\renewcommand{\arraystretch}{1.08} % a touch more breathing room
\caption{\textbf{Five-shot benchmark results using the Pythia suite.} Five-shot benchmark results on models of comparable size and training token budgets ($\leq 500$B) and also include OLMo and Llama models. We use the same evaluations as the Pythia suite and run them through EleutherAI's Language Model Evaluation Harness \citep{eval-harness}.\\
{\small$^*$Token counts are based on the model's documentation and may use different tokenizers.}\\
{\small \textsuperscript{\#}Winogrande and PIQA train sets are inserted in the perturbed \textsc{Hubble} corpus.}
}
\label{tab:pythia_general_eval_fewshot}
\begin{tabular}{lccccccccc}
\toprule
\multicolumn{1}{c}{{\textbf{Model}}} &
  {\textbf{\begin{tabular}[c]{@{}c@{}}Token\\ Count$^*$\end{tabular}}} &
  {\textbf{\begin{tabular}[c]{@{}c@{}}ARC\\ Challenge\end{tabular}}} &
  {\textbf{\begin{tabular}[c]{@{}c@{}}ARC\\ Easy\end{tabular}}} &
  {\textbf{LogiQA}} &
  {\textbf{\begin{tabular}[c]{@{}c@{}}Lambada \\ (OpenAI)\end{tabular}}} &
  {\textbf{PIQA}\textsuperscript{\#}} &
  {\textbf{SciQ}} &
  {\textbf{\begin{tabular}[c]{@{}c@{}}Wino\\ -Grande\textsuperscript{\#}\end{tabular}}} 
  % {\textbf{Winogrande\textsuperscript{\#}}}
  % {\textbf{Winogrande}}
  &
  {\textbf{WSC}} \\  
\midrule
\multicolumn{10}{c}{{\textbf{1B-Scale}}} \\ \midrule
{Hubble-1B} & {500B} & & & & & & & & \\
{-Standard} &
  &
  {0.40} &
  {0.72} &
  {0.25} &
  {7.43} &
  {0.76} &
  {0.95} &
  {0.63} &
  {0.41} \\
{-Perturbed} &
  &
  {0.40} &
  {0.72} &
  {0.25} &
  {7.23} &
  {0.76} &
  {0.94} &
  {0.63} &
  {0.45} \\
{Hubble-1B} & {100B} & & & & & & & & \\
{-Standard} &
  &
  {0.36} &
  {0.69} &
  {0.24} &
  {9.31} &
  {0.74} &
  {0.92} &
  {0.59} &
  {0.43} \\
{-Perturbed} &
  &
  {0.36} &
  {0.67} &
  {0.25} &
  {8.95} &
  {0.75} &
  {0.92} &
  {0.59} &
  {0.38} \\
{Pythia 1B} &
  {300B} &
  {0.28} &
  {0.57} &
  {0.25} &
  {10.86} &
  {0.70} &
  {0.92} &
  {0.53} &
  {0.43} \\
{Pythia 1.4B} &
  {300B} &
  {0.31} &
  {0.62} &
  {0.27} &
  {8.03} &
  {0.71} &
  {0.92} &
  {0.58} &
  {0.57} \\
{Bloom 1.1B} &
  {366B} &
  {0.28} &
  {0.53} &
  {0.25} &
  {24.84} &
  {0.68} &
  {0.90} &
  {0.53} &
  {0.37} \\
{Bloom 1.7B} &
  {366B} &
  {0.29} &
  {0.57} &
  {0.28} &
  {15.40} &
  {0.69} &
  {0.92} &
  {0.58} &
  {0.39} \\
{OPT 1.3B} &
  {180B} &
  {0.30} &
  {0.60} &
  {0.26} &
  {8.01} &
  {0.71} &
  {0.92} &
  {0.59} &
  {0.57} \\
{OLMo-2-1B} &
  {4T} &
  {0.46} &
  {0.76} &
  {0.27} &
  {6.26} &
  {0.77} &
  {0.96} &
  {0.66} &
  {0.45} \\
{Llama-3.2-1B} &
  {$\sim$9T} &
  {0.38} &
  {0.70} &
  {0.27} &
  {7.09} &
  {0.76} &
  {0.95} &
  {0.62} &
  {0.43} \\  

\midrule
\multicolumn{10}{c}{{\textbf{$\sim$ 8B-Scale}}} \\ \midrule
{Hubble-8B} & % -Standard
  {500B} &
  {0.58} &
  {0.84} &
  {0.32} &
  {3.71} &
  {0.82} &
  {0.98} &
  {0.77} &
  {0.56} \\
% {Hubble-8B-Perturbed} &
%   {500B} &
%   {0.57} &
%   {0.83} &
%   {0.30} &
%   {3.74} &
%   {0.82} &
%   {0.97} &
%   {0.79} &
%   {0.58} \\
{Hubble-8B} & % -Standard
  {100B} &
  {0.47} &
  {0.78} &
  {0.27} &
  {4.61} &
  {0.79} &
  {0.96} &
  {0.67} &
  {0.39} \\
% {Hubble-8B-Perturbed} &
%   {100B} &
%   {0.48} &
%   {0.78} &
%   {0.28} &
%   {4.66} &
%   {0.80} &
%   {0.96} &
%   {0.71} &
%   {0.47} \\
{Pythia 6.9B} &
  {300B} &
  {0.39} &
  {0.71} &
  {0.28} &
  {5.65} &
  {0.77} &
  {0.95} &
  {0.64} &
  {0.51} \\
{OPT 6.7B} &
  {180B} &
  {0.37} &
  {0.70} &
  {0.28} &
  {4.98} &
  {0.77} &
  {0.94} &
  {0.66} &
  {0.54} \\
{OLMo-2-7B} &
  {4T} &
  {0.63} &
  {0.85} &
  {0.34} &
  {3.90} &
  {0.81} &
  {0.97} &
  {0.77} &
  {0.78}  \\
{Llama-3.1-8B} &
  {15T+} &
  {0.58} &
  {0.85} &
  {0.33} &
  {3.93} &
  {0.82} &
  {0.98} &
  {0.77} &
  {0.63} \\
\bottomrule
\end{tabular}
\end{table}

\begin{table}[h]
\centering
\small
\caption{\textbf{Benchmark results using the DCLM v1 eval suite.} \textsc{DCLM-Baseline} and FineWeb edu results are copied from the official DCLM leaderboard. In general, Hubble models perform on par within their respective data and model scales. }
\label{tab:dclm_general_eval}
\begin{tabular}{lcccccc}
\toprule
Model & Params & Tokens & FLOPS & \textsc{Core} & MMLU & \textsc{Extended}\\ 
\midrule 
\multicolumn{7}{c}{\textbf{1B-Scale}} \\
\midrule
\textsc{DCLM-baseline} & 1.4B & 28.8B & 2.4e20 & 30.2 & 23.8 & 15.4 \\
FineWeb edu & 1.8B & 28B & 3.0e20 & 26.6 & 26.3 & 13.5 \\
\textsc{DCLM-baseline} & 1.4B & 144B & 1.2e21 & 36.1 & 26.4 & 18.6 \\
FineWeb edu & 1.8B & 140B & 1.5e21 & 33.8 & 25.5 & 17.6 \\
Pythia 1B & 1B & 300B & 1.8e21 & 24.8 & 25.1 & 13.5 \\
Pythia 1.4B & 1.4B & 300B & 2.5e21 & 27.8 & 25.4 & 14.2 \\
% Hubble 1B & 1.2B & 100B & 25.5 & 24.9 & 12.2 \\ V2
% Hubble 1B & 1.2B & 500B & 31.7 & 25.7 & 15.4 \\ V2
Hubble 1B & 1.2B & 100B & 7.2e20 & 27.8 & 24.9 & 14.5 \\ %v1
Hubble 1B & 1.2B & 500B & 3.6e21 & 34.2 & 25.7 & 17.7 \\ %v1

\midrule
\multicolumn{7}{c}{\textbf{$\sim$ 8B-Scale}} \\
\midrule 
\textsc{DCLM-baseline} & 6.9B & 138B & 5.7e21 & 44.8 & 42.2 & 28.8 \\
FineWeb edu & 7B & 138B & 5.8e21 & 38.7 & 26.3 & 22.1 \\
OPT 6.7B & 6.7B & 180B & 7.2e21 & 35.6 & 25.2 & 18.8 \\
\textsc{DCLM-baseline} & 6.9B & 276B & 1.1e22 & 48.9 & 50.8 & 31.8 \\
FineWeb edu & 7B & 276B & 1.2e22 & 41.9 & 37.4 & 24.5 \\
Pythia 6.9B & 6.9B & 300B & 1.2e22 & 35.7 & 25.4 & 19.6 \\
% Hubble 8B & 8.3B & 100B & 38.3 & 28.0 & 19.7 \\ V2
% Hubble 8B & 8.3B & 500B & 48.5 & 53.9 & 33.1 \\ V2
Hubble 8B & 8.3B & 100B & 5.0e21 & 40.8 & 28.0 & 22.0 \\ %V1
Hubble 8B & 8.3B & 500B & 2.5e22 & 50.0 & 53.9 & 34.6 \\ %V1

\bottomrule
\end{tabular}
\end{table}

% \subsection{Architecture design}
% % \todo{Aflah}{Describe base architecture, then describe the modified arch's}
% % what the architecture is (brief description of llama)
% % how does our architecture differ, and what is the justification.
% % training constraints (what we have right now)

% \todo{}{Explain Llama architecture more (rather than just stating ``Llama-3.2''). Include details on tied word embeddings that deviates from Llama. Better explain why we didn't copy tokenizers / used Olmo}

% We went through several phases of testing before converging on our final model architecture as described below - 

% \todo{}{- Add that 8B model has 36 layers to maximize the GPU usage with minor compute overhead - Reword the tokenizer choice details. It was made in parallel with the early testing and not in Phase 4}

% \begin{enumerate}
%     \item \textbf{Phase 1 – Choosing Model Sizes.} Our initial plan was to train a suite of models spanning sub-1B, 1B, ~3B, and 7B parameters. However, due to compute constraints, we narrowed this down to 1B and 7B models, aligning with the OLMo series \citep{groeneveld-etal-2024-olmo}.
%     \item \textbf{Phase 2 - Choosing the Model Architecture and Pretraining Framework.} The choice of model architecture was heavily dependent on the choice of pretraining framework as some frameworks made working with some architectures easier than others. E.g. GPT-NeoX did not support the non-parametric layernorm\footnote{It has since been added following discussion with the repository maintainers: \url{https://github.com/EleutherAI/gpt-neox/pull/1338}} used in OLMo nor did it support exporting the final model to the OLMo Model class in HuggingFace out of the box and we would have to bring the support for it. This wasn't the most challenging task but we decided to take the path of least resistance and since Llama is a fairly popular architecture we decided to go with it. In parallel we also tried using LitGPT with the OLMo architecture and integrated it in the framework\footnote{\url{https://github.com/Lightning-AI/litgpt/pull/1827}} however we ran into OOM errors which we couldn't fully debug\footnote{\url{https://github.com/Lightning-AI/litgpt/issues/1837}} and hence decided to stick with GPT-NeoX and the Llama architecture. Some other choices that we made are - 
%     \begin{enumerate}
%         \item \textbf{Parallel vs. Sequential Attention.} We opted for sequential architectures like Llama rather than parallel ones like Pythia, primarily because of their broader adoption.
%         \item \textbf{GQA vs. MQA vs. MHA.} We evaluated Grouped Query Attention (GQA), Multi-Head Attention (MHA), and Multi-Query Attention (MQA). GQA provided the best trade-off between speed and convergence, so we selected models that incorporated it.
%     \end{enumerate}
%     \item \textbf{Phase 3 – Finalizing the Model Design.} Within the chosen model family, we still faced the choice to whether build on Llama 2 or Llama 3. Our final choice rested on two considerations. First, we preferred the Llama 3 series since it incorporated more recent best practices and offered both 1B and 7B reference models. By contrast, the Llama 2 series only had a 1B-scale reference through TinyLlama\footnote{\url{https://huggingface.co/TinyLlama/TinyLlama-1.1B-Chat-v1.0}}, which was not trained by Meta. Second, we ran benchmarks at the 1B scale, comparing the Llama-3.2-1B model to TinyLlama, and observed a substantial training speed advantage ($\sim$1270 tokens/second/GPU) for Llama 3.2. This improvement stemmed from its wider but shallower architecture.
%     \item \textbf{Phase 4 – Tokenizer and Embedding Design.} This decision was made in parallel with the model architecture. We selected the OLMo tokenizer, which uses a much smaller vocabulary ($\sim$50K tokens) compared to the $\sim$128K vocabulary of the Llama 3 series. This choice substantially reduces the size of the embedding and output projection matrices. In addition, we opted not to tie input and output embeddings (unlike the Llama models but consistent with the Pythia suite) as this design facilitates interpretability research \citep{belrose2025elicitinglatentpredictionstransformers}.
% \end{enumerate}

% We base our 1B model on the Llama-3.2-1B architecture and our 8B model on Llama-3.1-8B, after comparing alternatives such as OLMo and Tiny-Llama. These choices offered the best balance of compatibility and performance: seamless integration with the GPT-NeoX framework, strong throughput during training, and broad ecosystem support across HuggingFace, inference libraries like vLLM and SGLang, and lightweight local runtimes such as llama.cpp. Notably, the Llama-3.2-1B design achieved a 1270+ tokens/sec/GPU speedup over Tiny-Llama, attributable to its wider and shallower structure.

% \subsection{Ablations}

% We conducted several ablations to finalize our architecture and training parameters:
% \begin{itemize}
%     \item \textbf{Batch size vs.\ activation checkpointing.} We explored multiple combinations of batch size and activation checkpointing \citep{chen2016trainingdeepnetssublinear}. Disabling activation checkpointing and using a smaller batch size with more gradient accumulation steps proved faster without hurting performance. Toggling individual checkpointing flags did not provide benefits, so we stuck to the default.
%     \item \textbf{Attention configurations.} We compared Grouped Query Attention (GQA), Multi-Head Attention (MHA), and Multi-Query Attention (MQA). GQA offered the best balance of speed and convergence, so we adopted it in our models.
%     \item \textbf{Distributed communication.} We tested different DeepSpeed bucket sizes for distributed communication but found negligible impact on training speed, so we retained the default.
%     \item \textbf{Kernel fusion.} We experimented with all relevant fusion kernels individually and in combination, observing consistent performance gains. We therefore enabled all available and relevant fusion kernels.

% \end{itemize}

% \todo{Aflah}{Add a table/plot with numbers across settings to compare and outline hardware used for the test experiments}


~\newpage
~\newpage
~\newpage
\section{Domain-specific results}
\label{app:domain_results}

\subsection{Copyright-specific Results}
\label{app:copyright_specific_results}

Additional evaluations for passages and paraphrases are shown in Figure \ref{fig:copyright-passages} and Figure \ref{fig:copyright-paraphrases}, respectively. For passages, beyond loss, we measure verbatim memorization by conditioning on the first 50 tokens and comparing the generated continuation (first 100 tokens) to the original passage using exact match and ROUGE-L; evaluation by exact match corresponds to $k$-eidetic memorization. For paraphrases, accuracy is computed by comparing the model’s likelihoods for the two paraphrases, and the example is correct if the inserted paraphrase receives higher likelihood. Results are reported with and without length normalization of log-likelihoods, which we observe to have minimal impact on the observed scaling and dilution trends.

% \textbf{The strength of memorization of passages is source dependent.} Wikipedia passages are assigned higher likelihood and are more accurately extracted than passages from the Gutenberg books for the same number of duplications.

\begin{figure}[h]
    \centering
    \includegraphics[width=\textwidth]{figures/domain-specific/copyright-passages.pdf}
    \caption{\textbf{Core results on Copyright Passages.} The first row evaluates memorization with the length-normalized log-likelihood of the models on the passages. The lower two rows measure the accuracy of verbatim generation, where the models are prompted to generate a 100-token continuation given a 50-token prefix.}
    \label{fig:copyright-passages}
\end{figure}

\begin{figure}[h]
    \centering
    \includegraphics[width=\textwidth]{figures/domain-specific/copyright-paraphrases.pdf}
    \caption{\textbf{Core results on Copyright Paraphrases.} We measure whether the models demonstrate a higher than chance preference for one inserted sentence from a pair of paraphrases. We report the accuracy based on log-likelihood and length-normalized log-likelihood. Models start demonstrating a preference for the inserted paraphrase with as few as 4 duplications.}
    \label{fig:copyright-paraphrases}
\end{figure}
\clearpage
% Not significant enough
% \paragraph{The strength of preference for the injected paraphrase is dataset-specific.} The standard models and the perturbed models on the set of unseen paraphrase pairs (0 duplicates) achieve a random-chance accuracy (50\%); this establishes that the models do not develop a preference for a particular paraphrase between two equivalent sentences. However, when one sentence from the pair is injected in the corpus, the models start assigning a higher likelihood to the seen paraphrase with a high accuracy. However, the accuracy on PAWS is higher than the accuracy on MRPC for the same number of duplicates.

\subsection{Privacy-specific Results}
\label{app:privacy_specific_results}

% In this section, we discuss additional evaluations on the \textbf{Biography} and \textbf{Chat} sub-domains.

\subsubsection{Direct PII Leakage} \label{app:direct_pii_leakage}

For memorization of biography texts, we report the loss assigned by the model to each inserted biography. Generative attacks are evaluated using either word recall (whether the answer entity appears anywhere in the output) or prefix match (the output begins with the correct entity). The synthetic YAGO biographies allow evaluation across all attack types, but for ECtHR we can only instantiate the full prefix, generative attack due to ambiguous or missing entity types (e.g., dates may refer to births or events). Figures \ref{fig:privacy-ecthr} and \ref{fig:privacy-yago} present attack success rates for ECtHR and YAGO, respectively. Figure \ref{fig:privacy-yago-meta} presents a breakdown of PII inference success by the type of private information.

\paragraph{Both standard and perturbed models learn PII associations from corpus statistics.} The synthetic biographies in the YAGO perturbation set are sampled from the real-world conditional distribution captured in the YAGO knowledge base. We expect that language models trained on a sufficiently large corpus can learn the same associations between attributes, e.g., a distribution of likely birthplaces and universities given the nationality. Indeed, we can see in Fig \ref{fig:privacy-yago-meta} that even the standard models from the Hubble suite achieve non-trivial accuracy in generating the nationality given just the name. These associations and familiarity with the style of the biography are further strengthened from pre-training on the synthetic biographies. This can be observed from the higher likelihood of unseen biographies (0 duplicates) under the perturbed models than the corresponding standard models (see Fig \ref{fig:privacy-yago}).

\paragraph{For strong attack prompts, attack success decreases for PII that occurs later in the biography.} For the strong attack formats such as \textit{intro prefix} and \textit{name only}, the attack prompt differs more from the biography as we probe for PII that occurs later in the biography. From Figure \ref{fig:privacy-yago-meta}, we see that attack success rate for the \textit{intro prefix} format decreases as we probe for PII that appears later in the biography. Two exceptions to this are UUID and email.  

\paragraph{Occupation, emails and UUIDs exhibit distinct memorization patterns.} There are three outliers from Figure \ref{fig:privacy-yago-meta}. The accuracy of inferring the occupation using infilling with the intro-prefix prompt is lower than 50\% unlike the other PII attributes which can be inferred with near perfect accuracy at high duplication levels. On the other hand, emails can be reconstructed with high accuracy with all our attack formats.
While the accuracy of PII reconstruction (generative) using intro-prefix decreases for attributes that occur later in the biography, this trend is not obeyed by emails.
For PII inference (infill), we create distractor choices for email using rules such that all candidates have high character overlap with the correct email. Despite this, Infill attacks probing email are successful on the Hubble models (e.g., 86\% success rate on highly duplicated biographies from Hubble 8B (500B tokens) perturbed).   
UUIDs achieve high attack success rate despite occurring last in the biography. Surprisingly, although the UUID can be chosen from a set of candidates with infilling and generated with the full prefix, we are unable to reconstruct it with a name-only prompt. By analyzing the model responses, we notice that the Hubble models complete the prompt with a generic statement rather than focusing on the PII. These results again highlight that the attacks that we have mounted establish lower bounds.

\begin{figure}[h]
    \centering
    \includegraphics[width=\textwidth]{figures/domain-specific/privacy-ecthr.pdf}
    \caption{\textbf{Attack success rates on ECtHR.} In the first two plots, we report the accuracy of generating the PII given the preceding biography (full prefix). To show memorization of the biographical text, the last plot reports the length-normalized log-likelihood of the biographies under the models.}
    \label{fig:privacy-ecthr}
\end{figure}
\clearpage

\begin{figure}[h]
    \centering
    \includegraphics[width=\textwidth]{figures/domain-specific/privacy-yago.pdf}
    \caption{\textbf{Attack success rates on YAGO}. Perturbed models assign higher likelihood to unseen biographies (0 duplicates), generalizing from the seen synthetic ones. Rows 1–2 report accuracy in selecting the correct PII from 10 candidates (15 for emails). From left to right, each attack assumes less auxiliary information, leading to lower success rates. Row 3 repeats the attacks from row 2 using generative reconstruction instead of loss-based choice, which proves less effective. Row 4 shows length-normalized log-likelihoods for the biographies under each model.}
    \label{fig:privacy-yago}
\end{figure}
\clearpage

\begin{figure}[h]
    \centering
    \includegraphics[width=0.97\textwidth]{figures/domain-specific/privacy-yago-meta.pdf}
    \caption{\textbf{Attack success rates on YAGO by PII type.} Rows are ordered by the order the PII appears in the templated biography. Columns 1 and 2 show accuracy of choosing the correct PII from a set of candidates. Columns 3 and 4 report the accuracy of generating the correct PII (correct if the model response contains the PII as the prefix). Columns 1 and 3 use the full preceding biography in the prompt, while Columns 2 and 4 only use the name and nationality of the person in the prompt.}
    \label{fig:privacy-yago-meta}
\end{figure}
\clearpage

% We instantiate loss-based choice attacks, in which the attacker narrows possible PII values and selects the most likely answer by comparing model likelihoods across options. When candidate options are unavailable, the attacker must use generative attacks, prompting the model to fill in the blank. Generative attacks are evaluated using either Word Recall (the answer entity appears anywhere in the output) or Prefix Match (the output begins with the correct entity). Table \ref{tab:pii_attack_definition} defines all instantiated attacks. The synthetic YAGO biographies allow evaluation across all attack types, while for ECtHR we can only instantiate the full prefix, generative attack due to ambiguous or missing entity types (e.g., dates may refer to births or events). Figures \ref{fig:privacy-ecthr} and \ref{fig:privacy-yago} present attack success rates for ECtHR and YAGO perturbation sets, respectively, and Figure \ref{fig:privacy-yago-meta} provides a detailed breakdown by PII type.

\subsubsection{Indirect PII Leakage} \label{app:indirect_pii_extraction}

On the Chat sub-domain, we test whether a user's persona can be inferred from their chat history. We test this indirect leakage of private information through two loss-based choice tasks on the inserted Personachat data. In the first task, \textit{Infill on Persona}, we test the models' accuracy on selecting the correct persona conditioned on the username from a set of 10 personas (distractors are drawn randomly from the other personas in the perturbation data). In the second task, \textit{Infill on Username}, we test whether the model can accurately select the correct username given the persona (distractor usernames are randomly drawn from the perturbation data). We illustrate the attacks in Table \ref{tab:chat_attack_definition}. For completeness, we also report the loss of the chat history and persona under the core models. We report findings in Figure \ref{fig:privacy-personachat}.

\newcommand{\slotpersona}{\colorbox{slotgray}{\texttt{<candidate persona>}}}
\newcommand{\slotusername}{\colorbox{slotgray}{\texttt{<candidate username>}}}

\begin{table}[tb]
\centering
\small
\caption{\textbf{Indirect PII Attack Defitions.} The instantiated indirect PII inference attacks are listed below. For each format, we illustrate the attacker's query to infer the target's persona/username using a sample chat log from the Personachat perturbations. Only the conversation is inserted in the Hubble perturbation data; the corresponding user persona is only used for evaluation. Candidates are drawn from other examples in the dataset.}
\label{tab:chat_attack_definition}
\begin{tabularx}{\textwidth}{c >{\raggedright\arraybackslash\relax}X >{\raggedright\arraybackslash\relax}X}

\toprule
\multicolumn{3}{>{\centering\hsize=\textwidth}X}{\textbf{Inserted Personachat conversation}}\\
\midrule
\multicolumn{3}{>{\centering\hsize=0.95\textwidth}X}{\texttt{chatbot: i like acting. i am in a telenovela now.
\attr{FloodBassoon371:} fun. dancing is my ticket to fame.
chatbot: what kind of dancing? were you in a show? i love musicals.
\attr{FloodBassoon371:} anything but dancing to country music, yuck, i hate it. chatbot: do you watch dancing with the stars?...}}\\
\midrule
\multicolumn{3}{>{\centering\hsize=\textwidth}X}{\textbf{Corresponding Personachat persona}}\\
\midrule
\multicolumn{3}{>{\centering\hsize=0.95\textwidth}X}{\texttt{i m an amazing dancer. i have blonde hair that reaches my knees. i volunteer at animal shelters. country music makes me cringe. i m a terrible speller.}}\\
\toprule
\textbf{Prompt Format} & \textbf{Example Query} & \textbf{Comments} \\
% \midrule
% Norm LL on Chat & \textit{chatbot: i like acting. i am in a telenovela now. FloodBassoon371: fun. dancing is my ticket to fame.
% chatbot: what kind of dancing? were you...} & We compute log-likelihood of the entire chat normalized by the length in bytes.\\
% \midrule
% Norm LL on Persona & chatbot: tell me a bit about yourself. InquiryTomb530:\textit{  i m an amazing dancer. i have blonde hair that reaches my knees. i volunteer at animal shelters...} & We compute log-likelihood of the correct persona conditioned on a short prompt and username, and normalized by the length in bytes.\\
\midrule
Infill on Persona & \texttt{FloodBassoon371: }\underline{\slotpersona} & We compare log-likelihood (with different normalizations) of the correct persona against 9 distractor personas conditioned on the username and report accuracy.\\
\midrule
(Prompted) Infill on Persona & \texttt{chatbot: tell me a bit about yourself. FloodBassoon371: }\underline{\slotpersona} & Same as Infill on Persona with an additional prompt.\\
\midrule
Infill on Username & \underline{\slotusername}\ul{\texttt{: i m an amazing dancer. i have blonde hair that reaches...}} & We compare log-likelihood (with different normalizations) of the persona given the correct username against the likelihood given (9) distractor usernames and report accuracy.\\
\midrule
(Prompted) Infill on Username & \texttt{chatbot: tell me a bit about yourself. }\underline{\slotusername}\ul{\texttt{: i m an amazing dancer. i have blonde hair that reaches...}} & Same as Infill on Username with an additional prompt.\\
\bottomrule
\end{tabularx}
\end{table}

\paragraph{Models assign lower likelihood to persona when memorizing chats.} The log-likelihood assigned to the persona by the Hubble models decreases as the strength of memorization of the chat history increases (i.e., with lower dilution). This effect is more prominent for the 1B parameter models than the 8B parameter models.

% Results also demonstrate dilution in the attack success rate as the corpus size increases. 

\begin{figure}[h]
    \centering
    \includegraphics[width=\textwidth]{figures/domain-specific/privacy-personachat.pdf}
    \caption{\textbf{Core results on Personachat.} Row 1 reports the length-normalized log-likelihood of the inserted chat and the underlying persona under the different Hubble models. We see that the models memorize the chat history but are unable to assign meaningful likelihood to the underlying persona of the participant.\\
    Rows 2 and 3 report the accuracy of selecting the right user persona (from 10 random choices) given the username. Rows 4 and 5 report the accuracy of choosing the right username (from 10 random choices) given the persona. Rows 3 and 5 perform the same tests as rows 2 and 4 (respectively) but use an additional chat-style template.
    }
    \label{fig:privacy-personachat}
\end{figure}
\clearpage

\subsection{Test set Contamination Results}
\label{app:testset_specific_results}

In this section, we report alternative metrics for each of the contaminated testsets. For \textbf{PopQA}, we report F1 score \citet{rajpurkar-etal-2018-know} in addition to the Exact Match (accuracy). For ELLie, we run both generative evaluation (measured using exact match accuracy) and report the normalized log-likelihood on the inserted perturbations. For all Infill-based tasks (WinoGrande-Infill, HellaSwag, PIQA, MUNCH), we report accuracy using alternative normalization schemes: \texttt{acc} directly compares the conditional log-likelihood of each choice, \texttt{acc\_norm} compares the conditional log-likelihood of each choice normalized by the byte-length of the choice, and \texttt{acc\_mutual\_info} compares the conditional log-likelihood of each choice after subtracting the unconditional log-likelihood of just the choice. For MCQ-style prompts, where the choices are part of the question and the expected answer is the label of the choice, we only report \texttt{acc} since the option lengths are all the same. We report the performance on PopQA, HellaSwag, MMLU, and PIQA in Figure \ref{fig:testset-set1}. We report the performance on different WinoGrande formats in Figure \ref{fig:testset-wg}. Finally, we report performance on the new test sets, MUNCH and ELLie, in Figure \ref{fig:testset-set2}.

% mentioned in main text
% \paragraph{Contamination can boost accuracy with very low duplication.} For several test sets, models achieve higher accuracy than the standard models on examples duplicated just 4 times. 

% \paragraph{Standard models demonstrate performance scaling based on model and corpus size.} Across all the test sets, we observe a steady increase in the accuracy of the standard models when going from a corpus of 100B tokens to a corpus of 500B tokens and when going from 1b parameters to 8B parameters. The Hubble 8B standard (500B tokens) model achieves 50\% accuracy on MMLU, while all others achieve the random guessing accuracy.


\textbf{Models do not generalize from contaminated examples to the corresponding minimal pairs.} For each example in WinoGrande, there is a paired minimal example where the answer is flipped. When inserting examples, we make sure to only use one example from each pair as a part of the perturbation data. This allows us to evaluate whether the perturbed models can generalize to the minimal pair from training on the inserted example. Our results on WinoGrande show that the models.

% i feel like this conflicts with the point that: models don't seem to reliably generalize from contamination.
% \paragraph{Contamination can improve, hurt, or leave unchanged within-task generalization.} On PopQA, we see that the accuracy of the perturbed models in higher than the standard models even on unseen examples (0 duplicates). On MMLU, we see that the performance on unseen examples is unchanged. However, on Winogrande, HellaSwag, and PIQA, we see that the accuracy on unseen examples is worse than the accuracy of the standard model. The lack of generalization is also demonstrated with the paraphrase experiments in Appendix \ref{app:paraphrase_results}, where we find that a perturbed model trained on paraphrased MMLU problems is unable to answer the original questions.

\begin{figure}[h]
    \centering
    \includegraphics[width=\textwidth]{figures/domain-specific/testset-set2.pdf}
    \caption{\textbf{Core results on ELLie and MUNCH.} }
    \label{fig:testset-set2}
\end{figure}

\paragraph{MUNCH is solved by standard models.} From Figure \ref{fig:testset-set2}, we see that both standard and perturbed models achieve very high accuracy on MUNCH. Each MUNCH example consists of two sentences, one of which is the original, valid sentence, and the other is modified by swapping one word from the original sentence for an inappropriate synonym. The task is to identify which sentence is meaningful and valid. Our core models are all competent at language modeling and thus can solve the task with high accuracy ($>$ 96\%). Even so, we see increased accuracy with perturbed models on the examples that are duplicated more than 16 times.

\paragraph{ELLie examples are minimal pairs making it isolate to disentangle the effect of duplication.} ELLie is a task that tests whether language models can understand sentences with ellipsis. From Figure \ref{fig:testset-set2}, we see that the standard model achieve near 0 accuracy on the task. On the other hand, perturbed models achieve accuracy greater than 50\% even on examples that were never duplicated. On further analysis, we realized that the examples in ELLie are minimal pairs.\footnote{Many examples in ELLie contain the same first sentence but different query sentences (the second sentence). Thus, they passed our deduplication check.} When we insert the examples in our corpus, examples with the same first sentence were put in different duplication bins, e.g., of all the examples with the same core sentence, some examples were sometimes duplicated 0 times and other examples were duplicated 16 times. Thus, we see that models achieve high accuracy on examples duplicated 0 times. This invalidates the use of ELLie for studying dilution.

\begin{figure}[h]
    \centering
    \includegraphics[width=\textwidth]{figures/domain-specific/testset-set1.pdf}
    \caption{\textbf{Core results on Test Sets (Part 1).} Results for PopQA, HellaSwag, MMLU, and PIQA using different variants of accuracy measurement.}
    \label{fig:testset-set1}
\end{figure}

\begin{figure}[h]
    \centering
    \includegraphics[width=\textwidth]{figures/domain-specific/testset-wg.pdf}
    \caption{\textbf{Core results and variants on WinoGrande.} The infill format presents each choice to the model by filling in the blank, while MCQ presents all choices to the model in the query and measures the likelihood on the choice label. Rows 1 and 2 evaluate accuracy on duplications \textit{inserted} with the Infill format. Rows 3 and 4 evaluate accuracy on duplications \textit{inserted} with the MCQ format. Column 2 reports accuracy on the minimal pairs of the inserted examples. Rows 1 and 4 use the Infill format for evaluation while rows 2 and 3 use the MCQ format for evaluation.}
    \label{fig:testset-wg}
\end{figure}

~\newpage
~\newpage
\section{Additional Results}
\label{app:additional_results}

\subsection{Timing Runs}
\label{app:timing_results}

To study how memorization evolves over training, we evaluate memorization on intermediate checkpoints every 2,000 steps up to 48,000. We also include Timing runs to analyze forgetting. Figure \ref{fig:checkpoint-forgetting-curves} reports normalized log-likelihood on Wikipedia passages and accuracy on MRPC paraphrases, each inserted 256 times. Across all four Timing runs, both metrics rise as duplicated data are encountered, peak once all perturbations have been seen, and then decay.

% \subsection{Temporal Effects}

\begin{figure}[h]
    \centering
    \includegraphics[width=\textwidth]{figures/injectrange/checkpoint-forgetting-curves.pdf}
    \caption{\textbf{Forgetting curves for the intermediate checkpoints of Timing runs.} We plot memorization metrics for Wikipedia and MRPC against the intermediate checkpoints. We report results on the subset of examples duplicated 256 times. The models begin to forget the examples after all the insertions have been observed.}
    \label{fig:checkpoint-forgetting-curves}
\end{figure}

\begin{figure}[h]
    \centering
    \includegraphics[width=\textwidth]{figures/injectrange/injectrange-full.pdf}
    \caption{\textbf{Evaluation on the InsertRange models.} Models that were trained on perturbations only in the early stages of training have lower performance on the memorization tasks than models trained on perturbations in the late stages of training. \texttt{InsertRange(x,y)} denotes a model trained on a corpus with perturbations inserted in batches between x\% and y\% of training.}
    \label{fig:injectrange-full}
\end{figure}
\clearpage

\subsection{Paraphrased Runs}
\label{app:paraphrase_results}

\begin{figure}[h]
    \centering
    \includegraphics[width=\textwidth]{figures/paraphrase/paraphrase-full.pdf}
    \caption{\textbf{Performance of Hubble perturbed models trained on paraphased insertions.} The models do not generalize from paraphrased examples seen in training to the original examples. However, PII can be reconstructed from models trained on paraphrased biographies, even with stronger attacks.}
    \label{fig:paraphrase-full}
\end{figure}

\textbf{Data Preparation for Paraphrase Runs.} We construct paraphrased variants of the YAGO biographies and MMLU test set with \texttt{gpt-4.1-mini}. Unless otherwise noted, generation uses \texttt{temperature=1} and \texttt{top\_p=1}.
For each original perturbation example to be inserted, we obtain as many paraphrases as its required duplication count.

\begin{itemize}
    \item \textbf{MMLU paraphrases.} We follow the paraphrasing instruction of \citet{mmlu_paraphrase}. When a paraphrase query is declined by \texttt{gpt-4.1-mini} API's safety filter, we use \texttt{gemini-2.5-flash-lite} with the same parameters. 
    
    \item \textbf{YAGO paraphrases.} We adopt the diverse-style watermarking generation instructions from \citet{yago_paraphrase}. Each paraphrase is checked with a string-matching validator to ensure all biographical attributes are preserved. A paraphrase is accepted only if every attribute appears. We follow the procedure until we obtain the required number of valid paraphrases.
\end{itemize}

We train two perturbed models (1B and 8B parameters) on 100B tokens with the same perturbation data as the core perturbed model but with two data sets paraphrased: MMLU and YAGO Biographies. We evaluate the behavior of the `paraphrase' models on MMLU and YAGO evaluations in Figure \ref{fig:paraphrase-full}.

\textbf{PII can be leaked from paraphrased biographies with loss-based choice and generative evaluations.} The weakest attacks, which assume that the attacker has access to all PII about a person except one fact, are successful on models trained with paraphrased biographies. However, they have lower effectiveness than extracting the facts from the model that was trained on the original biographies. PII can be extracted with 100\% accuracy from the core 8B perturbed model using the full prefix and full suffix MCQ format. This accuracy drops to 89\% when extracting PII from the paraphrase model. Surprisingly, when using stronger attacks (attacker has access to only the persons name), PII is more accurately extractable from the 8B model trained on paraphrased biographies compared to the core models. However, this finding depends on the format of the attack and scale; generative evaluations cannot extract PII from the 1B paraphrased model. 

\textbf{Models cannot generalize from paraphrased MMLU to the original examples.} We find that both models (1B and 8B parameters) obtain random accuracy on the MMLU MCQ evaluations when trained on paraphrased versions of the examples. 

\newpage
\subsection{Architecture Runs}
\label{app:architecture_results}

\begin{figure}[h]
    \centering
    \includegraphics[width=\textwidth]{figures/architecture/architecture-full.pdf}
    \caption{Deeper models memorize slightly more than shallower models. We train three 1B-parameter models with 8, 16, and 32 layers, adjusting width to keep total parameters constant ($\approx1.2$B). All models are pre-trained on 100B tokens. As shown in Figure \ref{fig:architecture-full}, the deeper (narrower) model memorizes slightly more than the base 16-layer model, while the shallower (wider) model memorizes less.}
    \label{fig:architecture-full}
\end{figure}




\newpage
\section{Additional \textsc{HubbleMIA} Results}
\label{app:mia_results}

We instantiate 6 variants of MIA benchmarks using the Hubble suite, using 4 models and 3 perturbation datasets (passages from Gutenberg Unpopular, biographies from YAGO, and contaminated examples from MMLU). As discussed in \S~\ref{sec:hubble_as_mia}, the standard models use entirely unseen data for both the seen and unseen sets, serving only as a reference point i.e. no method should achieve better-than-random accuracy in this setting.

%\begin{itemize}
    %\item Table \ref{tab:hubble_mia_8b_perturbed} reports MIA performance on the Hubble 8B Perturbed model.
    %\item Table \ref{tab:hubble_mia_8b_standard} reports MIA performance on the Hubble 8B Standard model.
    %\item Table \ref{tab:hubble_mia_1b_perturbed} reports MIA performance on the Hubble 1B Perturbed model.
    %\item Table \ref{tab:hubble_mia_1b_standard} reports MIA performance on the Hubble 1B Standard model.
%\end{itemize}

% \input{tables/mia_table_8b_perturbed_cont}

\begin{table}[h]
\centering
\small
\caption{\textbf{Membership inference performance on various benchmarks with Hubble 1B Perturbed.} The Dup values indicate the composition of the seen set: for example, \textit{Dup $\neq$ 0} means the attack compares all seen data against unseen data, whereas \textit{Dup = K} means the attack compares unseen data against data that was included exactly $K$ times in the seen set.}
\label{tab:hubble_mia_1b_perturbed}
\begin{tabular}{clcccccc}
\toprule
\multirow{2}{*}{\textbf{Evaluation}} & \multicolumn{1}{c}{\multirow{2}{*}{\textbf{MIA}}} & \multicolumn{6}{c}{\textbf{Hubble 1B Perturbed (500B tokens)}}                         \\ \cmidrule(lr){3-8} 
                                     & \multicolumn{1}{c}{}                                        & Dup $\neq$ 0 & Dup = 1 & Dup = 4 & Dup = 16 & Dup = 64 & Dup = 256 \\ \midrule
\multirow{4}{*}{\shortstack{Gutenberg\\Unpopular}} & Loss                                                        &   0.552       &   0.52      &  0.504       &   0.552       &   0.73       &  0.999           \\  
                                     & MinK\%                                                      &   0.552       &   0.52      &  0.504       &   0.552       &   0.729       &  0.999           \\  
                                     & MinK\%++                                                    &   0.575       &   0.513      &  0.53       &   0.605       &   0.825       &  1.0           \\  
                                     & ZLib                                                        &   0.543       &   0.511      &  0.497       &   0.533       &   0.729       &  1.0           \\ \midrule
\multirow{4}{*}{\shortstack{Yago\\Biographies}}    & Loss                                                        &   0.606       &   0.506      &  0.557       &   0.696       &   0.928       &  1.0           \\  
                                     & MinK\%                                                      &   0.606       &   0.506      &  0.556       &   0.695       &   0.927       &  1.0           \\  
                                     & MinK\%++                                                    &   0.615       &   0.509      &  0.565       &   0.715       &   0.947       &  1.0           \\  
                                     & ZLib                                                        &   0.596       &   0.499      &  0.551       &   0.679       &   0.899       &  1.0           \\ \midrule
\multirow{4}{*}{MMLU}                & Loss                                                        &   0.557       &   0.499      &  0.524       &   0.575       &   0.748       &  1.0           \\  
                                     & MinK\%                                                      &   0.557       &   0.5      &  0.524       &   0.575       &   0.747       &  1.0           \\  
                                     & MinK\%++                                                    &   0.605       &   0.522      &  0.556       &   0.681       &   0.887       &  0.996           \\  
                                     & ZLib                                                        &   0.548       &   0.502      &  0.521       &   0.556       &   0.67       &  0.998           \\ \bottomrule
\end{tabular}
\end{table}

\begin{table}[h]
\centering
\small
\caption{\textbf{Membership inference performance on various benchmarks with Hubble 8B Standard}. The Dup values indicate the composition of the seen set: for example, \textit{Dup $\neq$ 0} means the attack compares all seen data against unseen data, whereas \textit{Dup = K} means the attack compares unseen data against data that was included exactly $K$ times in the seen set.}
\label{tab:hubble_mia_8b_standard}
\begin{tabular}{clcccccc}
\toprule
\multirow{2}{*}{\textbf{Evaluation}} & \multicolumn{1}{c}{\multirow{2}{*}{\textbf{MIA}}} & \multicolumn{6}{c}{\textbf{Hubble 8B Standard (500B tokens)}}                         \\ \cmidrule(lr){3-8} 
                                     & \multicolumn{1}{c}{}                                        & Dup $\neq$ 0 & Dup = 1 & Dup = 4 & Dup = 16 & Dup = 64 & Dup = 256 \\ \midrule
\multirow{4}{*}{\shortstack{Gutenberg\\Unpopular}} & Loss                                                        &   0.507       &   0.522      &  0.486       &   0.495       &   0.54       &  0.545           \\  
                                     & MinK\%                                                      &   0.507       &   0.522      &  0.486       &   0.495       &   0.54       &  0.545           \\  
                                     & MinK\%++                                                    &   0.504       &   0.517      &  0.493       &   0.499       &   0.484       &  0.543           \\  
                                     & ZLib                                                        &   0.497       &   0.514      &  0.48       &   0.474       &   0.535       &  0.544           \\ \midrule
\multirow{4}{*}{\shortstack{Yago\\Biographies}}    & Loss                                                        &   0.499       &   0.489      &  0.499       &   0.519       &   0.486       &  0.516           \\  
                                     & MinK\%                                                      &   0.499       &   0.489      &  0.499       &   0.519       &   0.487       &  0.516           \\  
                                     & MinK\%++                                                    &   0.503       &   0.5      &  0.503       &   0.507       &   0.505       &  0.505           \\  
                                     & ZLib                                                        &   0.495       &   0.479      &  0.5       &   0.523       &   0.481       &  0.495           \\ \midrule
\multirow{4}{*}{MMLU}                & Loss                                                        &   0.502       &   0.506      &  0.503       &   0.512       &   0.459       &  0.476           \\  
                                     & MinK\%                                                      &   0.502       &   0.506      &  0.503       &   0.512       &   0.458       &  0.476           \\  
                                     & MinK\%++                                                    &   0.506       &   0.51      &  0.505       &   0.514       &   0.497       &  0.45           \\  
                                     & ZLib                                                        &   0.501       &   0.505      &  0.504       &   0.506       &   0.463       &  0.495           \\ \bottomrule
\end{tabular}
\end{table}

% \begin{table}[tb]
\centering
\small
\caption{\textbf{Membership inference performance on various benchmarks with Hubble 1B Standard.} The Dup values indicate the composition of the seen set: for example, \textit{Dup $\neq$ 0} means the attack compares all seen data against unseen data, whereas \textit{Dup = K} means the attack compares unseen data against data that was included exactly $K$ times in the seen set.}
\label{tab:hubble_mia_1b_standard}
\begin{tabular}{clcccccc}
\toprule
\multirow{2}{*}{\textbf{Evaluation}} & \multicolumn{1}{c}{\multirow{2}{*}{\textbf{MIA}}} & \multicolumn{6}{c}{\textbf{Hubble 1B Standard (500B tokens)}}                         \\ \cmidrule(lr){3-8} 
                                     & \multicolumn{1}{c}{}                                        & Dup $\neq$ 0 & Dup = 1 & Dup = 4 & Dup = 16 & Dup = 64 & Dup = 256 \\ \midrule
\multirow{4}{*}{\shortstack{Gutenberg\\Unpopular}} & Loss                                                        &   0.503       &   0.517      &  0.484       &   0.494       &   0.534       &  0.531           \\  
                                     & MinK\%                                                      &   0.502       &   0.517      &  0.483       &   0.494       &   0.534       &  0.531           \\  
                                     & MinK\%++                                                    &   0.5       &   0.509      &  0.493       &   0.497       &   0.481       &  0.529           \\  
                                     & ZLib                                                        &   0.493       &   0.509      &  0.477       &   0.471       &   0.529       &  0.533           \\ \midrule
\multirow{4}{*}{\shortstack{Yago\\Biographies}}    & Loss                                                        &   0.495       &   0.488      &  0.494       &   0.51       &   0.494       &  0.509           \\  
                                     & MinK\%                                                      &   0.495       &   0.487      &  0.494       &   0.51       &   0.494       &  0.508           \\  
                                     & MinK\%++                                                    &   0.5       &   0.499      &  0.501       &   0.494       &   0.518       &  0.497           \\  
                                     & ZLib                                                        &   0.494       &   0.481      &  0.498       &   0.516       &   0.489       &  0.49           \\ \midrule
\multirow{4}{*}{MMLU}                & Loss                                                        &   0.502       &   0.506      &  0.502       &   0.519       &   0.459       &  0.48           \\  
                                     & MinK\%                                                      &   0.503       &   0.506      &  0.502       &   0.519       &   0.459       &  0.481           \\  
                                     & MinK\%++                                                    &   0.509       &   0.512      &  0.509       &   0.53       &   0.475       &  0.448           \\  
                                     & ZLib                                                        &   0.501       &   0.504      &  0.503       &   0.508       &   0.465       &  0.494           \\ \bottomrule
\end{tabular}
\end{table}

\newpage
\section{Additional \textsc{HubbleUnlearning} results}
\label{sec:unlearn_grid_search} \label{sec:unlearn_full_results}
% \subsection{Grid Search Configurations}

Below are the detailed hyperparameters for each method: 

\begin{table}[h]
\centering
\small
\setlength{\tabcolsep}{6pt}
\renewcommand{\arraystretch}{1.2}
\begin{tabular}{lccc}
\toprule
\textbf{Hyperparameter} & \textbf{RMU} & \textbf{RR} & \textbf{SatImp} \\
\midrule
\textbf{Training type} & Layer FT & LoRA FT & Full FT \\
\textbf{Layers / Targets} & 5, 6, 7 & 10, 20 (transform all) & --- \\
\textbf{LoRA Rank / $\alpha$ / Dropout} & --- & 16 / 16 / 0.05 & --- \\
\textbf{LoRRA $\alpha$} & --- & 10 & --- \\
\textbf{Alpha ($\alpha$)} & 100, 1000, 10000 & --- & 0.01, 0.1, 1 \\
\textbf{Steering coefficient} & 5, 50, 500 & --- & --- \\
\textbf{$\beta_1$, $\beta_2$} & --- & --- & (5, 6), 1 \\
\textbf{Learning rate} & 5e-5, 1e-5, 5e-4 & 5e-5, 1e-4, 5e-4, 1e-3 & 1e-5, 5e-5, 1e-4 \\
\textbf{Effective batch size} & 4 & 8 & 16 \\
\textbf{Epochs} & 4, 8 & 4, 8 & --- \\
\textbf{Sample max length} & 512 & 256 & 256 \\
\bottomrule
\end{tabular}
\caption{\textbf{Grid search configurations for unlearning methods.} Each method is tuned over the listed hyperparameters. RMU and RR involve partial fine-tuning, while SatImp uses full fine-tuning.}
\label{tab:unlearn_grid_search}
\end{table}
\clearpage

% \subsection{Full Unlearning Results}

We provide the full scale unlearning results for Gutenberg in Figure \ref{fig:appendix_unlearn_gutenberg} and YAGO in Figure \ref{fig:appendix_unlearn_yago}.


\begin{figure}[h]
  \centering
  \begin{subfigure}{.48\textwidth}
    \centering
    \includegraphics[width=\linewidth]{figures/Hubble_unlearn/appendix/gutenberg_keep_keep.pdf}\\[0.5em]
    \includegraphics[width=\linewidth]{figures/Hubble_unlearn/appendix/gutenberg_wikitext_keep.pdf}
    % \caption{Gutenberg with Keep set as retain set}
  \end{subfigure}\hfill
  \begin{subfigure}{.48\textwidth}
    \centering
    \includegraphics[width=\linewidth]{figures/Hubble_unlearn/appendix/gutenberg_keep_test.pdf}\\[0.5em]
    \includegraphics[width=\linewidth]{figures/Hubble_unlearn/appendix/gutenberg_wikitext_test.pdf}
    % \caption{Gutenberg with Wikitext as retain set}
  \end{subfigure}
  \caption{\textbf{Unlearning results on Gutenberg Unpopular.} Unlearning results using (out-of-domain, unseen) Wikitext (lower row) and (in-domain, seen) Keep set (upper row) as the retain sets. None of the unlearning methods simultaneously achieve the target behavior on both the seen Keep set (left column) and the unseen Test set (right column).}
  \label{fig:appendix_unlearn_gutenberg}
\end{figure}

\begin{figure}[h]
  \centering
  \begin{subfigure}{.48\textwidth}
    \centering
    \includegraphics[width=\linewidth]{figures/Hubble_unlearn/appendix/yago_keep_keep.pdf}\\[0.5em]
    \includegraphics[width=\linewidth]{figures/Hubble_unlearn/appendix/yago_wikitext_keep.pdf}
    % \caption{YAGO with Keep set as retain set}
  \end{subfigure}\hfill
  \begin{subfigure}{.48\textwidth}
    \centering
    \includegraphics[width=\linewidth]{figures/Hubble_unlearn/appendix/yago_keep_test.pdf}\\[0.5em]
    \includegraphics[width=\linewidth]{figures/Hubble_unlearn/appendix/yago_wikitext_test.pdf}
    % \caption{YAGO with Wikitext as retain set}
  \end{subfigure}
  \caption{\textbf{Unlearning results on YAGO biographies.} Unlearning results using (out-of-domain, unseen) Wikitext (lower row) and (in-domain, seen) Keep set (upper row) as the retain sets. None of the unlearning methods simultaneously achieve the target behavior on both the seen Keep set (left column) and the unseen Test set (right column).}
  \label{fig:appendix_unlearn_yago}
\end{figure}
\clearpage

% \begin{figure*}[t]
%   \centering
%   \begin{subfigure}{.48\textwidth}
%     \centering
%     \includegraphics[width=\linewidth]{figures/Hubble_unlearn/appendix/mmlu_keep_keep.pdf}\\[0.5em]
%     \includegraphics[width=\linewidth]{figures/Hubble_unlearn/appendix/mmlu_wikitext_keep.pdf}
%     \caption{MMLU with Keep set as retain set}
%   \end{subfigure}\hfill
%   \begin{subfigure}{.48\textwidth}
%     \centering
%     \includegraphics[width=\linewidth]{figures/Hubble_unlearn/appendix/mmlu_keep_test.pdf}\\[0.5em]
%     \includegraphics[width=\linewidth]{figures/Hubble_unlearn/appendix/mmlu_wikitext_test.pdf}
%     \caption{MMLU with Wikitext as retain set}
%   \end{subfigure}
%   \caption{\textbf{MMLU results.} Full unlearning results across two retain sets.}
%   \label{fig:appendix_unlearn_mmlu}
% \end{figure*}

% \paragraph{RMU \citep{wmdp}:}
% \begin{itemize}
%   \item Layer Fine-tuning:
%   \begin{itemize}
%     \item Layers: 5, 6, 7
%   \end{itemize}
%   \item Alpha: 100, 1000, 10000
%   \item Steering coefficient: 5, 50, 500
%   \item Learning rate: \texttt{5e-5}, \texttt{1e-5}, \texttt{5e-4}
%   \item Effective batch size: 4
%   \item Epochs: 4, 8
%   \item Sample max length: 512
% \end{itemize}

% \paragraph{RR \citep{rr}:}
% \begin{itemize}
%   \item LoRA Fine-tuning:
%   \begin{itemize}
%     \item LoRA Rank: 16
%     \item LoRA $\alpha$: 16
%     \item LoRA dropout: 0.05
%   \end{itemize}
%   \item LoRRA Alpha: 10
%   \item Target layers: 10, 20
%   \item Transform layers: all
%   \item Learning rate: \texttt{5e-5}, \texttt{1e-4}, \texttt{5e-4}, \texttt{1e-3}
%   \item Effective batch size: 8
%   \item Epochs: 4, 8
%   \item Sample max length: 256
% \end{itemize}

% \paragraph{SatImp \citep{satimp}:}
% \begin{itemize}
%   \item $\alpha$: 0.01, 0.1, 1
%   \item $\beta_1$: 5, 6
%   \item $\beta_2$: 1
%   \item Learning rate: \texttt{1e-5}, \texttt{5e-5}, \texttt{1e-4}
%   \item Effective batch size: 16
%   \item Sample max length: 256
% \end{itemize}

% After grid search, we evaluate the unlearned checkpoints on tinyMMLU, tinyWinogrande, and tinyHellaswag from TinyBenchmarks \citep{tinybenchmarks} for general capabilities preservation, and discard checkpoints with average performance degradation exceeding 10\%.


%% Figures
% we will need to reorganize these
\newpage
\section{Additional Plots}

\begin{figure}[h]
    \centering
    \includegraphics[width=\textwidth]{figures/model_scale/model_scale-dclm_500B.pdf}
    \caption{\textbf{Larger models memorize at lower duplicates.} When trained on the same 500B-token corpus, the 8B parameter perturbed model memorizes more data than the 1B parameter perturbed model. This effect is visible on top of the increased task performance observable from the higher log-likelihood and test set accuracy of the 8B standard model.}
    \label{fig:model_scale-dclm_500b}
\end{figure}

% \subsection{Interference Results}
\begin{figure}[h]
    \centering
    \includegraphics[width=\textwidth]{figures/interference/interference-full.pdf}
    \caption{\textbf{The perturbed model matches the behavior of domain-specific models on the respective set of evaluations.} The perturbed model matches the \texttt{copyright\_only} model in memorizing the copyright passages and paraphrases, \texttt{privacy\_only} model in generating memorized PII from biographies and chat, and \texttt{testset\_only} model in memorizing the testsets. Thus, the perturbed model can be used to study individual domains despite being jointly trained on all three domains.}
    \label{fig:interference-full}
\end{figure}

% \subsection{Paraphrased Perturbations}
% what's this figure for?
% \begin{figure}[h]
%    \centering
%    \includegraphics[width=\textwidth]{figures/paraphrase/paraphrase-all.pdf}
%    \caption{\textbf{Sanity check for Paraphrase runs.} Paraphrasing only affects the changed perturbations. Other evaluations are unaffected.}
%    \label{fig:paraphrase-all}
%\end{figure}


\end{document}